\chapter{冒险、异常与性能}

\begin{keyidea}
并行执行只有在依赖、异常和提交边界都受控时才保持正确。本章把异常控制流与 CPI、吞吐、带宽和瓶颈分析放进同一套定量框架。
\end{keyidea}

\section{一、异常、中断与精确异常}

\topic{异常和中断是受控的非顺序 PC 更新}
异常、中断和系统调用都可以视为特殊控制流。普通分支由指令显式请求；异常由当前指令执行过程中发现的条件触发，例如非法 opcode、除零、未对齐访问、页错误或总线错误；中断来自外部设备或定时器。三者最终都要让 CPU 保存足够的旧状态，并把 PC 改到一个规定入口。

\begin{longtable}{L{0.18\linewidth}L{0.38\linewidth}L{0.34\linewidth}}
\toprule
事件 & 触发来源 & 入口动作 \\
\midrule
同步异常 & 当前指令导致，如非法指令、页错误、除零 & 保存故障 PC 和原因，跳异常向量 \\\tablerowrule
外部中断 & 设备、定时器或中断控制器请求 & 在可接收边界保存 PC/Flags，跳中断入口 \\\tablerowrule
系统调用 & 指令显式请求进入特权服务 & 切换权限和入口，保存返回位置 \\\tablerowrule
复位 & 电源或复位逻辑强制初始化 & PC 装入复位向量，控制状态清零或设默认值 \\
\bottomrule
\end{longtable}

异常路径最重要的是提交边界。若一条加载指令因页错误失败，目标寄存器不应保存随机数据；若一条存储指令访问设备前发生权限错误，设备不应看到写命令；若一条除法发生除零异常，后续指令不应像已经成功执行那样提交。简单 CPU 可以在每条指令结束时才接收中断；复杂流水线 CPU 则需要记录哪些指令已经可以提交，哪些仍可能被冲刷。

\begin{lstlisting}
// 注释（推进）：按行追踪数据来源、经过的部件和最终写入目标。
// 注释（边界）：只有标出的时钟沿、阶段或异常入口会改变体系结构可见状态。
faulting load:
  IF:  fetch instruction
  ID:  decode MOV
  EX:  form address
  MEM: address translation fails
  WB:  suppressed
  PC:  exception_vector
\end{lstlisting}

这个例子说明，异常不是“多跳一次分支”这么简单。它还要阻止错误写回，把失败原因带给入口，并保证重试或终止时系统能判断发生了什么。异常是硬件检测到无法继续普通提交后，选择另一条受控 PC 更新路径。

\begin{pdfdiagram}
  \node[hw compute, minimum width=3.4cm] (run) at (0,5.0) {指令正常执行\\IF→ID→EX→MEM→WB};
  \node[hw control, minimum width=4.0cm] (fault) at (0,3.5) {执行中检测到异常条件\\非法 opcode / 页错误 / 除零};
  \node[hw control, minimum width=3.4cm] (irq) at (5.6,3.5) {外部中断请求\\在指令边界接收};
  \node[hw state, minimum width=4.0cm] (save) at (0,1.9) {保存故障 PC、原因码\\与程序员可见现场};
  \node[hw compute, minimum width=3.6cm] (vec) at (0,0.4) {按事件类型查向量表\\PC ← 异常入口};
  \node[hw memory, minimum width=3.4cm] (handler) at (0,-1.1) {异常/中断处理程序};
  \node[hw compute, minimum width=3.0cm] (ret) at (-2.7,-2.7) {可恢复：恢复现场\\RET 回故障点};
  \node[hw control, minimum width=3.0cm] (halt) at (2.7,-2.7) {不可恢复：\\停机 / 错误 LED};
  \draw[hw bus] (run) -- (fault);
  \draw[hw bus] (fault) -- node[hw label, right]{同步异常} (save);
  \draw[hw bus] (irq.south) -- (5.6,1.9) -- node[hw label, above, pos=0.35]{异步} (save.east);
  \draw[hw bus] (save) -- (vec);
  \draw[hw bus] (vec) -- (handler);
  \draw[hw bus] (handler.south) -- ++(0,-0.4) -| (ret.north);
  \draw[hw bus] (handler.south) -- ++(0,-0.4) -| (halt.north);
  \draw[hw return] (ret.west) -- ++(-0.7,0) -- ++(0,7.35) node[hw label, rotate=90, pos=0.55, below]{处理完重新进入取指循环} -- (run.west);
\end{pdfdiagram}
\diagramnote{异常与中断共用同一条入口骨架：保存现场、查向量表、跳转处理程序。区别只在触发来源和接收时机——同步异常必须禁止当前指令的错误写回，异步中断只在指令边界接收。处理完后沿虚线回到取指循环；裸机程序若无法恢复，至少要把系统停在一个可诊断的状态，而不是带着错误现场继续跑。}

\topic{中断向量把异常号变成入口地址}
异常或中断被接收之后，硬件面对的问题很具体：PC 应该改成多少。普遍做法是把所有入口预先存成一张向量表——本质是一个“异常号映射入口地址”的数组。事件类型被编码成异常号（中断号、向量号），硬件拿它做下标查表：表项地址 = 表基址 + 异常号 × 表项宽度，一次访存就得到入口，CPU 不需要逐条询问设备。

慢的对照方案是轮询式。多个设备共享一条中断请求线时，CPU 进入公共处理程序，逐个读各设备的状态寄存器，读到哪家的请求标志置位就处理哪家。若共享 5 个设备，读一次状态寄存器算一次 20 个周期的总线事务，最坏情况下来源排在最后，仅确认来源就要 $5 \times 20 = 100$ 个周期，真正的处理还没开始；设备数增加，这笔开销线性上升。向量式让设备在中断响应周期直接把向量号送上总线（8086 系统中由 8259 中断控制器完成这个动作），确认来源的开销与设备数无关。

8051 是固定向量的例子：向量表从程序存储器地址 0 开始，每个入口间隔 8 字节——复位用 \code{0x0000}，外部中断 0 用 \code{0x0003}，定时器 0 用 \code{0x000B}，外部中断 1 用 \code{0x0013}，定时器 1 用 \code{0x001B}，串行口用 \code{0x0023}。8 字节只够放一条跳转指令加一点收尾，装不下完整处理程序，所以表项里通常就是一条跳往真正处理程序的跳转指令。8086 则是数组式向量表的完整形态：内存最低端 \code{0x00000}–\code{0x003FF} 的 1 KB 放 256 项，每项 4 字节（2 字节偏移加 2 字节段地址），中断号为 $n$ 的表项地址就是 $4n$。以 \code{int 21h} 为例，中断号 \code{0x21} 即十进制 33，$4 \times 33 = 132$，也就是 \code{0x84}，它的入口存在 \code{0x00084} 开始的 4 个字节里。

\begin{longtable}{L{0.14\linewidth}L{0.24\linewidth}L{0.26\linewidth}L{0.26\linewidth}}
\toprule
机器 & 表项定位 & 表项内容 & 组织特点 \\
\midrule
8051 & 固定地址，间隔 8 字节 & 通常是一条跳转指令 & 入口少；同一入口有多个来源时仍要软件再分辨 \\\tablerowrule
8086 & $4n$（$n$ 为中断号） & 偏移加段地址共 4 字节 & 256 项统一索引，软中断与硬件中断同表 \\\tablerowrule
教学模型 & 表基址 + 异常号 × 项宽 & 处理程序入口地址 & 本章异常入口沿用同一查表动作 \\
\bottomrule
\end{longtable}

把这条查表路径与两种真实组织画在同一张图上，查表动作就很具体了：

\begin{pdfdiagram}
  \node[hw state, minimum width=1.7cm, align=center] (vn) at (0,3.0) {向量号 n};
  \node[hw compute, minimum width=1.2cm, align=center] (mul) at (2.0,3.0) {$\times$ 项宽};
  \node[hw compute, minimum width=1.6cm, align=center] (addb) at (4.1,3.0) {$+$ 表基址};
  \node[hw memory, minimum width=1.8cm, align=center] (vt) at (6.5,3.0) {向量表};
  \node[hw state, minimum width=1.8cm, align=center] (npc) at (9.0,3.0) {PC ←\\入口地址};
  \draw[hw bus] (vn) -- (mul);
  \draw[hw bus] (mul) -- (addb);
  \draw[hw bus] (addb) -- node[hw label, above]{表项地址} (vt);
  \draw[hw bus] (vt) -- node[hw label, above]{查表} (npc);
  \node[hw memory, minimum width=4.8cm, align=center] (i51) at (2.4,0.6) {8051 固定入口\\复位 \code{0x0000}，外部中断 0 \code{0x0003}\\定时器 0 \code{0x000B}，每项间隔 8 字节};
  \node[hw memory, minimum width=4.8cm, align=center] (i86) at (7.9,0.6) {8086 IVT：\code{0x00000}--\code{0x003FF}\\256 项，中断号 n 的表项地址为 $4n$\\\code{int 21h}：$4\times33=132=$\code{0x84}};
  \draw[hw return] (vt.south) -- (6.5,1.9) -- (2.4,1.9) -- (i51.north);
  \draw[hw return] (vt.south) -- (6.5,1.55) -- (7.9,1.55) -- (i86.north);
\end{pdfdiagram}
\diagramnote{向量表的查表路径：向量号乘以表项宽度、加上表基址得到表项地址，一次访存取出入口地址装入 PC。8051 把少数入口固定在程序存储器低端、每项 8 字节只够放一条跳转指令；8086 在内存最低 1 KB 放 256 项、每项 4 字节，软中断与硬件中断共用同一张表、同一个下标规则。}

向量表把“谁请求”翻译成“去哪里”，此后的入口动作——保存现场、切换权限、阻止错误写回——对各类事件是公共的，这正是上一主题那张异常骨架图中“查向量表”一格的具体内容。还要留意一个边界：向量表本身也是内存，表项若被写坏，事件到来时 PC 会得到一个无意义地址。保护模式下的 x86 把向量表（IDT）的基址与界限交给特权寄存器管理，正是为这张表加上访问边界。

\topic{精确异常要求程序员可见状态像顺序执行}
流水线和乱序执行会让多条指令同时处在不同阶段。若较晚指令已经写回，较早指令随后发现异常，程序员可见状态就会变得难以恢复。精确异常要求异常发生时，状态看起来像按程序顺序执行到某条指令：异常之前的指令都已提交，异常指令和之后的指令都没有提交。教学 CPU 可以通过按序完成天然满足；现代 CPU 需要提交队列、重排序缓冲或类似机制。

把执行与提交画成两条时间轴，精确异常的边界位置就具体了：

\begin{pdfdiagram}
  % 精确异常边界
  \draw[hw return] (3.6,0.4) -- (3.6,3.3);
  \node[hw label, anchor=west] at (3.75,3.15) {精确异常边界：I3 报异常};
  % 执行完成行
  \node[hw label, anchor=east] at (-0.2,2.4) {执行};
  \draw[BookTeal, line width=0.9pt] (0,2.1) rectangle (1.2,2.7);
  \node[font=\small\sffamily] at (0.6,2.4) {I1};
  \draw[BookTeal, line width=0.9pt] (1.2,2.1) rectangle (2.4,2.7);
  \node[font=\small\sffamily] at (1.8,2.4) {I2};
  \draw[BookTeal, line width=0.9pt] (2.4,2.1) rectangle (3.6,2.7);
  \node[font=\small\sffamily] at (3.0,2.4) {I4};
  \draw[BookAccent, line width=0.9pt] (3.6,2.1) rectangle (4.8,2.7);
  \node[font=\small\sffamily] at (4.2,2.4) {I3 $\times$};
  \draw[BookTeal, line width=0.9pt] (4.8,2.1) rectangle (6.0,2.7);
  \node[font=\small\sffamily] at (5.4,2.4) {I5};
  % 提交行（严格按程序顺序）
  \node[hw label, anchor=east] at (-0.2,1.1) {提交};
  \draw[BookTeal, line width=0.9pt] (0,0.8) rectangle (1.2,1.4);
  \node[font=\small\sffamily] at (0.6,1.1) {I1};
  \draw[BookTeal, line width=0.9pt] (1.2,0.8) rectangle (2.4,1.4);
  \node[font=\small\sffamily] at (1.8,1.1) {I2};
  \draw[BookRule, line width=0.5pt, dashed] (2.4,0.8) rectangle (3.6,1.4);
  \node[font=\small\sffamily] at (3.0,1.1) {I3 未决};
  \draw[BookRule, line width=0.5pt, dashed] (3.6,0.8) rectangle (4.8,1.4);
  \node[font=\small\sffamily] at (4.2,1.1) {I4 弃};
  \draw[BookRule, line width=0.5pt, dashed] (4.8,0.8) rectangle (6.0,1.4);
  \node[font=\small\sffamily] at (5.4,1.1) {I5 弃};
  \node[hw label, anchor=west] at (7.4,1.1) {边界后一律不提交};
  % I4 先完成但必须排在 I3 之后
  \draw[hw return] (3.0,2.1) -- (3.0,1.4);
  \node[hw label, anchor=east] at (2.9,1.75) {I4 先完成仍按序等待};
\end{pdfdiagram}
\diagramnote{执行与提交是两条不同的时间轴。乱序机器里 I4 可能比 I3 先执行完成，但提交严格按程序顺序排队，I4 的结果只能先留在内部寄存器里等待；I3 在边界处报异常时，I1、I2 已经提交，I3 自身不提交，I4、I5 一律丢弃——程序员可见状态就像顺序执行到 I3 为止，这正是精确异常要求的形态。}

\begin{longtable}{L{0.18\linewidth}L{0.38\linewidth}L{0.34\linewidth}}
\toprule
机制 & 解决的问题 & 代价 \\
\midrule
按序流水线 & 后一条指令不越过前一条提交 & 吞吐有限，长延迟指令会拖住后续 \\\tablerowrule
提交队列 & 执行可乱序，提交按程序顺序 & 需要记录目的寄存器、异常原因和结果 \\\tablerowrule
寄存器重命名 & 消除部分假依赖，允许更多并行 & 需要映射表、空闲表和恢复机制 \\\tablerowrule
冲刷流水线 & 分支预测错或异常时丢弃错误路径 & 浪费已经做的工作，需要恢复 PC 和状态 \\
\bottomrule
\end{longtable}

虽然本书不展开乱序核心设计，但提前理解“执行”和“提交”的区别很重要。ALU 算出结果只是执行完成，结果进入程序员可见寄存器才是提交。加载指令从 cache 取回数据只是访存完成，目标寄存器被写回才是提交。存储更复杂：它可能先进入 store buffer，等确认没有异常、顺序允许、地址属性允许后才对 cache、内存或设备可见。这个区分把处理器核心相关章节 CPU trace、存储与设备事务相关章节内存顺序和经典芯片与平台案例平台异常连成一条线。

\topic{程序时间等于指令数乘 CPI 再乘时钟周期}
正确性边界讨论到此，接下来的问题是：这台机器跑一个程序有多快。评价它只需一个乘法：

\begin{lstlisting}
程序时间 = 指令数 x CPI x 时钟周期 = 指令数 x CPI / 主频
\end{lstlisting}

指令数由程序、编译器和 ISA 决定；CPI（每条指令的平均周期数）由微结构与访存行为决定；时钟周期由工艺和最长组合路径决定。三个因子各自独立地进入乘积，任何优化都必须先说清楚自己动的是哪一个。

把三个因子画成乘积结构，每个因子由什么决定就各有归属：

\begin{pdfdiagram}
  \node[hw state, minimum width=1.9cm, align=center] (f1) at (0,1.4) {指令数};
  \node[hw label] at (1.5,1.4) {$\times$};
  \node[hw state, minimum width=1.6cm, align=center] (f2) at (3.0,1.4) {CPI};
  \node[hw label] at (4.3,1.4) {$\times$};
  \node[hw state, minimum width=1.9cm, align=center] (f3) at (5.8,1.4) {时钟周期};
  \node[hw label] at (7.3,1.4) {$=$};
  \node[hw compute, minimum width=2.0cm, align=center] (tt) at (8.9,1.4) {程序时间};
  \node[hw label, align=center] at (0,0.3) {程序、编译器\\与 ISA 决定};
  \node[hw label, align=center] at (3.0,0.3) {微结构与\\访存行为决定};
  \node[hw label, align=center] at (5.8,0.3) {工艺与最长\\组合路径决定};
  \node[hw label, align=center] at (8.9,0.3) {优化效果的\\最终比较对象};
\end{pdfdiagram}
\diagramnote{程序时间的三因子分解。指令数由程序、编译器与 ISA 决定，CPI 由微结构与访存行为决定，时钟周期由工艺与最长组合路径决定；三个因子独立进入乘积，压低任何一个都可能抬高另一个。本节第一组演算中，同一程序 100 万条指令、周期同为 10 ns，CPI=1 的机器甲用 10 ms，CPI=4 的机器乙用 40 ms——时间差恰好等于 CPI 差。}

第一组演算：同一个程序（同一份二进制，100 万条指令）分别跑在两台 100 MHz 的机器上，时钟周期都是 10 ns。机器甲是理想单周期，CPI=1；机器乙是多周期，CPI=4。

\begin{longtable}{L{0.22\linewidth}L{0.12\linewidth}L{0.14\linewidth}L{0.38\linewidth}}
\toprule
机器 & CPI & 周期 & 程序时间 \\
\midrule
甲（单周期） & 1 & 10 ns & $10^6 \times 1 \times 10\text{ ns} = 10\text{ ms}$ \\\tablerowrule
乙（多周期） & 4 & 10 ns & $10^6 \times 4 \times 10\text{ ns} = 40\text{ ms}$ \\
\bottomrule
\end{longtable}

周期相同的前提下，时间差恰好等于 CPI 差：乙比甲慢 4 倍，多花的 30 ms 全部来自每条指令多用的 3 个周期。但不能反推出“CPI 低就一定快”。多周期设计敢于把周期做短，正是因为每个周期只做一小段：若甲按单周期实现，最长组合路径要求 40 ns 的周期，甲的实际时间变成 $10^6 \times 1 \times 40\text{ ns} = 40\text{ ms}$，与乙持平。这 40 ns 正是“从动作编码到控制字组织”一节五段模型的量化版本：取指、译码、ALU、访存、写回五段组合路径各约 8 ns，单周期必须在一拍内串行走完，多周期则把每一段各自安放进 10 ns 的短周期。只看 CPI 或只看主频都会误判，三个因子必须一起算。

第二组演算：cache 命中率对有效 CPI 的影响。cache 全命中时基础 CPI 为 1；设平均每条指令引起 1.2 次存储器访问（取指 1 次，约 20\% 的 load/store 指令再贡献 0.2 次），miss 率 2\%，每次 miss 罚 50 个周期，则：

\begin{lstlisting}
有效 CPI = 基础 CPI + 每条指令访存次数 x miss 率 x miss 代价
        = 1 + 1.2 x 0.02 x 50 = 1 + 1.2 = 2.2
\end{lstlisting}

98\% 的命中率听起来很高，却因为 50 周期的 miss 代价把有效 CPI 抬到 2.2，程序时间变成理想情形的 2.2 倍。miss 代价越大，命中率小数点后的每一位就越值钱，这也是存储层次要在本书后文独占一个完整部分的原因。

\begin{longtable}{L{0.18\linewidth}L{0.36\linewidth}L{0.34\linewidth}}
\toprule
提速途径 & 典型手段 & 付出的代价 \\
\midrule
减少指令数 & 更好的算法、编译器优化、表达力更强的 ISA & 单条指令更重，CPI 或周期可能变长 \\\tablerowrule
降低 CPI & 流水线、cache、旁路、分支预测 & 控制变复杂，冒险与精确异常都要额外机制 \\\tablerowrule
缩短周期 & 更深流水、更短组合路径、更先进工艺 & 级间寄存器开销增加，CPI 可能回升 \\
\bottomrule
\end{longtable}

三条途径共用同一个乘积，所以彼此牵制：复杂指令集用一条指令做更多事来减少指令数，往往抬高 CPI；深流水缩短周期，却让分支预测错误的冲刷更昂贵。本章的最小 CPU 在三条路上都走得很保守——指令数不省、CPI 接近 1、周期却必须包容最长路径。

\begin{longtable}{L{0.2\linewidth}L{0.42\linewidth}L{0.28\linewidth}}
\toprule
工具 & 回答的问题 & 常见误用 \\
\midrule
时延/吞吐拆分 & 用户等多久 vs 系统每秒做多少 & 用吞吐高掩盖时延长 \\\tablerowrule
总时间因子 & 指令数、CPI、周期各占多少 & 只优化占比小的因子 \\\tablerowrule
Amdahl 定律 & 优化一部分能换来多少整体加速 & 忘记未优化部分的拖累 \\\tablerowrule
roofline & 瓶颈在算力还是在带宽 & 拿峰值算力当可达性能 \\
\bottomrule
\end{longtable}

\begin{keyidea}
性能分析的铁律是“先测量，再优化，最后验证”。一切没有测量支持的优化都是猜测，一切不回测的优化都可能让系统更慢。
\end{keyidea}

\section{二、时延、吞吐与利用率}

\topic{时延与吞吐是两个不同的量}
时延（latency）是一个任务从发起到完成所经过的时间，单位是
“秒/条”——秒每指令、秒每请求、秒每次访问，回答“用户等多久”。
吞吐（throughput）是单位时间内完成的任务数，单位是“条/秒”——
指令每秒、请求每秒、字节每秒，回答“系统每秒做多少”。一个是单
个任务的视角，一个是系统管理者的视角。经典教材的说法是：在意
响应时间的用户该买更快的处理器，在意吞吐的管理员该买更多的处
理器。

\begin{longtable}{L{0.1\linewidth}L{0.32\linewidth}L{0.2\linewidth}L{0.24\linewidth}}
\toprule
量 & 定义 & 典型单位 & 回答的问题 \\
\midrule
时延 & 一个任务从发起到完成的时间 & 秒/条、秒/次、ns & 用户等多久 \\\tablerowrule
吞吐 & 单位时间完成的任务数 & 条/秒、MB/s & 系统每秒做多少 \\
\bottomrule
\end{longtable}

两个量的关系只在一个特殊情形下才简单：系统完全串行、无重叠、
无排队时，吞吐恰好等于时延的倒数。一旦允许重叠——流水线、并
行单元、批量传输——吞吐就可以脱离时延单独增长。把一件事切成
$n$ 段流水执行，单件事的时延不变，吞吐最多放大 $n$ 倍；把系统
原样复制 $n$ 份并排，吞吐再放大 $n$ 倍，单件事的时延还是不变。
这组不对称有物理根源：时延的下限由物理决定，信号要走完这段距
离、电平要穿过这些逻辑门，堆资源降不动；吞吐的上限主要由资源
决定，加并行度就能买。所以“更快”在工程上必须拆开问：是单件
更快，还是每秒更多。

四种组合因此全都存在，各配一个真实例子：

\begin{longtable}{L{0.17\linewidth}L{0.32\linewidth}L{0.37\linewidth}}
\toprule
组合 & 真实例子 & 为什么能这样搭配 \\
\midrule
快时延、高吞吐 & 乱序超标量服务器 CPU & 深流水加多发射：单条指令纳秒级完成，每秒退出上百亿条 \\\tablerowrule
快时延、低吞吐 & 同城专人直送快递 & 一单直达、半小时送到，但一人一次只能带一单，一天送不了几单 \\\tablerowrule
慢时延、高吞吐 & 远洋集装箱货轮 & 跨洋要一个月，一船装两万个箱子，折算每天的吞吐相当可观 \\\tablerowrule
慢时延、低吞吐 & 邮政平信 & 一封信走上好几天，一次也只送这一封 \\
\bottomrule
\end{longtable}

后三行在硬件里各有对应。中断控制器处理单个事件的时延很短，但
一次只处理一个事件，事件率一高就得排队，是“快时延、低吞吐”；
磁带冷备份定位一盘带要几十秒，持续写入却是每秒几百兆字节、一
夜能备份一个机房，是“慢时延、高吞吐”；机械电传打字机每个字
符都要上百毫秒、每秒也只有几个字符，是两头都不占。判断一个系
统属于哪一格，必须把两个量分别测出来——一个数字定不了位置。
也要提醒一句：给“专人直送”加人手能提高每天的单量，却缩短不了
任何一单的路程，这就是用并行度买吞吐、买不来时延的日常版本。

把这张表摆成坐标系，四格各居其位——两个量可以自由搭配，一格也不空：

\begin{pdfdiagram}
  % 四象限
  \draw[BookTeal, line width=0.8pt, fill=BookCodeBg] (1.0,2.6) rectangle (5.2,5.0);
  \node[align=center] at (3.1,3.8) {快时延 · 低吞吐\\专人直送快递\\中断控制器};
  \draw[BookTeal, line width=0.8pt, fill=BookCodeBg] (5.2,2.6) rectangle (9.4,5.0);
  \node[align=center] at (7.3,3.8) {快时延 · 高吞吐\\乱序超标量\\服务器 CPU};
  \draw[BookTeal, line width=0.8pt, fill=white] (1.0,0.2) rectangle (5.2,2.6);
  \node[align=center] at (3.1,1.4) {慢时延 · 低吞吐\\邮政平信\\电传打字机};
  \draw[BookTeal, line width=0.8pt, fill=white] (5.2,0.2) rectangle (9.4,2.6);
  \node[align=center] at (7.3,1.4) {慢时延 · 高吞吐\\远洋集装箱货轮\\磁带冷备份};
  % 坐标轴
  \draw[BookMuted, line width=0.6pt, ->] (1.0,-0.35) -- (9.6,-0.35);
  \node[hw label, anchor=north] at (5.3,-0.45) {吞吐：左低右高};
  \draw[BookMuted, line width=0.6pt, ->] (0.5,0.2) -- (0.5,5.1);
  \node[hw label, rotate=90, anchor=south] at (0.1,2.6) {时延：下慢上快};
\end{pdfdiagram}
\diagramnote{时延与吞吐的四象限：纵轴时延上快下慢，横轴吞吐左低右高。四格各有真实居民——快而少（专人直送）、快而多（超标量 CPU）、慢而多（集装箱货轮）、慢而少（邮政平信）。判断一个系统属于哪一格，两个量必须分别测出来，一个数字定不了位置。}

广告为什么只报一个数。原因之一是一个数字好传播，“每秒 10 亿
次”远比“在什么负载下、时延多少、吞吐多少”好记；原因之二是厂
商总是挑有利的那一个报：内存条把带宽印在包装上、把 CL 时延藏
进规格书深处，云服务报每秒请求数、不报第 99 百分位时延，手机
芯片报峰值算力、不报持续负载下的实测值。读完本节的读者应养成
一个反射：听到任何性能数字，先问它是时延还是吞吐，再问另一个
是多少、在什么负载下量的。

还有一点必须补在定义之后：时延本身也不是一个数。同一台机器上，
同类请求的时延会波动——存储与设备事务相关章节已经给出过来源：cache 命中与否、
DRAM 行命中与否、队列里排没排队。所以谈时延要谈分布，至少给
平均值加一个百分位。平均值会撒谎的很容易构造相应示例：100 次访问
里 99 次各花 20 ns、1 次花了 2.02 $\mu$s，平均值是整齐的
$4000/100=40$ ns，看上去一切正常，可那一次请求比平均值慢了
50 倍。云服务把第 99 百分位时延写进服务等级、而不是只写平均
值，就是因为用户记住的是最慢的那一次，不是平均的那一次。吞吐
没有这个问题：它本来就是总量对时间的平均，一个数基本够用——
这又是两个量性格不同的一处表现。

\topic{利用率接近 1 时延会发散}
总线仲裁器、内存控制器、中断处理程序，抽象以后是同一个东西：
一个服务台。请求以随机节奏到达，服务台以固定能力逐个处理。定
义利用率 $\rho=\text{到达率}/\text{服务率}$，即服务台有多忙。
$\rho<1$ 时系统跟得上，$\rho\geq1$ 时队列无限变长，这是常识。
不常识的部分是：即使 $\rho$ 离 1 还有距离，时延也已经开始非线
性膨胀，而且比直觉快得多。

排队论中最简单的 M/M/1 模型（泊松到达、指数服务时间、单服务台）
给出定量结论，本书只引用、不推导：

\begin{lstlisting}
平均排队等待 = 服务时间 x rho / (1 - rho)
平均停留时间 = 服务时间 x 1 / (1 - rho)   （等待 + 服务本身）
\end{lstlisting}

\begin{longtable}{L{0.13\linewidth}L{0.24\linewidth}L{0.25\linewidth}L{0.26\linewidth}}
\toprule
利用率 $\rho$ & 平均排队等待 & 平均停留时间 & 直观画面 \\
\midrule
50\% & 1 倍服务时间 & 2 倍服务时间 & 偶尔要等，等得不久 \\\tablerowrule
90\% & 9 倍服务时间 & 10 倍服务时间 & 大部分时间前面有队 \\\tablerowrule
99\% & 99 倍服务时间 & 100 倍服务时间 & 服务本身只占停留的零头 \\
\bottomrule
\end{longtable}

把三档利用率画成服务台前的队伍，表里的 1、9、99 倍就成了队伍的长度：

\begin{pdfdiagram}
  % 行一：rho = 50%
  \node[hw label, anchor=east] at (-0.25,3.4) {$\rho$=50\%};
  \draw[BookTeal, line width=0.7pt, fill=BookCodeBg] (0.3,3.1) rectangle (0.75,3.7);
  \node[hw compute, minimum width=1.3cm, minimum height=0.9cm] (sv1) at (1.8,3.4) {服务台};
  \draw[hw bus] (0.75,3.4) -- (sv1.west);
  \node[hw label, anchor=west] at (2.75,3.4) {等待 1 倍服务时间};
  % 行二：rho = 90%，9 个排队
  \node[hw label, anchor=east] at (-0.25,1.9) {$\rho$=90\%};
  \foreach \x in {0,...,8} {
    \draw[BookTeal, line width=0.7pt, fill=BookCodeBg] ({0.3+0.5*\x},1.6) rectangle ({0.75+0.5*\x},2.2);
  }
  \node[hw compute, minimum width=1.3cm, minimum height=0.9cm] (sv2) at (6.0,1.9) {服务台};
  \draw[hw bus] (4.75,1.9) -- (sv2.west);
  \node[hw label, anchor=west] at (6.95,1.9) {等待 9 倍服务时间};
  % 行三：rho = 99%，队伍望不到头
  \node[hw label, anchor=east] at (-0.25,0.4) {$\rho$=99\%};
  \foreach \x in {0,...,5} {
    \draw[BookTeal, line width=0.7pt, fill=BookCodeBg] ({0.3+0.5*\x},0.1) rectangle ({0.75+0.5*\x},0.7);
  }
  \node[hw label] at (3.75,0.4) {$\cdots$};
  \node[hw compute, minimum width=1.3cm, minimum height=0.9cm] (sv3) at (5.4,0.4) {服务台};
  \draw[hw bus] (4.1,0.4) -- (sv3.west);
  \node[hw label, anchor=west] at (6.35,0.4) {等待 99 倍服务时间};
\end{pdfdiagram}
\diagramnote{同一个服务台在三档利用率下的平均队列（小方块为排队请求，对应上表的平均排队等待，不含服务本身）：50\% 时平均 1 个在等，90\% 时 9 个，99\% 时 99 个。服务能力从一半用到九成，队伍却长了 8 倍——空闲时省下的能力存不起来，突发一到就只能排队消化。}

从 50\% 到 90\%，服务能力多用了不到一倍，平均等待却放大 9 倍；
从 90\% 到 99\%，再挤出一成能力，等待又放大 11 倍。代入具体数
字感受一次：设服务时间为 10 ns（一次行命中 DRAM 访问的量级），
到达率从服务率的一半提到九成，平均排队等待就从 10 ns 涨到 90
ns，平均停留从 20 ns 涨到 100 ns——服务本身一纳秒都没变慢，
慢的全是排队。直觉解释只有一句：服务台空闲时省下的能力存不起
来，空转的十分钟不能预支给下一波突发用；而到达是随机的，突发
必然出现。利用率越接近 1，系统越没有余量吸收突发，队伍一旦排
起来要很久才消化得掉，平均等待于是按 $1/(1-\rho)$ 发散。

把这张表画成曲线，发散的陡峭就成了形状而不是形容词：

\begin{pdfdiagram}
  % 横轴 rho：0-100%；纵轴平均停留/服务时间，对数刻度，y = 2*log10(倍数)
  \draw[BookMuted, line width=0.6pt, ->] (0,0) -- (10.2,0);
  \draw[BookMuted, line width=0.6pt, ->] (0,0) -- (0,4.95);
  \node[hw label, anchor=north] at (4.8,-0.62) {利用率 $\rho$};
  \node[hw label, rotate=90, anchor=south] at (-1.0,2.4) {平均停留 / 服务时间（对数）};
  \foreach \x/\t in {0/0, 1.92/20\%, 3.84/40\%, 5.76/60\%, 7.68/80\%, 9.6/100\%} {
    \draw[BookMuted, line width=0.5pt] (\x,0.06) -- (\x,-0.06);
    \node[hw label, anchor=north] at (\x,-0.1) {\t};
  }
  \foreach \y/\t in {0/1, 0.602/2, 1.398/5, 2.0/10, 2.602/20, 3.398/50, 4.0/100} {
    \draw[BookMuted, line width=0.5pt] (0.06,\y) -- (-0.06,\y);
    \node[hw label, anchor=east] at (-0.12,\y) {\t};
  }
  % M/M/1 曲线：平均停留 = 服务时间 / (1-rho)
  \draw[BookAccent, line width=1pt] plot[smooth] coordinates {
    (0,0) (1.92,0.194) (3.84,0.444) (4.8,0.602) (5.76,0.796) (6.72,1.046)
    (7.68,1.398) (8.16,1.648) (8.64,2.0) (9.12,2.602) (9.408,3.398) (9.504,4.0)
  };
  \node[hw label, anchor=west] at (5.9,1.62) {$1/(1-\rho)$};
  % 三个标记点：50% / 90% / 99%
  \draw[BookTeal, dashed, line width=0.5pt] (4.8,0) -- (4.8,0.602) -- (0,0.602);
  \draw[BookTeal, dashed, line width=0.5pt] (8.64,0) -- (8.64,2.0) -- (0,2.0);
  \draw[BookTeal, dashed, line width=0.5pt] (9.504,4.0) -- (0,4.0);
  \fill[BookRed] (4.8,0.602) circle (1.6pt);
  \fill[BookRed] (8.64,2.0) circle (1.6pt);
  \fill[BookRed] (9.504,4.0) circle (1.6pt);
  \node[hw label, anchor=east] at (4.62,1.0) {50\%：2 倍};
  \node[hw label, anchor=east] at (8.5,2.38) {90\%：10 倍};
  \node[hw label, anchor=east] at (9.35,4.38) {99\%：100 倍};
\end{pdfdiagram}
\diagramnote{M/M/1 模型下平均停留时间随利用率按 $1/(1-\rho)$ 发散，纵轴取对数。利用率 50\%、90\%、99\% 对应 2、10、100 倍服务时间；曲线越往右越陡，末段几乎直立——服务本身没变，涨的全是排队。}

工程含义可以写成一条纪律：任何共享服务台都要保持余量，稳态利
用率压在五到七成，不要拿峰值能力去对标平均负载。三条具体建议，
分别对应本书已经见过的三处硬件。

总线：仲裁与带宽按峰值负载的一倍半预留。存储与设备事务相关章节 DMA 周期窃取的
例子就是这种余量思想——软盘搬运只占总线约 6\% 的时间，其余
94\% 留给 CPU 取指，CPU 的访存时延才不会被搬运拖长。改用块模
式时搬运本身更快，却总线利用率逼近饱和，CPU 取指的排队时延立
刻沿上面的曲线上行——存储与设备事务相关章节说块模式要与周期窃取权衡，定量的
理由就是这张表。

内存控制器：它的“服务时间”本身不固定——行命中快、行冲突慢。
负载没有变、行命中率下降，等价于服务时间变长、$\rho$ 上升，排
队会突然恶化；有些机器“负载没怎么涨、响应突然变慢”，常见根因
就在这里。工程上有两个做法：监控请求队列深度，它比带宽百分比
更早报警；用写缓冲吸收突发写、给读请求更高优先级，因为读的后
面有一个正在停顿的 CPU 在等，写通常没有。

中断：中断到达是典型的随机事件，处理速率必须明显高于峰值到达
率，经验上至少留出一倍余量，否则一次中断风暴就把响应时延推上
发散曲线。处理程序只做最必要的事——读状态、记标志，重活推迟
到后台批量处理；网卡的中断合并（一次中断处理一批包）本质上是
把到达率除以批大小，直接降低 $\rho$，是同一规律的应用。

最后把“保持余量”这句话说透：余量买的不是平均速度，而是时延的
可预期性。$\rho$ 从 0.7 提到 0.9，吞吐只多了不到三成，平均等
待却从 2.3 倍服务时间涨到 9 倍——把三成的能力留作余量，等价于
把 $1/(1-\rho)$ 的值压在 10 以内。对终端用户来说，“每次都在
100 ns 内返回”往往比“平均 60 ns、偶尔 1 ms”值钱得多：前者来
自余量，后者来自把利用率拉满。批处理是光谱的另一极：把到达攒
成一批再服务，等效于服务率乘以批大小，代价是攒批本身增
加时延——吞吐与时延的交换在这里又出现了一次，只是换成了排队
的语言。也要声明适用范围：真实到达不总是泊松过程，服务时间也
不总是指数分布，M/M/1 的具体数字会随之变化；但“利用率逼近 1
则时延发散”这个单调性结论，对几乎所有到达与服务分布都成立，
这正是它能成为工程纪律的原因。

\topic{流水线提高吞吐不改变时延}
处理器核心相关章节用同一个 6 条指令的小程序把三种组织算过一遍：每段组合路
径 1 个时间单位时，单周期 30 单位、多周期 24 单位、理想流水 11
单位。当时的重点是“流水为什么快”；本节换一个问法：流水线里一
条指令自己的时延是多少？答案是不变——从取指到写回依然 5 级、
5 个单位，与多周期机的加载指令相同，与单周期机也相同。变快的不是任何一
条指令，而是“每拍完成一条”的稳态节奏。把两个量分列出来：

\begin{longtable}{L{0.13\linewidth}L{0.19\linewidth}L{0.21\linewidth}L{0.31\linewidth}}
\toprule
组织 & 单条指令时延 & 6 条程序总时间 & 稳态吞吐 \\
\midrule
单周期 & 5 单位 & 30 单位 & 0.20 条/单位 \\\tablerowrule
多周期 & 3 至 5 单位 & 24 单位 & 0.25 条/单位 \\\tablerowrule
理想流水 & 5 单位 & 11 单位 & 短程序 $6/11\approx0.55$，长程序趋近 1 \\
\bottomrule
\end{longtable}

把单周期与五级流水画到同一条时间轴上，“时延不变、吞吐改变”就直接可见：

\begin{waveform}[x=0.72cm,y=0.85cm]
  % 上：单周期，每条指令独占 5 个单位
  \node[hw label, anchor=east] at (-0.2,3.8) {单周期};
  \draw[BookTeal, line width=0.9pt, fill=BookCodeBg] (0,3.4) rectangle (5,4.2);
  \node at (2.5,3.8) {I1};
  \draw[BookTeal, line width=0.9pt, fill=BookCodeBg] (5,3.4) rectangle (10,4.2);
  \node at (7.5,3.8) {I2};
  \draw[BookMuted, line width=0.6pt, <->] (0,3.05) -- (5,3.05);
  \node[hw label, anchor=north] at (2.5,2.97) {单条时延 5 单位};
  \node[hw label, anchor=west] at (10.25,3.8) {每 5 单位完成 1 条};
  % 下：五级流水，I1-I6 按级错位一格
  \node[hw label, rotate=90, anchor=center] at (-1.15,0.975) {五级流水};
  \node[hw label, anchor=east] at (-0.2,2.175) {IF};
  \node[hw label, anchor=east] at (-0.2,1.575) {ID};
  \node[hw label, anchor=east] at (-0.2,0.975) {EX};
  \node[hw label, anchor=east] at (-0.2,0.375) {MEM};
  \node[hw label, anchor=east] at (-0.2,-0.225) {WB};
  \foreach \c/\xa/\xb in {I1/0/1, I2/1/2, I3/2/3, I4/3/4, I5/4/5, I6/5/6} {
    \draw[BookTeal, line width=0.7pt, fill=BookCodeBg] (\xa,1.9) rectangle (\xb,2.45);
    \node at ({\xa+0.5},2.175) {\c};
  }
  \foreach \c/\xa/\xb in {I1/1/2, I2/2/3, I3/3/4, I4/4/5, I5/5/6, I6/6/7} {
    \draw[BookTeal, line width=0.7pt, fill=BookCodeBg] (\xa,1.3) rectangle (\xb,1.85);
    \node at ({\xa+0.5},1.575) {\c};
  }
  \foreach \c/\xa/\xb in {I1/2/3, I2/3/4, I3/4/5, I4/5/6, I5/6/7, I6/7/8} {
    \draw[BookTeal, line width=0.7pt, fill=BookCodeBg] (\xa,0.7) rectangle (\xb,1.25);
    \node at ({\xa+0.5},0.975) {\c};
  }
  \foreach \c/\xa/\xb in {I1/3/4, I2/4/5, I3/5/6, I4/6/7, I5/7/8, I6/8/9} {
    \draw[BookTeal, line width=0.7pt, fill=BookCodeBg] (\xa,0.1) rectangle (\xb,0.65);
    \node at ({\xa+0.5},0.375) {\c};
  }
  \foreach \c/\xa/\xb in {I1/4/5, I2/5/6, I3/6/7, I4/7/8, I5/8/9, I6/9/10} {
    \draw[BookTeal, line width=0.7pt, fill=BookCodeBg] (\xa,-0.5) rectangle (\xb,0.05);
    \node at ({\xa+0.5},-0.225) {\c};
  }
  \draw[BookMuted, line width=0.6pt, <->] (0,-0.85) -- (5,-0.85);
  \node[hw label, anchor=north] at (2.5,-0.95) {单条时延同样 5 单位};
  \node[hw label, anchor=west] at (10.25,0.975) {稳态每 1 单位完成 1 条};
  % 时间轴
  \draw[BookMuted, line width=0.6pt, ->] (0,-1.75) -- (10.6,-1.75);
  \node[hw label, anchor=north] at (5.3,-1.85) {时间（时间单位）};
\end{waveform}
\diagramnote{单周期与五级流水执行同一指令流的对照。两种组织里，一条指令从取指到写回都要 5 个时间单位；不同的是完成节奏——单周期每 5 个单位退出 1 条，流水灌满后每 1 个单位退出 1 条。流水线买到的是吞吐，不是单条时延。图中画的是无冒险的理想情形，6 条指令 10 个单位走完；上表“理想流水 11 单位”另含 1 拍 load-use 停顿，比图示多出的 1 个单位正是它。}

流水的吞吐在短程序上还没灌满，程序越长越接近每单位 1 条；但无
论灌没灌满，单条时延始终是 5 个单位。写成公式：

\begin{lstlisting}
单条指令时延 ≈ 级数 x 级延迟
稳态吞吐     ≈ 1 / 级延迟        （无冒险的理想情形）
\end{lstlisting}

两个式子里只有一个含级数。把 5 级改成 10 级、级延迟减半：单条
时延 $10\times0.5=5$，原地不动；吞吐从每单位 1 条升到 2 条。这
就是“加深流水”真正买到的东西——不是单条更快，而是每秒退出的
指令更多。广告里的“20 级流水线”从不承诺你的某一条 LOAD 更快，
它承诺的是吞吐。

真实机器里深流水甚至让单条更慢。每级之间要插流水寄存器，设其
开销为 0.05 个单位：10 级流水的级延迟是 $0.5+0.05$，单条时延变
成 $10\times0.55=5.5$，比 5 级还多出一成；吞吐也不是理想的 2，
而是 $1/0.55\approx1.82$。级数越深，寄存器开销占比越大，边际收
益越薄，这是流水线深度不能无限加的第一个原因。第二个原因第四
章已经见过：级数越深，分支猜错时要冲刷的预取指令越多，有效
CPI 回升，使理论吞吐降低——“CPI 与指令级性能”一节会把这个代价写进公式。
商用 CPU 的流水深度在本世纪初曾一路冲到 30 级上下，随后又退回
十几级，正是这两笔账共同裁决的结果：深度换来的吞吐增量，盖不
过开销与冲刷带来的损失。

公式里还有一个容易读漏的下标：吞吐 ≈ 1/级延迟，其中的级延迟
取的是最慢的一级。同步流水的所有级共用同一个时钟，周期必须容
纳最差的一级——设五级的真实延迟分别是 1.0、1.0、1.2、1.0、
1.0 个单位，周期只能按 1.2 走，其余四级各浪费 0.2。此时单条时
延是 $5\times1.2=6.0$，比五级延迟之和 5.2 还多；吞吐是
$1/1.2\approx0.83$，而不是按平均延迟算出来的 $1/1.04\approx0.96$。不平衡同时伤
害两个量：时延被级数乘上最慢级，吞吐被最慢级单独封顶。这就是
处理器核心相关章节强调阶段平衡的原因，也是真实 CPU 把访存这种“天然慢”的
环节拆得更细、或干脆允许它独占多拍（多级流水化 cache）的原因
——把最慢的一级切小，比把任何快一级做得更快都值钱。

\topic{案例：教学总线的一次访问时延}
存储与设备事务相关章节把一次总线读事务拆成地址、等待、数据三段，并用时序预算
决定要不要插入等待状态。现在给这台教学总线定一组数值：时钟周
期 10 ns，一次读事务固定 4 拍——第 1 拍驱动地址与片选，第 2、
3 拍是等待状态（慢存储器的访问时间经 READY/WAIT 机制插入），
第 4 拍发起者采样数据。一次读的时延从地址生效算到数据被采样，
连续读的吞吐上限按“每个事务独占 4 拍、互不重叠”计算：

\begin{lstlisting}
读时延           = 4 拍 x 10 ns = 40 ns
连续读吞吐上限   = 1 次 / 40 ns = 2.5 x 10^7 次/秒 = 25M 次/秒
按 32-bit 折算   = 25M 次/秒 x 4 B = 100 MB/s
\end{lstlisting}

这组数字是同一台总线的两张面孔，用两个改动可以分别验证它们谁
动、谁不动：

\begin{longtable}{L{0.26\linewidth}L{0.15\linewidth}L{0.2\linewidth}L{0.17\linewidth}}
\toprule
访问形态 & 单次时延 & 连续吞吐 & 32-bit 带宽 \\
\midrule
基本事务（4 拍） & 40 ns & 25M 次/秒 & 100 MB/s \\\tablerowrule
多插 1 拍等待（5 拍） & 50 ns & 20M 次/秒 & 80 MB/s \\\tablerowrule
4 字突发（4+3 拍） & 首字 40 ns & 平均 17.5 ns/字 & 约 229 MB/s \\\tablerowrule
长突发极限 & 首字 40 ns & 10 ns/字 & 400 MB/s \\
\bottomrule
\end{longtable}

多插一拍等待，时延涨 25\%、吞吐跌 20\%，两个量同向恶化，因为
它们共享“事务总长”这个分母——这是“慢”的通常情形。突发传输
则是另一个方向的操作：首字时延一纳秒没省，平均吞吐却翻到接近
4 倍。存储与设备事务相关章节 DRAM 的行缓冲突发、cache 的整行填充，用的全是这
一种方法是：把每次访问的固定开销摊到一串数据上。摊薄固定开销提高的
是吞吐，受益的也只是第 2 个字起的时延；第 1 个字永远要等满 40
ns——首字时延由存储器的访问时间决定，和后面排多少数据无关。

写事务值得顺带一提：写可以把数据先丢进写缓冲就返回，时延被缓
冲藏住，只有缓冲写满时才暴露真实事务时延——存储与设备事务相关章节的写缓冲正
是利用“写后面没有人在等”这一点。最后补一条工程边界：等待拍
不能无限插下去，存储与设备事务相关章节强调过事务要有最大等待长度与超时返回，
否则一个不存在的目标会把整机永久挂在第一拍上——时延分析的前
提是时延有限。

教学总线还有一个故意不做的功能值得点名：地址与数据的解耦。现
代高性能总线允许事务拆分——地址拍发出后总线让给别人，数据准
备好后再占用总线送回——等待的时间不再独占总线，吞吐上限从
“1/事务总长”逼近“1/地址拍”。代价是每笔事务要带标签、目标
要能乱序完成，控制器复杂度远超教学总线的 READY/WAIT。本书的
总线停在教学模型是对的；但将来读真实总线的规格书时要记得，
“4 拍事务、25M 次/秒”描述的是总线的一种用法，不是总线的天花
板。

\topic{案例：教学计算机项目的时延与吞吐}
教学计算机项目中的面包板计算机是时延与吞吐的极端样本。它的时钟自动档约
1 Hz——面包板加人眼给频率划了上限，LED 每拍都要看得清；它的
T 状态计数器是模 8 的，没有提前回绕逻辑，每条指令固定走满 T0
到 T7 八拍。于是：

\begin{lstlisting}
单条指令时延 = 8 拍 x 1 秒 = 8 秒/条
吞吐         = 1 条 / 8 秒 = 0.125 条/秒
\end{lstlisting}

本书末章的教学计算机项目的数灯程序共执行 41 条指令，其中 HLT 在 T2 冻结、只占 3 拍，其余 40 条各 8 拍，共 $40\times8+3=323$ 拍，在这台
机器上要跑 323 秒，五分半钟。对照组：一颗 5 GHz、平均 IPC（instructions per cycle，CPI 的倒数）为 2
的真实 CPU 完成同样 41 条指令约需 21 个周期、4 纳秒。吞吐相差
将近 11 个数量级，时延相差 9 个数量级还多——本书最小与最快的
两台机器之间，隔着整个处理器工业几十年的优化史。

这台串行、无重叠的教学机也顺手验证了本节开头的命题：无重叠时
吞吐恰是时延的倒数，$1/8=0.125$。若把它的时钟从 1 Hz 提到 1
MHz（电路完全撑得住，只是人眼再也跟不上），时延变为 8 微秒、
吞吐变为 12.5 万条/秒——两个量按同一倍数变化。要打破这种同进
同退，只有引入重叠：流水、突发、并行，全是本书末章的教学计算机项目之后各章的主
题。

其实本书末章的教学计算机项目已经把这个旋钮做进了硬件：时钟分自动与单步两路，自
动档还能调速。往上调，时延与吞吐按比例变好，LED 先是看不清、
再是只剩残影，最后只剩调试器才知道机器在干什么——可观察性与
性能的交换是连续的，不是二选一。教学机把这个连续旋钮摆在面板
上，本身就是对性能工程最直白的一课：所有的快都有代价，区别只
在以什么形式支付。

教学机慢不是失败，它的设计目标本来就写在别处。慢到 1 Hz，是为
了让每一拍可观察：总线值、控制线、T 状态都能用 LED 逐拍核对；
固定 8 拍，是为了让微码表保持“地址 = opcode 乘 8 加 T”这一条
规律，省掉按指令提前结束的例外逻辑；单步按钮则把时延直接交到
人手里，按一下才走一拍。时延与吞吐被刻意牺牲，换来的是“没有
一处不可观察”。真实 CPU 的方向正相反：把每拍做短、把多条指令
的拍重叠起来、把访存等待藏起来——分别对应处理器核心相关章节的流水线、第
五章的 cache 与预取，以及本书未展开的乱序与多发射。读任何性能
数字都要连同设计目标一起读：在教学机的指标表上，可观察性排在
快慢前面。

\section{三、CPI 与指令级性能}

\topic{总时间三因子的完整展开}
处理器核心相关章节给出的基本乘法是本章一切计算的出发点：

\begin{lstlisting}
程序时间 = 指令数 x CPI x 时钟周期 = 指令数 x CPI / 主频
\end{lstlisting}

三个因子连乘，意味着每个因子都是总时间的一个完整维度：任何一
个翻一倍，总时间翻一倍；任何一个（假想着）降到零，总时间降到
零。同样重要的是，三个因子各有各的决定者，几乎没有哪个手段能
只动其中一个而不惊动另外两个：

\begin{longtable}{L{0.11\linewidth}L{0.3\linewidth}L{0.25\linewidth}L{0.22\linewidth}}
\toprule
因子 & 由什么决定 & 典型优化手段 & 常见的反向移动 \\
\midrule
指令数 & 算法、编译器、ISA 的表达力 & 更好的算法与编译优化、更高效的指令 & 单条指令变重，CPI 或周期变长 \\\tablerowrule
CPI & 微结构（流水、cache、旁路、预测）与程序行为 & 加 cache、改进预测、乱序执行 & 控制变复杂，面积、功耗、验证成本上升 \\\tablerowrule
周期 & 最长组合路径、工艺、电压与温度 & 切短关键路径、加深流水、先进制程 & 级间寄存器开销增加，CPI 回升 \\
\bottomrule
\end{longtable}

逐行看这张表。指令数一行说的是“程序要机器干多少活”：同一个
C 程序，关优化和开优化编译出的动态指令数可以差好几倍；同一源
码编译到两套 ISA 上，指令数也不同——所以比较指令数必须固定
编译器与 ISA，否则比的是软件不是机器。注意这里的指令数一律指
动态执行的条数：循环体里的一条加法执行几百万次就计几百万次，
静态代码里写了多少行不作数。CPI 一行说的是“机器干每条活平均
花几拍”，它一半由微结构决定（流水线多深、cache 多大、预测多
准），另一半由程序行为决定（访存多不多、分支好不好猜）——后
一半正是本节后文要逐项算的。周期一行说的是“一拍有多长”，由
最慢的那条组合路径封顶：处理器核心相关章节的单周期机就是反面教材，为了容
纳加载指令的五段串行路径，每条指令都得按 5 个单位的周期走。

把三因子画成一条乘法链会更直观：每个因子挂着自己的决定者，乘积才是程序时间。

\begin{pdfdiagram}[x=0.85cm]
  \node[hw state, minimum width=2.3cm, minimum height=1.2cm] (f1) at (0.9,0.4) {指令数\\$2\times10^9$};
  \node[hw control, minimum width=2.3cm, minimum height=1.2cm] (f2) at (4.3,0.4) {CPI\\1.5};
  \node[hw compute, minimum width=2.3cm, minimum height=1.2cm] (f3) at (7.4,0.4) {时钟周期\\2 ns};
  \node[hw compute, minimum width=2.0cm, minimum height=1.2cm, draw=BookRed] (res) at (10.6,0.4) {程序时间\\6 s};
  \draw[hw bus] (f1.east) -- (f2.west);
  \draw[hw bus] (f2.east) -- (f3.west);
  \draw[hw bus] (f3.east) -- (res.west);
  \node[hw label] at (2.6,0.78) {$\times$};
  \node[hw label] at (5.75,0.78) {$\times$};
  \node[hw label] at (9.05,0.78) {$=$};
  % 各因子的决定者（虚线下挂）
  \draw[hw return, dashed] (f1.south) -- (0.9,-0.35);
  \node[hw label, align=center] at (0.9,-0.8) {算法、编译器、\\ISA 表达力};
  \draw[hw return, dashed] (f2.south) -- (4.3,-0.35);
  \node[hw label, align=center] at (4.3,-0.8) {微结构\\+ 程序行为};
  \draw[hw return, dashed] (f3.south) -- (7.4,-0.35);
  \node[hw label, align=center] at (7.4,-0.8) {最长组合路径、\\工艺与电压};
\end{pdfdiagram}
\diagramnote{程序时间三因子连乘：指令数由算法、编译器与 ISA 决定，CPI 由微结构与程序行为共同决定，时钟周期被最长组合路径封顶。数值取自本节的完整演算：$2\times10^9$ 条 $\times$ 1.5 $\times$ 2 ns $=$ 6 s。任何优化手段都要先说清动的是哪个因子、另外两个动不动。}

“优化必须说清动的是哪个因子”不是形式主义，因为三个因子互相
牵制，压下一个常常顶起另一个。两个经典例子。其一是 ISA 之争：
复杂指令用一条指令干更多事，指令数降了，每条却更重、CPI 更高；
精简指令把单条做轻，CPI 接近 1，指令数又升回去——净效果只有
算完乘积才知道，这正是上世纪 RISC 之争最终靠实测而不是靠口号
裁决的原因。其二是深流水：级数加深把周期切短，分支冲刷却更贵、
流水寄存器更多，CPI 回升，“时延、吞吐与利用率”一节已经算过寄存器开销那笔账。再
给一个假想的评审案例：某编译器优化宣称“动态指令数减少 30\%”，
但新指令序列让平均 CPI 从 1.0 涨到 1.5——乘积
$0.7\times1.5=1.05$，总时间不降反升 5\%。单看任何一个因子的
“改善”都不构成优化，这是本章第一条纪律；任何一个优化提案，
评审时只问一句：另外两个因子动不动、动多少，乘积降不降。

三个因子凑齐，先完整算一遍再往下拆：某程序动态执行
$2\times10^9$ 条指令，机器平均 CPI 为 1.5，主频 500 MHz（周期
2 ns）：

\begin{lstlisting}
程序时间 = 2 x 10^9 条 x 1.5 x 2 ns = 6 x 10^9 ns = 6 s
\end{lstlisting}

6 秒这个结果拆得毫无悬念：$2\times10^9\times1.5=3\times10^9$
拍，每拍 2 ns。有价值的是反过来的问法——要把它跑进 2 秒，有
三条独立的路：指令数降到三分之一（换算法或编译器）、CPI 降到
0.5（换微结构或换程序形态）、周期降到 0.67 ns（换工艺或重切
关键路径）。三条路的难度天差地别，定量框架的作用就是先告诉你
每条路要降多少，再去评估哪条走得通。再给一个编译器动指令数的
实感：同一程序开优化后动态指令数从 $2\times10^9$ 降到
$1.2\times10^9$，若 CPI 不变，6 秒直接变成 3.6 秒——编译器一
分钱硬件没动，只动了第一个因子。这也是对比机器时必须固定编译
选项的原因：你测的是机器，不是编译器。

关于周期这个因子还要补一段史实：2005 年前后，商用 CPU 的主频
在几 GHz 处撞上了功耗墙——频率再往上，功耗与散热先撑不住，
“周期逐年减半”的时代就此结束。性能增长的主战场被迫转向另外
两个因子：用更宽的微结构压 CPI，用更多的核分摊指令数。本章后
面谈多核的 Amdahl 极限、谈功耗约束下的设计取舍，根子都在这一
段历史上。

\topic{CPI 按指令类型加权}
CPI 的全称是每条指令的平均周期数，“平均”必须落实到具体程序
上：同一台 CPU 跑不同程序，指令混合不同，平均 CPI 就不同。计
算方法是按指令类型加权——每类指令的动态占比乘上该类自己的
CPI，再求和。权重用动态条数（执行了多少次），不用静态条数
（程序里写了多少行）。

取一个有代表性的整数程序混合，逐类给出该类 CPI，加权求和：

\begin{longtable}{L{0.17\linewidth}L{0.16\linewidth}L{0.15\linewidth}L{0.32\linewidth}}
\toprule
指令类型 & 动态占比 & 该类 CPI & 对平均 CPI 的贡献 \\
\midrule
ALU 运算 & 40\% & 1 & $0.40\times1=0.40$ \\\tablerowrule
加载 & 20\% & 2 & $0.20\times2=0.40$ \\\tablerowrule
存储 & 10\% & 2 & $0.10\times2=0.20$ \\\tablerowrule
分支 & 20\% & 1.5 & $0.20\times1.5=0.30$ \\\tablerowrule
乘除 & 10\% & 8 & $0.10\times8=0.80$ \\
\midrule
合计 & 100\% & --- & 平均 CPI $=2.1$ \\
\bottomrule
\end{longtable}

\begin{lstlisting}
CPI = 0.40x1 + 0.20x2 + 0.10x2 + 0.20x1.5 + 0.10x8
    = 0.40 + 0.40 + 0.20 + 0.30 + 0.80
    = 2.1
\end{lstlisting}

把五个贡献按长度画成一条带，“条数占比”与“时间占比”的错位就直接可见：

\begin{pdfdiagram}
  % 平均 CPI 2.1 按指令类型的贡献，条带全长 10 cm
  \draw[BookTeal, line width=0.8pt, fill=BookCodeBg] (0,0) rectangle (1.905,1.0);
  \node at (0.9525,0.5) {ALU};
  \draw[BookTeal, line width=0.8pt, fill=BookCodeBg] (1.905,0) rectangle (3.81,1.0);
  \node at (2.8575,0.5) {加载};
  \draw[BookTeal, line width=0.8pt, fill=BookCodeBg] (3.81,0) rectangle (4.762,1.0);
  \node at (4.286,0.5) {存储};
  \draw[BookTeal, line width=0.8pt, fill=BookCodeBg] (4.762,0) rectangle (6.19,1.0);
  \node at (5.476,0.5) {分支};
  \draw[BookAmber, line width=0.8pt, fill=BookAmber!18] (6.19,0) rectangle (10,1.0);
  \node at (8.095,0.5) {乘除};
  \node[hw label, anchor=south west] at (0,1.14) {对平均 CPI 的贡献（总长 = 2.1）};
  \node[hw label, anchor=south] at (8.095,1.14) {0.80 · 占时间 38\%};
  % 每段的 占比 x 该类 CPI
  \node[hw label, anchor=north] at (0.9525,-0.12) {40\%$\times$1};
  \node[hw label, anchor=north] at (2.8575,-0.12) {20\%$\times$2};
  \node[hw label, anchor=north] at (4.286,-0.12) {10\%$\times$2};
  \node[hw label, anchor=north] at (5.476,-0.12) {20\%$\times$1.5};
  \node[hw label, anchor=north] at (8.095,-0.12) {10\%$\times$8};
  % 刻度
  \foreach \x/\t in {0/0, 4.762/1.0, 9.524/2.0, 10/2.1} {
    \draw[BookMuted, line width=0.5pt] (\x,-0.55) -- (\x,-0.75);
    \node[hw label, anchor=north] at (\x,-0.82) {\t};
  }
  \node[hw label, anchor=north] at (5,-1.45) {乘除只占条数 10\%，却贡献约 38\% 的执行时间：占比小、单类慢，照样主导平均值};
\end{pdfdiagram}
\diagramnote{平均 CPI 2.1 按指令类型拆成五段：ALU 与加载各贡献 0.40，存储 0.20，分支 0.30，乘除独占 0.80。乘除只占动态条数的 10\%，却占执行时间约 38\%——条数占比与时间占比的错位从条带上直接读出，这正是“先弄清时间花在哪”的定量形态。}

这张表里最值得停一秒的是最后一行：乘除指令只占条数的 10\%，
却贡献了 $0.8/2.1\approx38\%$ 的执行时间。占比小、单类慢，照
样能主导平均值——这是“先弄清时间花在哪”的定量版本，也是下
一节 Amdahl 定律的主题。在这个混合上比较两个候选优化：把乘除
单元加快一倍（该类 CPI 从 8 降到 4），省下 $0.4/2.1\approx19\%$
的总时间；把占条数最多的 ALU 加快一倍（CPI 从 1 降到 0.5），
只省 $0.2/2.1\approx10\%$。条数最多不等于最该优化，加权表把
这句话变成了算术。

反过来，换一个几乎没有乘除的程序（比如纯控制流代码），把其余
四类按比例归一，同一台机器的平均 CPI 立刻降回 1.4 上下。CPI
是“机器乘程序”的性质，从来不是硬件常数——处理器核心相关章节章末练习里
“CPI 是动态混合的平均”那句话，到这里有了完整算法。工程上还
要知道权重从哪来：动态占比靠指令级模拟器、插桩或性能计数器统
计得到，这正是第五节测量工具要回答的第一类问题。

权重的“动态”二字值得用反例再强调一次。某程序静态代码 1000
条指令，乘除只占 5 条（0.5\%），单看静态混合会以为乘除无关紧
要；可这 5 条全部在最内层循环里，动态占比 10\%——正是上表那
个混合。拿静态占比加权会算出完全不同的 CPI，进而错判优化方向。
静态条数回答“程序写了什么”，动态条数回答“机器跑了什么”，性
能计算只认后者。

\topic{有效 CPI 要加上 memory 与分支代价}
上一张表的“该类 CPI”还不是全部真相：它默认每次访存都命中、
每个分支都猜对。真实机器要把存储系统和分支预测的失效率折算进
去，得到有效 CPI：

\begin{lstlisting}
有效 CPI = 基础 CPI
         + 每条指令访存次数 x miss 率 x miss 代价
         + 每条指令分支数 x 误预测率 x 冲刷代价
\end{lstlisting}

形式上与处理器核心相关章节的 cache 演算一脉相承，只是把分支项一并写进来。
给一组完整数值：基础 CPI 取 1.0（理想流水、全命中、全猜对）；
每条指令平均访存 1 次（取指与数据访问合并，按此简化；处理器核心相关章节按
1.2 次算过更细的版本），cache miss 率 2\%，每次 miss 罚 50 拍；
分支占指令流的 25\%，误预测率 20\%——折合每条指令摊到 5\%
次误测——每次误测冲刷 2 拍。

\begin{lstlisting}
存储项   = 1 x 0.02 x 50      = 1.0
分支项   = 0.25 x 0.20 x 2    = 0.1
有效 CPI = 1.0 + 1.0 + 0.1    = 2.1
有效 IPC = 1 / 2.1            ≈ 0.48
\end{lstlisting}

拆开看占比：基础执行只占 48\%，存储系统占 48\%，分支占
5\%。一台“IPC 设计值 1”的机器实测 0.48——IPC 曲线总比理想
低，原因就在这里：理想值假设失效率为零，而每个失效事件都带一
个大乘数。98\% 的命中率听上去接近完美，可 50 拍的 miss 代价让
那 2\% 变成整整 1.0 个 CPI。敏感性也算一笔：把 miss 率从 2\%
压到 1\%，有效 CPI 从 2.1 降到 1.6，程序快 31\%——命中率小数
点后面每一位都明码标价，这就是存储与设备事务相关章节用一整章讲存储层次的性能
理由。分支项同理：深流水机的冲刷代价不是 2 拍而是 4、5 拍，同
样的误预测率立刻把分支项放大两三倍，机器对预测器的依赖随流水线
深度一起加深。超标量也逃不开同一公式：理想 IPC 为 4 的机器，基
础 CPI 只有 0.25，失效项的相对权重天然更大——每条指令 0.3 次数
据访存、miss 率 2\%、miss 代价 50 拍，存储项就是 0.30；分支占
两成、误预测率 4\%、每次罚 15 拍，分支项再加 0.12；有效 CPI 约
0.67，IPC 只剩约 1.5——基础越宽，失效项越显眼。
工程读法两条：其一，降失效率比继续堆基础宽度便宜，因为大头在
失效项；其二，报性能必须连失效率一起报——只说 IPC 0.48，不说
miss 率和误预测率，既没解释机器为什么是这个速度，也没给下一步
指出该动哪一项。

把这三项按占比画成一条带，“一半去了哪”比公式更直观：

\begin{pdfdiagram}
  % 有效 CPI = 2.1：基础 1.0 + 存储 1.0 + 分支 0.1，全长 10 cm
  \draw[BookTeal, line width=0.8pt, fill=BookCodeBg] (0,0) rectangle (4.76,1.0);
  \node at (2.38,0.5) {基础执行 1.0 · 48\%};
  \draw[BookAmber, line width=0.8pt, fill=white] (4.76,0) rectangle (9.52,1.0);
  \node at (7.14,0.5) {存储代价 1.0 · 48\%};
  \draw[BookRed, line width=0.8pt, fill=white] (9.52,0) rectangle (10,1.0);
  \node[hw label, anchor=south west] at (0,1.12) {有效 CPI 2.1 的构成（每指令周期数）};
  \node[hw label, anchor=south east] at (10,1.12) {分支 0.1 · 5\%};
  \foreach \x/\t in {0/0, 4.76/1.0, 9.52/2.0, 10/2.1} {
    \draw[BookMuted, line width=0.5pt] (\x,0) -- (\x,-0.09);
    \node[hw label, anchor=north] at (\x,-0.13) {\t};
  }
  \node[hw label, anchor=north] at (5,-0.68) {失效项合计 1.1，使 IPC 降低约一半：$1/2.1\approx0.48$};
\end{pdfdiagram}
\diagramnote{有效 CPI 2.1 按来源切成三段：基础执行与存储代价各占 48\%，分支占 5\%。设计值 IPC 1.0 的机器实测约 0.48，差值几乎全由失效项贡献——98\% 命中率的存储系统，占用的周期与执行本身一样多。优化预算该拨给哪一项，从条带上直接读。}

把三项的占比与出处列成表，本书前后每一章的性能话题都能在其中
找到自己的格子：

\begin{longtable}{L{0.13\linewidth}L{0.12\linewidth}L{0.16\linewidth}L{0.43\linewidth}}
\toprule
项 & 数值 & 占有效 CPI & 对应的硬件机制（章节） \\
\midrule
基础执行 & 1.0 & 48\% & 数据通路与理想流水（处理器核心相关章节） \\\tablerowrule
存储代价 & 1.0 & 48\% & cache、存储层次与行调度（存储与设备事务相关章节） \\\tablerowrule
分支代价 & 0.1 & 5\% & 分支预测与冲刷控制（处理器核心相关章节） \\
\bottomrule
\end{longtable}

这张表还有两种用法。设计早期没有机器可测，就用它做纸面估算：
替候选参数各算一遍有效 CPI，谁的乘积小选谁，把昂贵的仿真留给
入围方案。机器造出来以后，表里的失效率不再靠假设——miss 率、
误预测率都能用性能计数器直接读出，那是第五节的主题；读出来以后
代回本表，就知道下一阶段的优化预算该拨给哪一项。

\topic{基准程序的思想}
CPI 依赖程序，“这台机器有多快”于是绕不开一个前置问题：用哪
个程序测？单一 kernel（计算内核）有两个经典陷阱。一是代表性：一个矩阵乘
测不出分支预测的好坏，一个整数排序压不出浮点单元；任何单一程
序都只覆盖机器能力的一个切片。二是可被定向优化：编译器厂商认
出 benchmark 的热循环，专门为它生成特殊代码，分数涨了，普通
程序没有变快——度量一旦变成目标，就不再是好的度量。

SPEC 类基准套件的思想就是正面回应这两个陷阱：不用单一 kernel，
而选一组真实的、有公开源码的程序，覆盖整数与浮点、计算密集与
存储密集，按用途分成若干组；每个程序固定输入、单独计时，报告
相对参照机的比值；最后把全部程序的比值聚合成总分。整套路数里，
“用哪个程序测”与“怎么加权”同样重要：程序选偏了，加权再公道
也没用；加权选错了，真实的分项结果也会被平均数扭曲。

聚合为什么用几何平均而不是算术平均，一个小例子就能看清。设机
器乙相对机器甲在两个程序上的加速比分别是 4 与 0.25（一个快四
倍，一个慢四倍）：

\begin{lstlisting}
算术平均 = (4 + 0.25) / 2 = 2.125      （看似乙大胜）
几何平均 = sqrt(4 x 0.25) = 1          （实际是平手）
\end{lstlisting}

算术平均被那个 4 主导；几何平均对“快多少”与“慢多少”一视同
仁，快四倍与慢四倍恰好抵消。几何平均还有两个实用性质：比值换
一台参照机重算，各机器的名次不变（算术平均会随参照机换人）；
连乘结构的加速比与连乘结构的优化过程相容，分步改进可以直接相
乘再开方。代价也要说清：几何平均描述的是“相对表现的典型水平”，
不预测任何一个程序的真实运行时间；规范的做法是永远同时给出每
个程序的分项，几何平均只当索引用。测试纪律也是套件的一部分：
固定输入数据、公开编译选项、禁止针对个别程序的特判——缺了这
三条，平均分再漂亮也只是广告。

基准套件自己也会过时，这是容易被忽略的一点。程序会随时代失去
代表性：十几年前的热点程序放到今天，可能整个被 cache 装下，存
储子系统毫无压力；编译器会逐代学会套件里的惯用写法，分数的水
分逐年增加；负载形态本身也在迁移，从单核整数程序到并行服务再
到 AI 推理。所以严肃的套件都定期换血、严格带版本号——引用一
个分数而不带版本与分项，等于引用一个没有时间地点的测量值。本
书的立场是：套件用来训练判断力，不是用来背分数。

套件的内部结构同样值得看一眼。常见做法是按整数与浮点分成两组，
因为两类负载几乎不共享瓶颈：整数程序考分支预测与指针追踪，浮
点程序考运算管线与存储带宽。有的套件还区分 speed 与 rate 两种
测法——前者测一个程序单独跑多快，是时延视角；后者测同时跑若
干份时的总完成量，是吞吐视角——一套基准把“时延、吞吐与利用率”一节分开的两个量
都照顾到了。读分时先弄清看到的是哪一组、哪一种测法：整数
speed 领先、浮点 rate 落后并不矛盾，它只是说明瓶颈长在不同程
序上。这也回答了“为什么不干脆只报一个总分”：总分是给采购的
一行摘要，分组与分项才是给工程师的地图。

\topic{案例：两个设计点的完整对比}
把本节全部工具用在一次设计裁决上。同一套 ISA 有两个实现：设计
A 走高频路线，2 GHz、平均 CPI 1.2，靠深流水把周期做短；设计 B
走宽发射路线，1.5 GHz、平均 CPI 0.9，每拍干得更多但频率上不去。
先跑同一个程序，动态指令数 $10^9$ 条：

\begin{longtable}{L{0.09\linewidth}L{0.12\linewidth}L{0.1\linewidth}L{0.11\linewidth}L{0.24\linewidth}L{0.19\linewidth}}
\toprule
设计 & 主频 & CPI & IPC & 每秒指令数 & 程序时间 \\
\midrule
A & 2 GHz & 1.2 & 0.83 & $2\times0.83\approx1.67$（G 条/秒） & $1.2/2=0.60$ s \\\tablerowrule
B & 1.5 GHz & 0.9 & 1.11 & $1.5\times1.11\approx1.67$（G 条/秒） & $0.9/1.5=0.60$ s \\
\bottomrule
\end{longtable}

精确打平。每秒指令数等于主频乘 IPC：A 用频率补 IPC，B 用 IPC
补频率，乘积相同，程序时间自然相同——“高频低 IPC”与“低频高
IPC”在这个程序上分不出胜负，这正是单看主频或单看 IPC 都会误
判的原因。

换一个程序立刻分出胜负。程序乙分支密集：每 4 条指令 1 个分支，
误预测率 10\%，A 的深流水每次误测冲刷 8 拍，B 只冲刷 2 拍：

\begin{lstlisting}
A: 有效 CPI = 1.2 + 0.25 x 0.10 x 8 = 1.40   时间 = 1.40/2   = 0.70 s
B: 有效 CPI = 0.9 + 0.25 x 0.10 x 2 = 0.95   时间 = 0.95/1.5 = 0.63 s
\end{lstlisting}

分支预测一进场，浅而宽的 B 赢了约 10\%。再换程序丙：无分支、
全命中的规则数值循环，两台机器的有效 CPI 都退化为 1.0，频率说
了算——A 用 0.50 s，B 用 0.67 s，A 反赢 33\%。三个程序三种结
果：打平、B 赢、A 赢。谁赢取决于程序，而程序通过 CPI 这一列进
入乘积。换成采购语言就是：挑机器之前先弄清自己的程序长什么样
——分支密集选预测好、冲刷便宜的，规则计算选频率高的；厂商的
综合分可以当初筛，不能当裁决，原因上一个主题已经说透。

还有一个读表细节：程序甲上 0.60 s 对 0.60 s 是这组参数下的精
确结果，不是设计哲学的必然——CPI 只差 0.05，结论就会翻转。两
个设计点离得越近，越要回到失效率的分项数据上做判断：谁的
miss 率更低、谁的误测冲刷更便宜，这些小数点后一位的差别，比
主频整数位的差别更能决定胜负。只比总分，等于在噪声里掷硬币。

这张对比表的使用前提也要写明：两台机器同一套 ISA、跑同一份二
进制，指令数这一列才相同，裁决才能只交给 CPI 与周期两列。一旦
连编译器或 ISA 也变了，指令数列必须进表重算——设 B 平台的编
译器把同一程序编出 $1.1\times10^9$ 条指令，程序甲的时间就变成
$1.1\times0.9/1.5=0.66$ s，原本的打平立刻反转成 A 赢。三因子
缺一个都算不完：读任何性能对比，先检查表里有没有指令数、CPI、
周期这三列，缺列的对比不是简化，是隐瞒。

本节收束为一条纪律：主频不是性能，IPC 不是性能，指令数也不是
性能，只有“指令数 × CPI × 周期”的乘积是性能。“Amdahl 定律与加速极限”一节回答它
的对偶问题：已经知道时间花在哪了，优化其中一部分，整体能快多
少——Amdahl 定律给出上限。

\section{四、Amdahl 定律与加速极限}

本节回答“优化值不值得做、先做哪一个”。直觉的回答是“把最慢的部件换快”，但整机的快慢取决于时间花在哪里，而不取决于哪个部件孤立地看最慢。Amdahl 定律把这件事写成一条公式：只需要两个输入——被优化部分原来占总时间的比例、这部分能快多少倍——就能算出整机快多少。它是性能分析里被引用最多、也被误用最多的一条公式：只报局部加速比、不报占比，是最常见的半截话；只记得“加速有上限”、忘了上限怎么算，是第二常见的半截话。

\topic{Amdahl 公式的推导}
设机器跑某程序的总时间为 $T$。用测量手段（“性能测量与案例”一节的性能计数器，或存储与设备事务相关章节的总线监视）把 $T$ 分成两份：打算优化的部分原来占 $p$（$0\le p\le1$），其余部分占 $1-p$。现在让前一部分快 $s$ 倍（$s>1$），后一部分原封不动。推导只有三步：

\begin{lstlisting}
优化前:  T  = p*T + (1-p)*T          (两份时间相加)
优化后:  T' = p*T/s + (1-p)*T        (优化部分快 s 倍，其余不变)
加速比:  S  = T / T' = 1 / ((1-p) + p/s)
\end{lstlisting}

第三步把 $T$ 约掉，剩下的式子里只有 $p$ 和 $s$ 两个数，这就是 Amdahl 定律：$S=1/((1-p)+p/s)$。它没有引入任何近似，唯一的假设是“优化不影响其余部分”——被优化的部分单独计时，没碰的部分原样计时。这个假设后面还要反复回来检查。

两个输入缺一不可。只说“这个优化让内存访问快 2 倍”而不说它占多少时间，等于什么都没说：内存访问占 60\% 时这是整机 1.43 倍，占 5\% 时只有约 1.03 倍，同一个“快 2 倍”差出十几倍的整体效果。反过来只说“这部分占 80\%”同样不够，还得知道能快几倍。占比和局部加速比是同一个乘积里的两个因子，缺了任何一个，另一个再大也定不了积。

\begin{lstlisting}
同一个"局部快 2 倍"，占比不同，整体完全不同:
  p = 60%:  S = 1/(0.4  + 0.6/2)  = 1/0.700 = 1.43
  p = 25%:  S = 1/(0.75 + 0.25/2) = 1/0.875 = 1.14
  p = 5%:   S = 1/(0.95 + 0.05/2) = 1/0.975 = 1.026
\end{lstlisting}

再给一个完整演算。某部分占 40\%，优化后快 2 倍：$T'=0.4T/2+0.6T=0.8T$，$S=1/0.8=1.25$。局部快 2 倍只换来整机快 1.25 倍，这个“打折”不是工程损耗，是算术必然：优化后那 40\% 缩成了 20\%，省下的 20\% 就是全部收益，而其余 80\% 的时间一分不少。比例看不懂时，把 $T$ 换成具体毫秒再算一遍就清楚了：

\begin{lstlisting}
设 T = 200 ms，被优化部分占 40% 即 80 ms，其余 120 ms
优化后: 80/2 + 120 = 160 ms
加速比: 200/160 = 1.25        (与比例算法一致)
\end{lstlisting}

处理器核心相关章节按拍数逐条核算指令时间时，读者已经见过同样的思想：总时间永远被所有部分分完，加快任何一部分，收益都以它占的份额为限。本节只是把“拍数”换成“时间份额”，把同一思想写成通用公式。也正因如此，$p$ 和 $s$ 的来源必须分别说清楚：

\begin{longtable}{L{0.16\linewidth}L{0.36\linewidth}L{0.34\linewidth}}
\toprule
输入 & 含义 & 从哪里来 \\
\midrule
占比 $p$ & 被优化部分原来占总时间的份额 & 只能测量：profiler、性能计数器、总线监视 \\\tablerowrule
局部加速 $s$ & 该部分优化前后的时间比 & 微基准实测，或按部件参数推算 \\
\midrule
整体加速 $S$ & \multicolumn{2}{l}{由前两者唯一决定：$S=1/((1-p)+p/s)$，不是第三个独立输入} \\
\bottomrule
\end{longtable}

把推导浓缩成一句使用说明：先测 $p$，再估 $s$，最后才谈整体。实践中这条公式的误用也集中在这两个输入上，对照如下：

\begin{longtable}{L{0.30\linewidth}L{0.30\linewidth}L{0.28\linewidth}}
\toprule
常见说法 & 问题出在哪 & 正确问法 \\
\midrule
“新部件快 10 倍，值得上” & 只有 $s$，没有 $p$ & 这部分原来占总时间的多少 \\\tablerowrule
“这段代码占了 80\% 的时间” & 只有 $p$，没有 $s$ & 它能被加快多少倍 \\\tablerowrule
“整机应该能快 3 倍” & 把 $s$ 当成了 $S$ & 代入 $1/((1-p)+p/s)$ 再报数 \\\tablerowrule
“优化后这部分只占 5\% 了” & 拿优化后的占比当 $p$ & $p$ 必须取优化前的份额 \\
\bottomrule
\end{longtable}

末行是最隐蔽的一处误用：公式里的 $p$ 是“优化前”的份额。优化见效之后这部分已经缩水，拿缩水后的占比再代一次公式，等于把同一笔收益算了第二遍。所有迭代式优化（做一轮、测一轮、再做一轮）都必须回到新一轮测量重新取 $p$，不能沿用上一次的数字。

\topic{加速上限表}
把 $p$ 取六档、$s$ 取四档代入公式，逐个算出来，就是下面这张表。它是性能工程里最值得背下来的一张表：背下来之后，任何“某部件快了几倍”的宣传都能在十秒内翻译成整机收益。

\begin{longtable}{L{0.22\linewidth}L{0.16\linewidth}L{0.16\linewidth}L{0.16\linewidth}L{0.18\linewidth}}
\toprule
占比 $p$ & 局部快 2 倍 & 局部快 4 倍 & 局部快 8 倍 & 无限加速上限 \\
\midrule
10\% & 1.05 & 1.08 & 1.10 & 1.11 \\\tablerowrule
25\% & 1.14 & 1.23 & 1.28 & 1.33 \\\tablerowrule
50\% & 1.33 & 1.60 & 1.78 & 2.00 \\\tablerowrule
75\% & 1.60 & 2.29 & 2.91 & 4.00 \\\tablerowrule
90\% & 1.82 & 3.08 & 4.71 & 10.00 \\\tablerowrule
95\% & 1.90 & 3.48 & 5.93 & 20.00 \\
\bottomrule
\end{longtable}

把这张表的四列画成曲线，“占比决定上限”就从数字变成了形状：

\begin{pdfdiagram}
  % 横轴 p：0-100%；纵轴整机加速比 S，y = (S-1)*0.22
  \draw[BookMuted, line width=0.6pt, ->] (0,0) -- (10.2,0);
  \draw[BookMuted, line width=0.6pt, ->] (0,0) -- (0,4.75);
  \node[hw label, anchor=north] at (4.8,-0.62) {优化前占比 $p$};
  \node[hw label, rotate=90, anchor=south] at (-0.95,2.3) {整机加速比 $S$};
  \foreach \x/\t in {0/0, 2.4/25\%, 4.8/50\%, 7.2/75\%, 9.6/100\%} {
    \draw[BookMuted, line width=0.5pt] (\x,0.06) -- (\x,-0.06);
    \node[hw label, anchor=north] at (\x,-0.1) {\t};
  }
  \foreach \y/\t in {0/1, 0.88/5, 1.98/10, 3.08/15, 4.18/20} {
    \draw[BookMuted, line width=0.5pt] (0.06,\y) -- (-0.06,\y);
    \node[hw label, anchor=east] at (-0.12,\y) {\t};
  }
  % S = 1/((1-p)+p/s)，四条曲线 s = 2/4/8/无穷
  \draw[BookMuted, line width=0.9pt] plot[smooth] coordinates {
    (0,0) (0.96,0.012) (2.4,0.031) (4.8,0.073) (7.2,0.132) (8.64,0.180) (9.12,0.199)
  };
  \draw[BookTeal, line width=0.9pt] plot[smooth] coordinates {
    (0,0) (0.96,0.018) (2.4,0.051) (4.8,0.132) (7.2,0.283) (8.64,0.457) (9.12,0.545)
  };
  \draw[BookAccent, line width=0.9pt] plot[smooth] coordinates {
    (0,0) (0.96,0.022) (2.4,0.062) (4.8,0.171) (7.2,0.420) (8.64,0.815) (9.12,1.084)
  };
  \draw[BookRed, line width=1pt] plot[smooth] coordinates {
    (0,0) (0.96,0.024) (2.4,0.073) (4.8,0.22) (7.2,0.66) (8.64,1.98) (9.12,4.18)
  };
  \node[hw label, anchor=west] at (9.22,0.199) {$s{=}2$};
  \node[hw label, anchor=west] at (9.22,0.545) {$s{=}4$};
  \node[hw label, anchor=west] at (9.22,1.084) {$s{=}8$};
  \node[hw label, anchor=west] at (9.22,4.18) {$s\to\infty$};
  \node[hw label, anchor=west] at (5.5,1.25) {上限 $1/(1-p)$};
\end{pdfdiagram}
\diagramnote{Amdahl 定律的四条加速曲线。$s\to\infty$ 那条就是上限 $1/(1-p)$，$p=0.9$ 时也只到 10；其余三条被它压在下面。读法：$s$ 决定曲线弯成什么形状，$p$ 决定曲线能爬到多高——占比才是主变量。}

最后一列是 $s\to\infty$ 的极限，公式退化成一个更简单的式子：

\begin{lstlisting}
s -> 无穷 时 p/s -> 0，分母只剩 (1-p):
  p = 50%:  上限 = 1/0.5  = 2.00
  p = 90%:  上限 = 1/0.1  = 10.00
  p = 95%:  上限 = 1/0.05 = 20.00
\end{lstlisting}

“无限加速”的含义是把这部分优化到不花时间——瞬时完成。即使做到这一步，整机也只快 $1/(1-p)$ 倍：占一半的部分优化到无穷快，整机快 2 倍；占九成的部分优化到无穷快，整机也才快 10 倍。上限被未优化的 $1-p$ 限定，与被优化部分究竟多快毫无关系。这就是“Amdahl 上限”的准确含义：它不是工程上暂时够不着的理想，而是数学上的天花板。

读这张表要竖着读，也要横着读。竖着看：同一列里占比每涨一档，收益跳一档——占比是收益的主变量，想拿大加速先找大占比。横着看：同一行里 $s$ 从 2 提到 4 收益明显，从 8 提到无穷反而所剩无几。把 $p=75\%$ 那一行单独展开，收益递减看得更清楚：

\begin{longtable}{L{0.18\linewidth}L{0.18\linewidth}L{0.24\linewidth}L{0.26\linewidth}}
\toprule
局部加速 $s$ & 整体加速 & 比上一档多拿 & 关键问题 \\
\midrule
2 & 1.60 & — & 第一笔投入最划算 \\\tablerowrule
4 & 2.29 & 0.69 & 仍然值得 \\\tablerowrule
8 & 2.91 & 0.62 & 开始递减 \\\tablerowrule
16 & 3.37 & 0.46 & 投入翻倍，增量缩水 \\\tablerowrule
32 & 3.66 & 0.29 & 接近天花板 \\\tablerowrule
$\infty$ & 4.00 & 0.34 & 理论极限，与 32 倍只差 0.34 \\
\bottomrule
\end{longtable}

工程含义很直接：局部加速比做到一个量级之后，就该回头检查占比，而不是继续提高同一个部件的 $s$。把 $s$ 从 8 做到 16 往往要花一代工艺，把同样的精力转去扩大优化的覆盖面——让更多代码路径受益于这次优化——常常能多拿一整档。教学计算机项目中的教学机是这张表最直白的注脚：它每条访存都要等一个完整总线周期，访存在总时间里占大头，任何“把 ALU 做快”的改法都受表中前几行所示占比限制；反过来，存储与设备事务相关章节给主存加一级 cache、让大多数访存不再等总线，才是对它真正有效的那一行。

上限表还有一个用法：快速否决。发布会和宣传页上的数字，都可以代入公式反推，看占比是否对得上：

\begin{lstlisting}
宣传: "新浮点单元快 10 倍，整机快 3 倍"
反推占比: 3 = 1/(1 - p + p/10)  ->  1/3 = 1 - 0.9p  ->  p = 0.74
含义: 浮点要占到总时间的 74%，这个宣传才成立
若 profile 实测浮点只有 40%:  S = 1/(0.6 + 0.4/10) = 1/0.64 = 1.56
结论: 3 倍不成立，真实收益 1.56 倍
\end{lstlisting}

反推出的 74\% 与实测的 40\% 差了近一倍，宣传里的“整机快 3 倍”只能在某个浮点极密集的特例程序上成立。这类反推不需要任何仪器，一张上限表加一条公式就够，是读评测报告时最省钱的过滤器。

上限表还可以反过来用：已知整机加速和局部加速，反推占比是否说得通。某团队报告“cache 改造让访存快 4 倍，整机快了 1.5 倍”，反推一下：

\begin{lstlisting}
1.5 = 1/(1 - p + p/4)  ->  2/3 = 1 - 0.75p  ->  p = 0.444
含义: 访存原来要占 44.4%，这组数字才互相说通
核对: 该程序的访存占比实测约 70%
若 p = 0.70:  S = 1/(0.3 + 0.7/4) = 1/0.475 = 2.11，不是 1.5
\end{lstlisting}

对不上本身就是信息：要么“快 4 倍”只是微基准里的理想值，真实程序没有兑现；要么改造影响了其他部分（例如 cache 结构变了之后取指行为也跟着变了），违反了“其余部分不动”的假设。无论哪一种，下一步都是回到测量，而不是在公式里调数字。

\topic{多核扩展的 Amdahl 极限}
多核是“局部加速”最引人注目的特例：把程序的可并行部分摊到 $n$ 个核上，相当于这部分 $s=n$；而串行部分——启动、收尾、归约、临界区、I/O——一个核也帮不上，原样留在 $1-p$ 里。设串行部分占 5\%，可并行部分占 95\%，代入公式：

\begin{lstlisting}
S(n) = 1 / (0.05 + 0.95/n)
n = 64:  S = 1/(0.05 + 0.95/64) = 1/0.0648 = 15.42
n -> 无穷:  S = 1/0.05 = 20
\end{lstlisting}

\begin{longtable}{L{0.14\linewidth}L{0.22\linewidth}L{0.20\linewidth}L{0.18\linewidth}}
\toprule
核数 $n$ & 并行部分耗时 & 总耗时 & 整体加速 \\
\midrule
2 & 0.4750 & 0.5250 & 1.90 \\\tablerowrule
4 & 0.2375 & 0.2875 & 3.48 \\\tablerowrule
8 & 0.1188 & 0.1688 & 5.93 \\\tablerowrule
16 & 0.0594 & 0.1094 & 9.14 \\\tablerowrule
32 & 0.0297 & 0.0797 & 12.55 \\\tablerowrule
64 & 0.0148 & 0.0648 & 15.42 \\\tablerowrule
128 & 0.0074 & 0.0574 & 17.41 \\\tablerowrule
256 & 0.0037 & 0.0537 & 18.62 \\\tablerowrule
$\infty$ & 0 & 0.0500 & 20.00 \\
\bottomrule
\end{longtable}

64 个核只换来约 15.4 倍加速——不是直觉里的 64 倍，离上限 20 倍也还差一截。这张表和上一张上限表是同一族曲线：95\% 占比那一行，局部加速取 2、4、8 时恰好就是 2 核、4 核、8 核的整机加速，核数只是 $s$ 的另一种写法。

把这张表画成曲线，串行部分像地板一样接住整条曲线的形状就清楚了：

\begin{pdfdiagram}
  % 横轴核数（log2，每档 1.1）；纵轴加速比，y = S * 0.21
  \draw[BookMuted, line width=0.6pt, ->] (0,0) -- (10.0,0);
  \draw[BookMuted, line width=0.6pt, ->] (0,0) -- (0,4.8);
  \node[hw label, anchor=north] at (4.6,-0.62) {核数 $n$（对数刻度）};
  \node[hw label, rotate=90, anchor=south] at (-0.95,2.3) {整机加速比 $S$};
  \foreach \x/\t in {0/1, 2.2/4, 4.4/16, 6.6/64, 8.8/256} {
    \draw[BookMuted, line width=0.5pt] (\x,0.06) -- (\x,-0.06);
    \node[hw label, anchor=north] at (\x,-0.1) {\t};
  }
  \foreach \y/\t in {0.21/1, 1.05/5, 2.1/10, 3.15/15, 4.2/20} {
    \draw[BookMuted, line width=0.5pt] (0.06,\y) -- (-0.06,\y);
    \node[hw label, anchor=east] at (-0.12,\y) {\t};
  }
  % S(n) = 1/(0.05 + 0.95/n)
  \draw[BookAccent, line width=1pt] plot[smooth] coordinates {
    (0,0.21) (1.1,0.40) (2.2,0.73) (3.3,1.245) (4.4,1.92) (5.5,2.64) (6.6,3.24) (7.7,3.66) (8.8,3.91)
  };
  \node[hw label, anchor=west] at (0.5,2.3) {$S(n)=1/(0.05+0.95/n)$};
  % 渐近上限 20
  \draw[BookRed, dashed, line width=0.6pt] (0,4.2) -- (9.6,4.2);
  \node[hw label, anchor=east] at (9.6,4.4) {渐近上限 20（串行 5\%）};
  % 64 核标记
  \draw[BookTeal, dashed, line width=0.5pt] (6.6,0) -- (6.6,3.24) -- (0,3.24);
  \fill[BookRed] (6.6,3.24) circle (1.6pt);
  \node[hw label, anchor=west] at (6.75,3.0) {64 核：15.42};
\end{pdfdiagram}
\diagramnote{串行占比 5\% 时整机加速比随核数的曲线（横轴对数刻度）：8 核 5.93、16 核 9.14、64 核 15.42，之后核数每翻一番，增量越缩越小。红色虚线是 $n\to\infty$ 的渐近上限 20——串行部分像地板一样等在那里，核数再多也只是无限贴近，永远穿不过去。}

核数翻倍不等于速度翻倍，而且差距越拉越大：16 核加到 32 核，整机快 $12.55/9.14\approx1.37$ 倍；64 核加到 128 核，只快 $17.41/15.42\approx1.13$ 倍；128 核加到 256 核，只剩 $18.62/17.41\approx1.07$ 倍。曲线弯下来的原因在分母：核数越多，总耗时越接近串行部分的 0.05，新增的核全部砸在已经只剩零头的并行部分上。串行部分像地板一样等在那里，核数再多也只是无限贴近它，永远穿不过去。换算成采购语言更刺眼：

\begin{longtable}{L{0.14\linewidth}L{0.18\linewidth}L{0.24\linewidth}L{0.28\linewidth}}
\toprule
核数 & 整体加速 & 每核效率（加速/核数） & 关键问题 \\
\midrule
16 & 9.14 & 0.57 & 效率尚可 \\\tablerowrule
32 & 12.55 & 0.39 & 效率明显下滑 \\\tablerowrule
64 & 15.42 & 0.24 & 四分之三的核在陪跑 \\
\bottomrule
\end{longtable}

所以多核工程的真实收益不来自“加核”，而来自把串行部分也改小：全局锁拆成分段锁，归约改成树形，启动和收尾改成与计算流水交叠，串行 I/O 改成批量异步。效果可以直接算出来——同样 64 核，串行占比从 5\% 改到 1\%，加速从 15.42 跳到 39.26：

\begin{longtable}{L{0.16\linewidth}L{0.20\linewidth}L{0.20\linewidth}L{0.28\linewidth}}
\toprule
串行占比 & 64 核加速 & 无限核上限 & 关键问题 \\
\midrule
10\% & 8.77 & 10 & 核多无用武之地 \\\tablerowrule
5\% & 15.42 & 20 & 服务器常见量级 \\\tablerowrule
2\% & 28.32 & 50 & 需要专门改串行段 \\\tablerowrule
1\% & 39.26 & 100 & 动串行一档，胜过加核四倍 \\
\bottomrule
\end{longtable}

末行那句话值得再算一遍：串行 5\% 时从 64 核加到 256 核，核数翻两番，加速只从 15.42 涨到 18.62；串行从 5\% 改到 1\%，加速从 15.42 涨到 39.26。并行的真实收益来自把串行部分也变小，这不是口号，是第二列对第一列的斜率。还要提醒一句：5\% 的串行占比已经算乐观——真实并行程序还有负载不均、通信和同步开销，它们的效果等价于把串行占比再抬高几档，表里的数字只会更差。

“把串行部分改小”也不是一句口号，它对应一组具体手段，每种手段削的都是串行占比的不同来源：

\begin{longtable}{L{0.20\linewidth}L{0.34\linewidth}L{0.30\linewidth}}
\toprule
串行来源 & 典型场景 & 缩小手段 \\
\midrule
启动与收尾 & 主线程分发任务、回收结果 & 分发与计算交叠，收尾流水化 \\\tablerowrule
全局锁 & 所有核排队进同一个临界区 & 分段锁、无锁结构、线程私有副本 \\\tablerowrule
归约 & 各核结果最后串行汇总 & 树形归约，深度 $\log n$ 代替 $n$ \\\tablerowrule
串行 I/O & 计算前一次读入、算完一次写出 & 预取、双缓冲、批量异步 \\\tablerowrule
负载不均 & 最快的核等最慢的核 & 动态调度、任务切细 \\
\bottomrule
\end{longtable}

负载不均值得单独说一句：它不在公式里出现，效果却等价于串行占比升高——64 个核里 63 个做完了等最后 1 个，等待期间等价于 63 个核被串行化。所以表里的“任务切细”不是锦上添花：任务粒度越细，调度器越容易把各核的完工时间拉齐，等价的串行尾巴越短。

串行占比本身也不是猜的，可以测：同一程序分别在 1、2、4、8 核上测总时间，算出各档加速 $S(n)$，理论上 $1/S(n)=f+(1-f)/n$（$f$ 为串行份额），把 $1/S$ 对 $1/n$ 描点应近似一条直线——截距是 $f$，斜率是 $1-f$。若点明显跑到直线上方，说明存在随核数增长的额外开销（通信、竞争），串行份额本身在随核数上升。这是公式之外、实测之内才能看到的东西，也是“带宽、roofline 与瓶颈定位”一节的测量方法要回答的问题。

\begin{lstlisting}
串行份额的实测拟合:
  1/S(n) = f + (1-f)/n
  以 1/S 对 1/n 描点: 截距 = f，斜率 = 1-f
  点上弯(高于直线) = 有随核数增长的额外开销
\end{lstlisting}

教学计算机项目中的教学机是这条曲线的另一个端点：它整机只有一个执行流，串行占比 100\%，加核对它毫无意义——它能讲清楚并行，恰恰是因为自己一点也不并行。

\topic{Gustafson 的对照视角}
Amdahl 的假设是“问题规模固定”：同一道题，核越多越好。Gustafson 换了一个假设：机器变大了，题也跟着变大——加核不是为了把同一道题算得更快，而是为了在同样的时间里算更细的网格、更大的模型、更多的粒子。这时串行部分占墙上时间的比例 $f$ 不随规模增长（启动还是那几秒），并行部分却随核数线性放大。推导同样只有几步：

\begin{lstlisting}
n 核机器上，墙上时间归一化为 1: 串行部分占 f，并行部分占 1-f
同一问题搬到单核: 串行部分仍是 f
                并行部分一个核要做 n 份，为 n*(1-f)
单核总时间 = f + n*(1-f) = n - (n-1)*f    <- 放大规模下的加速比
代入 f = 0.05, n = 64:  64 - 63*0.05 = 60.85
\end{lstlisting}

同样 5\% 的串行占比、同样 64 个核，Amdahl 算出 15.42，Gustafson 算出 60.85，约为 Amdahl 的四倍。逐档对照：

\begin{longtable}{L{0.14\linewidth}L{0.24\linewidth}L{0.26\linewidth}L{0.22\linewidth}}
\toprule
核数 $n$ & Amdahl（规模固定） & Gustafson（规模放大） & 关键问题 \\
\midrule
2 & 1.90 & 1.95 & 核少时两者接近 \\\tablerowrule
8 & 5.93 & 7.65 & 分歧开始显现 \\\tablerowrule
16 & 9.14 & 15.25 & 差约三分之二 \\\tablerowrule
32 & 12.55 & 30.45 & 差出一倍以上 \\\tablerowrule
64 & 15.42 & 60.85 & 假设差异完全显形 \\\tablerowrule
$\infty$ & 20.00 & $\infty$ & 一个有顶，一个没顶 \\
\bottomrule
\end{longtable}

把两列画到同一张图上，分歧就成了两条曲线的形状：一条被串行地板接住，一条近似直线向上。

\begin{pdfdiagram}
  % 横轴核数 0..64（线性，每核 0.14）；纵轴加速比，y = S * 0.068
  \draw[BookMuted, line width=0.6pt, ->] (0,0) -- (10.0,0);
  \draw[BookMuted, line width=0.6pt, ->] (0,0) -- (0,4.9);
  \node[hw label, anchor=north] at (4.8,-0.62) {核数 $n$};
  \node[hw label, rotate=90, anchor=south] at (-0.95,2.4) {整机加速比};
  \foreach \x/\t in {1.12/8, 2.24/16, 4.48/32, 8.96/64} {
    \draw[BookMuted, line width=0.5pt] (\x,0.06) -- (\x,-0.06);
    \node[hw label, anchor=north] at (\x,-0.1) {\t};
  }
  \foreach \y/\t in {1.36/20, 2.72/40, 4.08/60} {
    \draw[BookMuted, line width=0.5pt] (0.06,\y) -- (-0.06,\y);
    \node[hw label, anchor=east] at (-0.12,\y) {\t};
  }
  % Gustafson：S = n - (n-1)*0.05，近似直线，64 核 60.85
  \draw[BookTeal, line width=1pt] (0.14,0.068) -- (8.96,4.138);
  \fill[BookTeal] (8.96,4.138) circle (1.6pt);
  % Amdahl：S = 1/(0.05+0.95/n)，64 核 15.42，渐近 20
  \draw[BookAccent, line width=1pt] plot[smooth] coordinates {
    (0.14,0.068) (0.28,0.129) (1.12,0.403) (2.24,0.62) (4.48,0.853) (8.96,1.049)
  };
  \fill[BookRed] (8.96,1.049) circle (1.6pt);
  \draw[BookRed, dashed, line width=0.6pt] (0,1.36) -- (9.6,1.36);
  \node[hw label, anchor=west] at (0.15,1.52) {Amdahl 上限 20};
  % 图例
  \draw[BookTeal, line width=1pt] (0.4,4.55) -- (1.1,4.55);
  \node[hw label, anchor=west] at (1.25,4.55) {Gustafson：规模放大，64 核 60.85};
  \draw[BookAccent, line width=1pt] (0.4,4.15) -- (1.1,4.15);
  \node[hw label, anchor=west] at (1.25,4.15) {Amdahl：规模固定，64 核 15.42};
\end{pdfdiagram}
\diagramnote{同样 5\% 串行占比的两条加速曲线：Amdahl 固定问题规模，64 核 15.42，渐近上限 20，核数越多越撞上串行地板；Gustafson 固定墙上时间、让规模随核数放大，64 核 60.85，没有顶。两个公式都没有错，分歧只在“什么保持不变”——题固定，还是墙上时间固定。}

两个公式都没有错，分歧只在“什么保持不变”。Amdahl 固定问题规模，是悲观的：规模不变时串行占比原地不动，核越多越撞上它。Gustafson 固定墙上时间、放大问题规模，是乐观的：规模放大后并行部分变大，串行占比被稀释——原本占 5\% 的启动开销，在放大一千倍的问题里可能只占 0.005\%。历史背景也有意思：Amdahl 定律 1967 年提出时的语境是为单处理器辩护，论证并行机前途有限；Gustafson 在 1988 年用上千个处理器的机器按放大规模实测，拿到接近线性的加速，说明那场争论的双方各抓住了问题的一半。

适用场合按“用户等的是什么”来分：交互程序、实时控制、固定的夜间批处理，题就这么大，用户等的是同一道题快多少，用 Amdahl 估计才诚实；超算中心、气象预报、AI 训练这类“机器多大题就多大”的负载，预算批的是算力规模，用 Gustafson 更贴近真实采购逻辑。反过来也有题不会变大的场合：编译一次工程、应答一次数据库查询，没人因为机器大了就把查询写复杂十倍，这些负载永远服从 Amdahl。成熟的做法是两个都算：Amdahl 给出固定规模的下限，Gustafson 给出放大规模的上限，真实收益位于两者之间，由“题到底会不会跟着变大”决定靠近哪一端——买机器之前先回答这个问题，比争论哪个公式“对”有用得多。

两个公式还可以用同一个具体场景各走一遍，看它们到底在算什么。同样 5\% 的串行占比、同样 64 个核、同样 60 分钟，两个场景的 60 分钟不是一回事：

\begin{lstlisting}
场景一(Amdahl，题固定): 单核跑完要 60 分钟，其中串行 3 分钟
  搬上 64 核: 3 + 57/64 = 3.89 分钟
  加速 = 60/3.89 = 15.42
场景二(Gustafson，题随机器放大): 64 核跑 60 分钟，串行段占 3 分钟
  这 60 分钟里并行段干了 57*64 = 3648 核·分钟的活
  同一道题搬回单核: 3 + 3648 = 3651 分钟
  加速 = 3651/60 = 60.85
\end{lstlisting}

场景一的 57 分钟是总工作量，场景二的 57 分钟是墙上时间（背后是 64 份并行工作量）。把两个场景的数字混着用——拿场景二的占比代场景一的公式——是围绕 Gustafson 最常见的计算错误；分清“题固定”还是“时间固定”，是两个公式共同的前提检查。

\topic{案例：三个候选优化点的排序}
把本节全部工具用在一个完整问题上。某程序 profile 一轮，时间分布为：内存访问 60\%，整数计算 30\%，分支代价 10\%。手头有三个候选方案，每个只动一处：方案 A 做 cache 优化，让内存访问快 2 倍；方案 B 重做 ALU，让整数计算快 3 倍；方案 C 加分支预测，消除 80\% 的分支代价。哪个先做？不许凭感觉，逐个代入公式：

\begin{lstlisting}
方案 A: 内存占 60%，快 2 倍
  T' = 0.6/2 + 0.4 = 0.70    S = 1/0.70 = 1.43
方案 B: 整数占 30%，快 3 倍
  T' = 0.3/3 + 0.7 = 0.80    S = 1/0.80 = 1.25
方案 C: 分支占 10%，消除 80% 代价 = 这部分只剩 20%，即快 5 倍
  T' = 0.1/5 + 0.9 = 0.92    S = 1/0.92 = 1.09
\end{lstlisting}

\begin{longtable}{L{0.26\linewidth}L{0.20\linewidth}L{0.12\linewidth}L{0.14\linewidth}L{0.10\linewidth}}
\toprule
方案 & 作用部分（占比） & 局部加速 & 整体加速 & 排序 \\
\midrule
A：cache 优化 & 内存（60\%） & 2 & 1.43 & 1 \\\tablerowrule
B：重做 ALU & 整数（30\%） & 3 & 1.25 & 2 \\\tablerowrule
C：分支预测 & 分支（10\%） & 5 & 1.09 & 3 \\
\bottomrule
\end{longtable}

排序结果 A$>$B$>$C，恰好与局部加速比的直觉排序相反：B 的局部 3 倍比 A 的 2 倍漂亮，C 的“消除 80\%”听起来最猛，但收益是占比与局部效果的乘积，两个因子谁小谁卡谁——60\% 配上平平的 2 倍，照样赢过 10\% 配上惊人的 5 倍。方案 C 还藏着本案例唯一的翻译陷阱：“消除 80\% 代价”不是 $p=0.8$，而是分支那 10\% 里消掉 80\%，即这部分时间缩到原来的五分之一；把自然语言翻译成 $p$ 和 $s$ 时用错档位，是 Amdahl 计算里最常见的错误。

这张表的输入全部来自 profile：60、30、10 是测出来的，不是看代码猜的。凭直觉很多人会先做 B，因为“ALU 快 3 倍”听起来最像硬件进步；数字说先做 A。这也是章首“先测量，再优化”的具体含义：没有占比这个输入，公式只剩一半，而一半公式给出的排序往往是错的。分支代价在处理器核心相关章节的流水线里见过实体——误预测的冲刷拍；存储与设备事务相关章节的 cache 则是方案 A 的实体。占比测量告诉工程师该把晶体管花在哪一章的技术上。排序对测量误差有多敏感，也可以当场检查：

\begin{longtable}{L{0.26\linewidth}L{0.20\linewidth}L{0.20\linewidth}L{0.18\linewidth}}
\toprule
内存占比实测 & 方案 A 整体加速 & 方案 B（不变） & 预期结果 \\
\midrule
50\%（低估 10 点） & 1.33 & 1.25 & A 仍第一 \\\tablerowrule
60\%（基准） & 1.43 & 1.25 & A 第一 \\\tablerowrule
70\%（高估 10 点） & 1.54 & 1.25 & A 第一 \\
\bottomrule
\end{longtable}

占比误差正负十个点之内排序不变，这个结论是稳健的；若三个方案的整体加速本来就只差一两个点，才需要回去提高测量精度。最后检查公式的边界：三个方案互斥的假设是“动一处不影响其余两处”。若 cache 优化同时改变了分支行为（例如预取挤掉了代码行），就要重新测量、重新代入，不能把三张单子各算一遍再相加——Amdahl 是线性拆分，交叉影响不在公式里。还要记住上限表的告诫：即使方案 A 做到内存访问无穷快，整机也只有 $1/0.4=2.5$ 倍；若业务需要 4 倍，单靠优化内存永远不够，必须同时动 B 和 C 让 $1-p$ 缩小。排序做完不是终点：做完 A 之后内存时间减半，新的时间分布约成 43/43/14，下一轮要重新 profile、重新排序，优化的顺序可能因此改变。把后两轮也算完，看这个迭代实际怎么走：

\begin{lstlisting}
第二轮(做完 A 之后，总时间已从 100 缩到 70):
  新分布: 内存 30、整数 30、分支 10
  方案 B: 30/3 + 40 = 50    本轮加速 = 70/50 = 1.40
  方案 C: 10/5 + 60 = 62    本轮加速 = 70/62 = 1.13
  仍选 B；做完 B 后分布: 内存 30、整数 10、分支 10(共 50)
第三轮只剩 C 可做:
  方案 C: 10/5 + 40 = 42    本轮加速 = 50/42 = 1.19
累计: 100 -> 70 -> 50 -> 42，整机共快 100/42 = 2.38 倍
\end{lstlisting}

三个数字各有一处值得停一下。第一，方案 B 的本轮加速从第一轮的 1.25 涨到 1.40：A 把总时间做小之后，整数计算的相对占比变大了——同一个方案的价值随工程进程变化，这正是每轮必须重测占比的原因。第二，累计 2.38 倍正好等于三轮加速的连乘 $1.43\times1.40\times1.19$，因为三个方案作用在互不重叠的部分上；若方案之间有交叉影响，连乘关系就不成立，必须回到合并后的分布重算。第三，2.38 倍离“内存无限快”的 2.5 倍上限已经很近：这个程序再往下做，每一轮的可选方案会越来越少、单轮收益越来越小，优化的尽头是被各部分 $1-p$ 轮流接力的天花板。

\section{五、带宽、roofline 与瓶颈定位}

“Amdahl 定律与加速极限”一节把“时间”拆成份额，这一节把“字节”摆上秤。现代机器的算力和带宽是两个独立的物理量：算力由片上晶体管和时钟决定，带宽由引脚、走线和 DRAM 工艺决定，两者的增长速度从来不一样。性能分析必须回答的第二个问题因此是：这个程序的瓶颈在算力侧，还是在带宽侧？方向判断错了，优化就是把钱花在错误的屋顶上。

\topic{Little 定律的直觉版}
排队论里的 Little 定律说：系统内平均在途数等于平均吞吐率乘以平均时延。它不假设任何排队规则，稳态下恒成立。用到内存系统上，形式特别直白：要维持某个带宽，必须随时有足够多的字节“在路上”——在途字节数等于带宽乘以时延。直觉画面是水管：水速一定时，管子越长，管中存的水越多；想维持出水量不变，加长水管就必须同时多灌水。

\begin{lstlisting}
要喂满带宽，需要维持在途的字节数 = 带宽 * 时延
  = 100 GB/s * 100 ns
  = 10^11 B/s * 10^-7 s
  = 10^4 B = 10 KB
按 64 B 一条 cache 行折算: 10000 / 64 = 156 个在途请求
结论: 这台机器必须随时有约 156 行访存"在路上"，
      在途数低于此值，带宽就吃不满
\end{lstlisting}

这个数字立刻解释了一个常见困惑：为什么单个核跑不出整机带宽。一个核同时在途的 miss 是有限的——主流乱序核的 miss 处理资源（MSHR）大约在十几条的量级，取 16 条、每条 64 B，在途预算就是 1 KB；$1\text{ KB}/100\text{ ns}\approx10.2$ GB/s，这就是单核带宽的封顶。整机 100 GB/s 需要约 10 个这样的核同时发出访存，或者靠预取把单核的在途数撑大。

把这段算术画成水管，闸门宽度与管中存水的关系就一眼可见：

\begin{pdfdiagram}
  \node[hw compute, minimum width=1.5cm, minimum height=1.0cm] (core) at (0,0) {CPU 核};
  \node[hw state, minimum width=2.0cm, minimum height=1.0cm] (mshr) at (2.55,0) {MSHR 闸门\\16 条 $\times$ 64 B};
  \node[hw memory, minimum width=1.7cm, minimum height=1.0cm] (dram) at (9.7,0) {DRAM\\时延 100 ns};
  % 水管：在途字节
  \draw[BookTeal, line width=0.9pt, fill=BookCodeBg] (3.75,-0.42) rectangle (8.6,0.42);
  \node[hw label] at (6.17,0) {在途字节 = 带宽 $\times$ 时延};
  \draw[hw bus] (core.east) -- (mshr.west);
  \draw[hw bus] (mshr.east) -- (3.75,0);
  \draw[hw bus] (8.6,0) -- (dram.west);
  % 返回通道：先下探再折回，竖直进入 core 底边
  \draw[hw bus] (dram.south) |- (4.9,-1.25) -| (core.south);
  \node[hw label, anchor=north] at (4.9,-1.35) {64 B 数据返回};
  \node[hw label, anchor=north] at (2.55,-0.72) {单核在途 1 KB $\Rightarrow$ 封顶 10.2 GB/s};
  \node[hw label, anchor=south] at (6.17,0.55) {整机 100 GB/s：管中需常有约 10 KB};
\end{pdfdiagram}
\diagramnote{Little 定律的直观版：时延 100 ns 的水管里要维持 100 GB/s，线上必须随时有约 10 KB 字节在途。单核的 MSHR 只有 16 条、在途预算 1 KB，于是单核带宽被闸门封顶在 $1\text{ KB}/100\text{ ns}\approx10.2$ GB/s——水管再粗，闸门只开这么宽。}

\begin{longtable}{L{0.16\linewidth}L{0.18\linewidth}L{0.18\linewidth}L{0.34\linewidth}}
\toprule
在途字节数 & 折合在途行 & 带宽上限（100 ns 时延） & 典型来源 \\
\midrule
256 B & 4 & 2.6 GB/s & 保守的顺序小核 \\\tablerowrule
512 B & 8 & 5.1 GB/s & 较小的 MSHR 配置 \\\tablerowrule
1 KB & 16 & 10.2 GB/s & 主流乱序核的 MSHR 量级 \\\tablerowrule
2 KB & 32 & 20.5 GB/s & 加强的 miss 处理 \\\tablerowrule
4 KB & 64 & 41.0 GB/s & 核加硬件预取缓冲合计 \\\tablerowrule
10 KB & 160 & 102.4 GB/s & 喂满整机带宽所需总量 \\
\bottomrule
\end{longtable}

这张表把乱序执行、MSHR、硬件预取三件事的统一本质说穿了：维持在途数。乱序让一条 miss 后面的独立指令继续走，走着走着往往又触发新的 miss，于是一个 miss 变成多个 miss 同时在途；MSHR 是在途请求的登记表，容量直接划定单核带宽上限；预取更彻底，把将来才会发生的 miss 提前变成现在的在途。存储与设备事务相关章节算预取提前量时见过同一个公式：miss 时延 200 拍、每 32 拍消化一行，提前量取 $200\times\frac{1}{32}\approx6.25$ 行、向上取 7 行——那正是“在途数等于时延乘消耗速率”的 Little 定律写法。教学计算机项目中的教学机则是另一个极端：它一次只发起一笔访存，发起到完成是一个完整总线周期，期间没有任何别的请求在飞，在途数恒为 1，于是吞吐被时延直接限定为 $1/\text{时延}$。它慢得坦白，也把“为什么在途数重要”衬得最清楚。

使用这条定律前先核对单位：带宽是“字节每秒”，时延是“秒”，乘积是“字节”——在途数是一个库存量，不是速率。任何“带宽乘时延”算出来不是字节的，都是单位代错了。

再把它翻译成一个日常类比。收银台每小时接待 60 位顾客（吞吐），每位顾客在店里平均待 6 分钟即 0.1 小时（时延），店里平均就有 $60\times0.1=6$ 位顾客（在途数）。想让店里人少，要么收银更快（提高吞吐、缩短停留），要么确实少来人——没有哪种办法能绕过这个乘积。内存系统同理：时延由 DRAM 物理决定，带宽目标由整机决定，唯一可调的旋钮就是在途数。

公式还可以反着用。已知单核在途预算 4 KB、DRAM 时延 100 ns，单核带宽上限是多少？$4\text{ KB}/100\text{ ns}\approx41$ GB/s。若工艺改进把时延降到 50 ns，同样 4 KB 的在途预算立刻值 82 GB/s——时延减半，同样的在途数能喂两倍的带宽。这正是各代内存一边提带宽、一边压时延的原因：两个参数通过同一个乘积耦合在一起。

\begin{lstlisting}
反向演算: 在途预算 4 KB
  时延 100 ns: 带宽上限 = 4096 B / 100 ns = 41 GB/s
  时延 50 ns:  带宽上限 = 4096 B / 50 ns  = 82 GB/s
结论: 在途预算固定时，带宽上限与时延成反比
\end{lstlisting}

把三种机器摆在一起，“在途数”这条线索就完整了：

\begin{longtable}{L{0.20\linewidth}L{0.18\linewidth}L{0.20\linewidth}L{0.28\linewidth}}
\toprule
机器 & 同时在途的访存 & 吞吐与时延的关系 & 关键问题 \\
\midrule
教学计算机项目 & 恒为 1 笔 & 吞吐 $=1/$ 时延 & 每个总线周期只做一笔，慢得坦白 \\\tablerowrule
简单顺序核 & 几笔 & 吞吐 $\approx$ 在途数 $/$ 时延 & MSHR 规模直接划定上限 \\\tablerowrule
乱序核加预取 & 几十到上百笔 & 同上，但在途数被刻意撑大 & 窗口、预取都在为在途数服务 \\
\bottomrule
\end{longtable}

从教学机到现代核，访存子系统的全部复杂化都可以读成一句话：在改变不了 DRAM 时延的前提下，把同时在途的请求数从 1 撑到上百。

\topic{内存墙：算力与带宽的历史分叉}
为什么要费这么大劲维持在途数？因为算力和带宽的差距是几十年复利滚出来的。教科书反复引用的量级估计是：处理器算力大约每年增长 50\%，DRAM 带宽（以及时延）大约每年改善 7\%。这两个数字是说明斜率差的量级，不是某一年的精确统计，但它们的复利后果必须算出来看：

\begin{longtable}{L{0.16\linewidth}L{0.24\linewidth}L{0.24\linewidth}L{0.22\linewidth}}
\toprule
经过年数 & 算力倍数（年增 50\%） & 带宽倍数（年增 7\%） & 算力/带宽差距 \\
\midrule
0 & 1 & 1 & 1 \\\tablerowrule
5 & 7.6 & 1.40 & 约 5.4 \\\tablerowrule
10 & 57.7 & 1.97 & 约 29 \\\tablerowrule
15 & 437.9 & 2.76 & 约 159 \\\tablerowrule
20 & 3325 & 3.87 & 约 859 \\
\bottomrule
\end{longtable}

差距本身也是一条复利曲线：每年乘以 $1.5/1.07\approx1.4$，十年约 29 倍，二十年约 859 倍。“内存墙”不是某一年突然砌起来的墙，而是两条指数曲线的分叉口。它的后果链条在整机结构里随处可见：随着时延差扩大，cache 从一级扩展到多级且容量不断增加；随着带宽差扩大，预取器、乱序窗口和 MSHR 的规模也相应增加。这套机制只能在程序具有局部性时缓解内存墙的影响，无法消除处理器与存储器之间的速度差距。

缓存无法隐藏所有负载的存储时延：没有复用的流式负载直接受到内存墙限制，视频处理仅顺序访问一次帧数据，AI 推理仅顺序读取一次权重，大数据扫描也仅顺序读取一次数据表；工作集一旦超过 cache 容量，miss 率回升，内存墙效应重新显现；预取本身消耗带宽，错误预取是纯浪费，存储与设备事务相关章节已经量化过这一权衡。按局部性对负载分类后，可以明确区分内存墙对不同负载的影响：

\begin{longtable}{L{0.24\linewidth}L{0.20\linewidth}L{0.22\linewidth}L{0.22\linewidth}}
\toprule
负载类型 & 局部性 & 墙的表现 & 有效手段 \\
\midrule
稠密数值（分块矩阵） & 高，复用充分 & 被 cache 遮住 & 分块、向量化 \\\tablerowrule
指针追踪（链表、图） & 差，访问不可预测 & 时延墙全暴露 & 布局重排、池化分配 \\\tablerowrule
流式扫描（视频、日志） & 空间有，时间无 & 带宽墙全暴露 & 预取、压缩、减少趟数 \\\tablerowrule
随机小消息（网络服务） & 几乎无 & 双墙齐撞 & 批量、合并、缓存热集 \\
\bottomrule
\end{longtable}

这张表也预告了“性能测量与案例”一节的结论：对局部性差的负载，优化手段几乎全部围绕“少移动数据”展开，因为墙的那一侧没有更快的零件可买——DRAM 的 7\% 不会因为你的程序着急而变成 50\%。

这条历史曲线要读对两处。第一，50\% 和 7\% 是斜率不是现状：它们描述“差距为什么越拉越大”，不描述“今天差多少倍”——今天的绝对差距取决于各代产品的具体参数，但斜率差决定了这个差距明年一定更大。第二，7\% 指的主要是 DRAM 核心的访问时延与单位引脚带宽：行激活、预充电这些物理动作受电容充放电约束，不是堆晶体管就能解决的，这正是 DRAM 与逻辑工艺走势分叉的根源。

\begin{lstlisting}
复利演算(年增 50% vs 年增 7%):
  算力: 1.5^10 = 57.7      带宽: 1.07^10 = 1.97
  差距本身的年增长率: 1.5/1.07 = 1.40
  10 年: 1.40^10 ≈ 29      20 年: (1.5/1.07)^20 ≈ 859
  读法: 每过 5 年，同一程序里"等内存"的时间份额再上一档
\end{lstlisting}

把这两行复利画成半对数曲线，“分叉口”就成了两条直线的夹角：

\begin{pdfdiagram}
  % 半对数坐标：x 每 0.42 为一年（0-20 年），y = log10(倍数) x 1.2
  \draw[BookMuted, line width=0.6pt, ->] (0,0) -- (9.0,0);
  \draw[BookMuted, line width=0.6pt, ->] (0,0) -- (0,4.8);
  \node[hw label, anchor=north] at (4.2,-0.6) {经过年数};
  \node[hw label, anchor=south west] at (0,4.85) {相对倍数（对数轴）};
  \foreach \x/\t in {0/0, 2.1/5, 4.2/10, 6.3/15, 8.4/20} {
    \draw[BookMuted, line width=0.5pt] (\x,0.06) -- (\x,-0.06);
    \node[hw label, anchor=north] at (\x,-0.1) {\t};
  }
  \foreach \y/\t in {0/1, 1.2/10, 2.4/100, 3.6/1000} {
    \draw[BookMuted, line width=0.5pt] (0.06,\y) -- (-0.06,\y);
    \node[hw label, anchor=east] at (-0.12,\y) {\t};
  }
  % 算力 1.5^n 与带宽 1.07^n：对数轴上都是直线
  \draw[BookAccent, line width=1pt] (0,0) -- (8.4,4.23);
  \draw[BookTeal, line width=1pt] (0,0) -- (8.4,0.71);
  \node[hw label, anchor=west] at (8.45,4.23) {算力 $\times$3325};
  \node[hw label, anchor=west] at (8.45,0.71) {带宽 $\times$3.87};
  \fill[BookAccent] (2.1,1.06) circle (1.4pt);
  \fill[BookAccent] (4.2,2.11) circle (1.4pt);
  \fill[BookAccent] (6.3,3.17) circle (1.4pt);
  \fill[BookTeal] (2.1,0.18) circle (1.4pt);
  \fill[BookTeal] (4.2,0.35) circle (1.4pt);
  \fill[BookTeal] (6.3,0.53) circle (1.4pt);
  % 差距与遮羞手段（位于两线之间的楔形区）
  \draw[{Stealth[length=1.8mm]}-{Stealth[length=1.8mm]}, BookRed, line width=0.6pt] (8.4,0.71) -- (8.4,4.23);
  \node[hw label, anchor=east] at (8.25,1.6) {20 年：约 859 倍};
  \node[hw label] at (2.7,0.62) {cache 遮时延};
  \node[hw label] at (4.6,1.15) {预取、MSHR 撑在途};
  \node[hw label] at (6.9,2.2) {多核凑在途};
\end{pdfdiagram}
\diagramnote{内存墙的历史分叉（半对数坐标）：算力年增 50\%、DRAM 带宽年增 7\%，两条指数曲线在对数轴上都是直线，斜率差就是分叉角。20 年后算力领先约 859 倍；楔形区里的多级 cache、预取与 MSHR、多核，是整机绕墙而行的遮羞手段——墙本身还在原地，且每年长高一点。}

墙在整机设计上留下的痕迹，按出现顺序排大致是：先加 cache（遮时延），再加预取（遮带宽），再加深乱序窗口和 MSHR（撑在途数），再堆多核（用核数凑在途数）。每一步在前面两节的公式里都有对应的位置：

\begin{longtable}{L{0.24\linewidth}L{0.24\linewidth}L{0.36\linewidth}}
\toprule
设计手段 & 遮的是哪一面墙 & 在本节公式里的位置 \\
\midrule
多级 cache & 时延差 & 把有效时延从 100 ns 压向几 ns \\\tablerowrule
硬件预取 & 带宽差 & 把将来的 miss 变成现在的在途数 \\\tablerowrule
乱序窗口、MSHR & 在途数不足 & 直接抬升单核的在途预算 \\\tablerowrule
多核 & 单核在途数有限 & 用核数乘以每核在途数 \\
\bottomrule
\end{longtable}

注意这张表里没有一行是“让 DRAM 本身变快”——整机的全部聪明都花在绕过墙上，墙本身还在原地，且每年长高一点。

也要公平地说，工业界并非没有正面攻过墙，只是攻法都带限定条件：

\begin{longtable}{L{0.24\linewidth}L{0.30\linewidth}L{0.32\linewidth}}
\toprule
修补手段 & 做法 & 改变了什么、没改变什么 \\
\midrule
宽接口存储（HBM 等） & 把存储堆叠到逻辑芯片近旁，引脚以千计 & 带宽密度抬一个量级；容量与成本受限，拐点位置变了，墙还在 \\\tablerowrule
加大片上缓存 & 把更大的工作集留在片内 & 可延后内存墙出现；无局部性的负载仍受内存墙限制 \\\tablerowrule
近存计算 & 把部分运算搬进存储侧 & 省掉来回搬运；适用面窄，编程模型未定型 \\
\bottomrule
\end{longtable}

这些手段的共同点是抬高局部的带宽密度，而不是抹平斜率差：逻辑与存储的复利曲线没有因此合拢，分叉仍在继续。

\topic{roofline 模型}
roofline 模型把前面的讨论收成一条拐线。先给定义：算术强度（arithmetic intensity）$I$ 等于程序完成的浮点操作数除以它从 DRAM 搬运的字节数，单位 FLOP/B，这是程序的属性；可达性能等于峰值算力与“带宽乘算术强度”两者的较小者，这是机器给的上限；两者的交点叫拐点，拐点强度 $I^{*}$ 等于峰值算力除以带宽，这是机器的属性。

\begin{lstlisting}
可达性能(FLOP/s) = min( 峰值算力, 内存带宽 * 算术强度 )
拐点强度 I* = 峰值算力 / 内存带宽
I < I*: 带宽瓶颈(拐线左侧)，性能 = 带宽 * I，随强度线性上升
I > I*: 算力瓶颈(拐线右侧)，性能 = 峰值算力，平顶
\end{lstlisting}

给一台具体机器：峰值算力 100 GFLOP/s，内存带宽 25 GB/s，拐点 $I^{*}=100/25=4$ FLOP/B。取两个 kernel：甲的算术强度 0.25 FLOP/B，乙的算术强度 16 FLOP/B。甲：带宽允许 $25\times0.25=6.25$ GFLOP/s，小于峰值 100，可达性能 6.25 GFLOP/s，带宽瓶颈，峰值只用到 6.25\%；乙：带宽允许 $25\times16=400$ GFLOP/s，超过峰值，可达性能 100 GFLOP/s，算力瓶颈，峰值用满。

\begin{longtable}{L{0.10\linewidth}L{0.16\linewidth}L{0.20\linewidth}L{0.16\linewidth}L{0.16\linewidth}L{0.10\linewidth}}
\toprule
kernel & 算术强度 & 带宽允许上限 & 峰值算力 & 可达性能 & 瓶颈 \\
\midrule
甲 & 0.25 FLOP/B & 6.25 GFLOP/s & 100 GFLOP/s & 6.25 GFLOP/s & 带宽 \\\tablerowrule
乙 & 16 FLOP/B & 400 GFLOP/s & 100 GFLOP/s & 100 GFLOP/s & 算力 \\
\bottomrule
\end{longtable}

把这台机器画成 roofline 图，两个 kernel 落在哪一侧、离屋顶多远，一眼可读：

\begin{pdfdiagram}
  % 双对数坐标：x = (log10 I + 2)*2.3，y = (log10 P + 1)*1.15
  \draw[BookMuted, line width=0.6pt, ->] (0,0) -- (9.8,0);
  \draw[BookMuted, line width=0.6pt, ->] (0,0) -- (0,5.0);
  \node[hw label, anchor=north] at (4.6,-0.62) {算术强度 $I$（FLOP/B，对数）};
  \node[hw label, rotate=90, anchor=south] at (-1.0,2.3) {可达性能（GFLOP/s，对数）};
  \foreach \x/\t in {0/0.01, 2.3/0.1, 4.6/1, 6.9/10, 9.2/100} {
    \draw[BookMuted, line width=0.5pt] (\x,0.06) -- (\x,-0.06);
    \node[hw label, anchor=north] at (\x,-0.1) {\t};
  }
  \foreach \y/\t in {0/0.1, 1.15/1, 2.3/10, 3.45/100, 4.6/1000} {
    \draw[BookMuted, line width=0.5pt] (0.06,\y) -- (-0.06,\y);
    \node[hw label, anchor=east] at (-0.12,\y) {\t};
  }
  % 带宽斜线 P = 25*I（I 从 0.01 到拐点 4），算力屋顶 P = 100
  \draw[BookAccent, line width=1pt] (0,0.458) -- (5.985,3.45);
  \draw[BookAccent, line width=1pt] (5.985,3.45) -- (9.2,3.45);
  \draw[BookTeal, dashed, line width=0.5pt] (5.985,0) -- (5.985,3.45);
  \draw[BookInk, line width=0.7pt, fill=white] (5.985,3.45) circle (1.6pt);
  \node[hw label, anchor=south east] at (5.85,3.56) {$I^{*}=4$};
  \node[hw label, anchor=west] at (0.5,0.3) {带宽斜线 25 GB/s};
  \node[hw label, anchor=west] at (6.15,3.8) {算力屋顶 100 GFLOP/s};
  \node[hw label] at (4.6,0.72) {带宽瓶颈区};
  \node[hw label] at (8.3,2.35) {算力瓶颈区};
  % 两个 kernel 落点
  \fill[BookRed] (2.077,1.496) circle (1.6pt);
  \node[hw label, anchor=west] at (2.28,1.42) {向量加法 $I\approx0.08$};
  \fill[BookRed] (7.37,3.45) circle (1.6pt);
  \node[hw label, anchor=north] at (7.37,3.26) {矩阵乘 $I\approx16$};
\end{pdfdiagram}
\diagramnote{roofline 模型（双对数坐标）。斜线是带宽屋顶，水平线是算力屋顶，交点即拐点 $I^{*}=4$ FLOP/B。向量加法 $I\approx0.08$ 落在带宽瓶颈区，可达约 2 GFLOP/s；分块矩阵乘 $I\approx16$ 落在算力瓶颈区，用满峰值。}

左右两侧的优化顺序完全不同。左侧的 kernel 缺的是字节：先动数据——布局、分块、压缩、减少遍历趟数——把 $I$ 往右推；右侧的 kernel 缺的是运算吞吐：这时才轮到动指令——向量化、FMA、多核——把峰值用起来。在左侧做向量化的结果，是把等待带宽的时间花得更快；在右侧做数据布局的收益，早就拿完了。roofline 最常见的误用是拿峰值算力评估左侧 kernel：“机器标称 100 GFLOP/s，我的程序怎么只有 6？”——因为程序在拐线左侧，峰值与它无关。章首表格里“拿峰值算力当可达性能”那条误用，指的就是这件事。

屋顶还会随机器移动：换一台机器，拐点位置就变了。三台机器的对照可以直接算：

\begin{longtable}{L{0.10\linewidth}L{0.16\linewidth}L{0.14\linewidth}L{0.14\linewidth}L{0.16\linewidth}L{0.16\linewidth}}
\toprule
机器 & 峰值算力 & 带宽 & 拐点强度 & 甲（0.25） & 乙（16） \\
\midrule
A & 100 GFLOP/s & 25 GB/s & 4 FLOP/B & 6.25，带宽 & 100，算力 \\\tablerowrule
B & 200 GFLOP/s & 25 GB/s & 8 FLOP/B & 6.25，带宽 & 200，算力 \\\tablerowrule
C & 100 GFLOP/s & 50 GB/s & 2 FLOP/B & 12.5，带宽 & 100，算力 \\
\bottomrule
\end{longtable}

甲在机器 B 上毫无变化——算力翻倍与它无关；在机器 C 上直接翻倍——带宽才是它的屋顶。乙正好相反。一句话总结采购逻辑：左侧 kernel 为带宽付费，右侧 kernel 为算力付费，升级机器之前先定位 kernel 在拐线的哪一侧。

把三台机器的屋顶画进同一张图，同一 kernel 的位置随机器怎么移动，就直接可见：

\begin{pdfdiagram}
  % 双对数坐标：x = (log10 I + 2)*2.3，y = (log10 P + 1)*1.15
  \draw[BookMuted, line width=0.6pt, ->] (0,0) -- (9.8,0);
  \draw[BookMuted, line width=0.6pt, ->] (0,0) -- (0,5.0);
  \node[hw label, anchor=north] at (4.6,-0.62) {算术强度 $I$（FLOP/B，对数）};
  \node[hw label, rotate=90, anchor=south] at (-1.0,2.3) {可达性能（GFLOP/s，对数）};
  \foreach \x/\t in {0/0.01, 2.3/0.1, 4.6/1, 6.9/10, 9.2/100} {
    \draw[BookMuted, line width=0.5pt] (\x,0.06) -- (\x,-0.06);
    \node[hw label, anchor=north] at (\x,-0.1) {\t};
  }
  \foreach \y/\t in {0/0.1, 1.15/1, 2.3/10, 3.45/100} {
    \draw[BookMuted, line width=0.5pt] (0.06,\y) -- (-0.06,\y);
    \node[hw label, anchor=east] at (-0.12,\y) {\t};
  }
  % 机器 A：峰值 100、带宽 25
  \draw[BookAccent, line width=1pt] (0,0.458) -- (5.985,3.45) -- (9.2,3.45);
  % 机器 B：峰值 200、带宽 25——共享带宽斜线，屋顶更高
  \draw[BookRed, line width=1pt] (5.985,3.45) -- (6.677,3.796) -- (9.2,3.796);
  % 机器 C：峰值 100、带宽 50——屋顶与 A 重合，以虚线示出
  \draw[BookTeal, line width=1pt] (0,0.804) -- (5.292,3.45);
  \draw[BookTeal, line width=1pt, dashed] (5.292,3.45) -- (9.2,3.45);
  % 图例
  \draw[BookAccent, line width=1pt] (0.35,4.75) -- (1.0,4.75);
  \node[hw label, anchor=west] at (1.15,4.75) {A：峰值 100 · 带宽 25};
  \draw[BookRed, line width=1pt] (0.35,4.4) -- (1.0,4.4);
  \node[hw label, anchor=west] at (1.15,4.4) {B：峰值 200 · 带宽 25};
  \draw[BookTeal, line width=1pt] (0.35,4.05) -- (1.0,4.05);
  \node[hw label, anchor=west] at (1.15,4.05) {C：峰值 100 · 带宽 50};
  % kernel 甲 I=0.25（x=3.216）：A/B 上 6.25，C 上 12.5
  \fill[BookInk] (3.216,2.065) circle (1.6pt);
  \node[hw label, anchor=west] at (3.4,1.8) {甲在 A/B：6.25};
  \fill[BookTeal] (3.216,2.412) circle (1.6pt);
  \node[hw label, anchor=east] at (3.05,2.6) {甲在 C：12.5};
  % kernel 乙 I=16（x=7.37）：A/C 上 100，B 上 200
  \fill[BookInk] (7.37,3.45) circle (1.6pt);
  \node[hw label, anchor=north] at (7.37,3.26) {乙在 A/C：100};
  \fill[BookRed] (7.37,3.796) circle (1.6pt);
  \node[hw label, anchor=west] at (7.6,3.95) {乙在 B：200};
\end{pdfdiagram}
\diagramnote{三台机器的 roofline 对照（双对数坐标）。A 与 B 共享同一条带宽斜线，B 的算力屋顶更高；C 与 A 峰值算力相同、带宽斜线更陡，C 的屋顶与 A 重合，以虚线示出。同一个 kernel 甲（$I=0.25$）在 A、B 上都是 6.25 GFLOP/s——算力翻倍与它无关；在 C 上变成 12.5——带宽才是它的屋顶。kernel 乙（$I=16$）正好相反：A、C 上 100，B 上 200。升级机器之前，先定位 kernel 在拐线的哪一侧。}

算力屋顶其实不止一条：标量峰值、向量峰值、单核峰值、多核峰值各是一条水平线，而带宽屋顶通常只有一条。把算力屋顶抬高（加向量单元、加核），拐点就向右移，原先在右侧的 kernel 可能被划到左侧。看一个具体例子：给机器 A 再加三个核，峰值算力变 400 GFLOP/s，带宽仍是 25 GB/s，拐点从 4 移到 16 FLOP/B。

\begin{longtable}{L{0.12\linewidth}L{0.16\linewidth}L{0.14\linewidth}L{0.14\linewidth}L{0.16\linewidth}L{0.16\linewidth}}
\toprule
配置 & 峰值算力 & 带宽 & 拐点强度 & 乙（$I=16$） & 丙（$I=10$） \\
\midrule
单核 & 100 GFLOP/s & 25 GB/s & 4 FLOP/B & 100，算力 & 100，算力 \\\tablerowrule
四核 & 400 GFLOP/s & 25 GB/s & 16 FLOP/B & 400，恰在拐点 & 250，带宽 \\
\bottomrule
\end{longtable}

kernel 丙在单核机器上是算力瓶颈，换到四核机器上反而成了带宽瓶颈：可达性能从 100 涨到 250，但离 400 的新峰值还差一截，而且再怎么优化指令也摸不到 400——带宽 25 GB/s 乘以强度 10 只允许 250。升级机器把瓶颈也升级了：算力侧的投入越多，剩下的 kernel 越集中到带宽一侧，这正是内存墙在 roofline 图上的样子。

另一个反例同样值得记住：向量点积 \code{c += a[i]*b[i]}，每对元素 2 次浮点操作、读 8 字节，算术强度 0.25 FLOP/B，与矩阵规模无关。给这种 kernel 做向量化，可达性能一格不涨——它在拐线左侧，上限是 $25\times0.25=6.25$ GFLOP/s，向量指令只让“等待”执行得更快。它的唯一出路是改变算法本身：把点积融合进一个更大的、有复用的计算里，人为抬高算术强度。

再补一个有复用但复用不足的中间态例子：五点 stencil（图像模糊、流体网格里到处都是）。每个输出点读 4 个邻居加自身、写回 1 个点，约 6 次浮点操作；若完全无复用，每点读 5 个 float、写 1 个，共 24 B，强度 $6/24=0.25$ FLOP/B；cache 利用行间复用后，平均每个输出点只需从 DRAM 新读约 1 个元素再加写回，约 8 B，强度升到约 0.75 FLOP/B。在机器 A 上可达 $25\times0.75\approx18.8$ GFLOP/s——比点积好三倍，仍然在拐线左侧。

\begin{lstlisting}
五点 stencil 在机器 A 上:
  无复用:   6 FLOP / 24 B = 0.25 FLOP/B  ->  6.25 GFLOP/s
  行间复用: 6 FLOP / 8 B  = 0.75 FLOP/B  ->  18.8 GFLOP/s
  结论: cache 帮了三倍，瓶颈仍是带宽(0.75 < 4)
\end{lstlisting}

点积、stencil、分块矩阵乘恰好排成算术强度的三个台阶：0.25（无复用）、0.75（行间复用）、16（块内复用）。三个台阶的名字都一样——“用了 cache”——效果差出两个数量级；问“cache 有没有用”之前，先问复用发生在哪一行、哪一块。

最后交代 roofline 模型的使用边界：它假设 kernel 处于稳态、算术强度按 DRAM 流量计算、忽略时延对短 kernel 的影响。启动阶段的冷 cache、毫秒级的小 kernel、时延受限的指针追踪，都不在这条拐线的适用范围里；那些场合要回到 Little 定律和“Amdahl 定律与加速极限”一节的时间拆分，一项一项算。

\topic{案例：矩阵乘法两个版本的算术强度}
同一数学、两种写法，是演示 roofline 的经典案例。$n\times n$ 单精度矩阵乘 $C=A\times B$，数学上是 $2n^3$ 次浮点操作，这一点两个版本完全相同；不同的是数据从 DRAM 走几趟，也就是算术强度。

\begin{lstlisting}
// 注释（推进）：每轮只推进一个元素或一个控制步，并保持中间状态的一致性。
// 注释（边界）：退出条件成立后才提交结果；空集合、越界和失败要单独处理。
朴素版(i-j-k 顺序, float):
  for (i) for (j) for (k)
      c[i][j] += a[i][k] * b[k][j];
  内层每迭代: 2 FLOP；若无 cache 复用，读 a、b 各 4 B，共 8 B
  算术强度 = 2 FLOP / 8 B = 0.25 FLOP/B

分块版(64 x 64 子块驻留 cache):
  for (ii) for (jj) for (kk)         // 步长 64
      三个 64x64 子块驻留 cache
      for (块内 i) for (块内 j) for (块内 k)
          c[i][j] += a[i][k] * b[k][j];
  每对分块: 2*64^3 = 524288 FLOP
           DRAM 读入 2*64*64*4 B = 32768 B = 32 KB
  算术强度 = 524288 / 32768 = 16 FLOP/B
\end{lstlisting}

朴素版的强度 0.25 FLOP/B 与矩阵规模无关：矩阵再大，每次乘法加法还是要等新的字节，它永远停在拐线左侧，在机器 A 上只拿 6.25 GFLOP/s。分块版的强度 16 FLOP/B 来自复用：一个 $64\times64$ 的块载入一次，参与 64 次乘加，字节被反复使用；三个块的工作集 $3\times16$ KB $=48$ KB，驻留 cache 毫无压力。注意一个整齐的关系：分块边长是 64，强度恰好从 0.25 提到 16，提升倍数就是分块边长——每维复用次数直接折算成算术强度。

把两个版本的数据流画出来，0.25 与 16 的差距就是字节走几趟的差距：

\begin{pdfdiagram}
  % 左：朴素版，每迭代都从 DRAM 新取
  \node[hw label] at (1.6,3.6) {朴素版：每迭代新取 8 B};
  \node[hw compute, minimum width=2.2cm, minimum height=1.0cm] (c1) at (1.6,2.4) {乘加\\2 FLOP};
  \node[hw memory, minimum width=2.2cm, minimum height=0.9cm] (d1) at (1.6,0.4) {DRAM};
  \draw[hw bus] (d1.north) -- (c1.south);
  \node[hw label, anchor=west] at (1.75,1.4) {每次迭代 8 B};
  \node[hw label, anchor=north] at (1.6,-0.3) {算术强度 0.25 FLOP/B};
  % 右：64 分块，块驻留 cache 反复使用
  \node[hw label] at (7.3,4.4) {64 分块：载入一次，复用 64 次};
  \node[hw compute, minimum width=2.4cm, minimum height=1.0cm] (c2) at (7.3,3.2) {块内乘加\\$2\times64^3$ FLOP};
  \node[hw memory, minimum width=2.4cm, minimum height=0.9cm, draw=BookAmber] (ca) at (7.3,1.5) {cache：3 块 48 KB};
  \node[hw memory, minimum width=2.4cm, minimum height=0.9cm] (d2) at (7.3,-0.2) {DRAM};
  \draw[hw bus] (d2.north) -- (ca.south);
  \node[hw label, anchor=west] at (7.5,0.62) {每对块 32 KB，仅一次};
  \draw[hw bus] (ca.north) -- (c2.south);
  \node[hw label, anchor=west] at (7.5,2.32) {块内每元素复用 64 次};
  \node[hw label, anchor=north] at (7.3,-0.9) {算术强度 16 FLOP/B};
\end{pdfdiagram}
\diagramnote{同一段矩阵乘的两种数据流。左：朴素版内层每次乘加都从 DRAM 新取两个元素，2 FLOP 配 8 B，强度 0.25 FLOP/B，与矩阵规模无关。右：$64\times64$ 分块后三个块共 48 KB 驻留 cache，每对块从 DRAM 读入 32 KB 只有一次，块内完成 $2\times64^3$ 次乘加——字节被反复使用，强度升到 16 FLOP/B。优化没有改动任何一次浮点运算，只改了字节走的路程。}

\begin{longtable}{L{0.14\linewidth}L{0.22\linewidth}L{0.14\linewidth}L{0.16\linewidth}L{0.14\linewidth}L{0.08\linewidth}}
\toprule
版本 & DRAM 流量 & 算术强度 & 机器 A 可达 & 峰值利用 & 瓶颈 \\
\midrule
朴素三重循环 & 每迭代 8 B & 0.25 FLOP/B & 6.25 GFLOP/s & 6.25\% & 带宽 \\\tablerowrule
64 分块 & 每 524288 FLOP 仅 32 KB & 16 FLOP/B & 100 GFLOP/s & 100\% & 算力 \\
\bottomrule
\end{longtable}

同一段数学，位于屋顶两侧。这次优化没有改动任何一个浮点运算，只改了字节走的路程——这就是存储与设备事务相关章节局部性原理在性能侧的兑现：cache 的价值不是“让访存变快”，而是让算术强度变大，把 kernel 从拐线左侧搬到右侧。工程次序也因此固定：先做分块把数据侧的问题解决，再谈向量化和多核把算力侧的峰值用满；顺序反了，就是把左侧 kernel 的等待做得更快。

64 这个数字不是魔法，它来自一个通式。设分块边长为 $B$，每对分块做 $2B^3$ 次浮点操作，从 DRAM 读入两个 $B\times B$ 块共 $8B^2$ 字节，强度就是 $B/4$；代价是三个分块的工作集 $12B^2$ 字节必须驻留 cache。把几档 $B$ 摆出来：

\begin{lstlisting}
分块边长 B 的通式(单精度):
  每对分块 FLOP  = 2B^3
  DRAM 流量      = 2 * B^2 * 4 = 8B^2 字节
  算术强度 I     = 2B^3 / 8B^2 = B/4
  工作集         = 3 个分块 = 12B^2 字节
  B = 32:  I = 8,  工作集 12 KB
  B = 64:  I = 16, 工作集 48 KB
  B = 128: I = 32, 工作集 192 KB
约束: 工作集必须驻留 cache，否则复用落空，强度退回左侧
\end{lstlisting}

强度随 $B$ 线性上涨，工作集随 $B^2$ 上涨，选 $B$ 就是在这两条曲线之间找最大可行点：$B$ 取到工作集刚好能驻留的那一档为止，再大一档，块本身开始被挤出 cache，复用落空，强度不但不涨反而退回左侧。这与存储与设备事务相关章节“cache 容量决定工作集上限”是同一条约束的两种说法。

还有一个容易混淆的优化要分清：不调分块、只调换三重循环的次序（例如 i-k-j），能让内层循环改为沿行顺序访问，空间局部性改善、冲突 miss 减少，朴素版的 0.25 FLOP/B 会变得更“扎实”；但换序不产生跨迭代的批量复用，强度的量级不变，kernel 仍在拐线左侧。换序是把 0.25 做扎实，分块才是把 0.25 变成 16——前者修访问模式，后者改数据流量，两件不同的事。最后从测量侧看一眼：两个版本的浮点操作数完全一样，性能计数器里不同的只是 cache miss 次数和 DRAM 流量；“性能测量与案例”一节的测量工具会把这个差别直接摆成数字。

\topic{数据移动成本表}
矩阵乘法的例子不是特例，它背后是整机存储结构的统一图景。把全书零散出现过的访问时延收进一张表，按数量级记——这些是每个硬件工程师该能脱口而出的预算数：

\begin{longtable}{L{0.12\linewidth}L{0.14\linewidth}L{0.14\linewidth}L{0.14\linewidth}L{0.32\linewidth}}
\toprule
层级 & 典型时延 & 相对寄存器 & 比上一级慢 & 一句话 \\
\midrule
寄存器 & 约 0.5 ns & 1 倍 & — & 几乎不计价，但容量以百计 \\\tablerowrule
L1 cache & 约 1 ns & 2 倍 & 2 倍 & 容量以万字节计 \\\tablerowrule
L2 cache & 约 5 ns & 10 倍 & 5 倍 & 容量以十万字节计 \\\tablerowrule
DRAM & 约 100 ns & 200 倍 & 20 倍 & 容量以十亿字节计 \\\tablerowrule
SSD & 约 100 µs & 200000 倍 & 1000 倍 & 第一级“出板”的存储 \\\tablerowrule
网络 & 约 1 ms & 2000000 倍 & 10 倍 & 第一级“出机器”的存储 \\
\bottomrule
\end{longtable}

读这张表要看相邻级之间的断崖：不是百分之几十的差别，而是 2 倍到 1000 倍的跳变，DRAM 到 SSD 那一级最显著，整整三个数量级——所以一次意外的缺页换页至今仍是事故级事件。能耗图景与时延同构：教科书的能耗表里，把一个字送到片外再取回的代价，同样比一次片上运算高出几个数量级。时延和能耗两条曲线指向同一个结论：“优化等于把数据推近”，这是大多数性能工作的实质。寄存器分配把热变量推进寄存器，分块把热数据推进 L1 和 L2，预取把将来的访问提前推进 cache，压缩用更少字节把同样的信息推过带宽，批量 I/O 把一千次小读换成一次大读——手段各异，方向全部相同：沿这张表向上走。

沿这张表向上走的具体手段，按目标层级各归其位：

\begin{longtable}{L{0.16\linewidth}L{0.32\linewidth}L{0.32\linewidth}}
\toprule
目标层级 & 把数据推近的手段 & 关键问题 \\
\midrule
寄存器 & 循环展开、寄存器分配、减少溢出 & 编译器与手工安排的分工 \\\tablerowrule
L1/L2 & 分块、布局重排、结构体拆分 & 工作集是否小于容量 \\\tablerowrule
DRAM & 预取、压缩、减少遍历趟数 & 在途数是否喂满带宽 \\\tablerowrule
SSD & 顺序批量代替随机小读 & 是否避免了同步等待 \\\tablerowrule
网络 & 合并消息、缓存热集、异步流水 & 是否把等待藏进计算 \\
\bottomrule
\end{longtable}

上面那张时延成本表的日常用法是心算预算。$1\,\mu$s 的预算里允许约 10 次 DRAM 访问、约 1000 次 L1 访问、约 2000 次寄存器操作；一个要求 1 ms 内应答的在线服务，理论上最多约一万次 DRAM 访问，任何一次 SSD 读就直接占用预算的十分之一。性能评审时先问“这 1 ms 里各级访问各有多少次”，比先问“算法复杂度多少”更接近瓶颈。教学计算机项目中的教学机在这张表里的位置也值得一看：它没有 cache，每次访存都按总线周期全额付费，等于所有访问都停在 DRAM 那一行。它能用 74HC 器件搭出来，是因为慢对它不是缺陷而是教学目标；现代机把同样的访问藏进表的前三行，才有了本节全部问题的起点。

数量级还差一个直觉。把寄存器的一次访问当成 1 秒，其余层级按同一比例放大，这张表立刻变成日常经验：

\begin{longtable}{L{0.12\linewidth}L{0.16\linewidth}L{0.20\linewidth}L{0.36\linewidth}}
\toprule
层级 & 真实时延 & 人类尺度 & 直观画面 \\
\midrule
寄存器 & 约 0.5 ns & 1 秒 & 抬手就到 \\\tablerowrule
L1 cache & 约 1 ns & 2 秒 & 伸手拿桌上的东西 \\\tablerowrule
L2 cache & 约 5 ns & 10 秒 & 起身去书架 \\\tablerowrule
DRAM & 约 100 ns & 3 分半 & 下楼去仓库取货 \\\tablerowrule
SSD & 约 100 µs & 约 2.3 天 & 等一趟跨省快递 \\\tablerowrule
网络 & 约 1 ms & 约 23 天 & 等一班跨洋海运 \\
\bottomrule
\end{longtable}

记住“3 分半”和“23 天”这两个数，很多设计决策就不用再讨论：为取一个值等 3 分半，当然要把顺路能拿的都拿上（cache 行一次取 64 B 而不是 4 B）；等 23 天的请求，当然不能在原地空等（网络请求必须异步、批量）。前四行由硬件自动管理——cache、预取器、MMU 替软件做决定；后两行的决策必须由软件显式做出——I/O 调度、批量、异步、流水。数量级的断崖划出了软硬件分工的边界：断崖之上，同步等待可以接受；断崖之下，同步等待就是事故。

预算分配演算把分工再具体化。一个 1 ms 应答的在线服务，一种可行的分配是：

\begin{lstlisting}
1 ms 在线服务预算的一种分配:
  DRAM 访问 2000 次: 200 us
  L2 级访问 10000 次: 50 us
  留给计算与其余开销: 750 us
  红线一: 任何一次 SSD 读(约 100 us)占用预算的十分之一
  红线二: 任何一次跨机同步调用(约 1 ms)直接超支
\end{lstlisting}

这张预算表也解释了为什么高性能服务的架构评审总围绕“热数据是否在内存里、热路径是否不出本机”展开——这并非风格偏好，而是由访问时延的数量级差异决定的。

\section{六、性能测量与案例}

测量不是“运行一次并记录墙钟时间”。同样是“测一下”，墙钟计时、性能计数器、采样分析看到的是三个不同的世界；测到的数字里混着冷启动、随机噪声和混淆变量；连“两组数字能不能比”都有一份要先过的清单。

\topic{测量的三个层次}
测量手段按“能看到什么”分三层，三层回答三个不同的问题，谁也不能替代谁。

第一层是墙钟计时（wall-clock time）：在程序起点和终点各读一次时钟求差，得到第一节定义的时延。它回答“用户等了多久”，这是性能的最终定义——所有其他测量都只是为它找解释。

墙钟计时的代码只有三行，任何语言都写得出：

\begin{lstlisting}
t0 = now();          /* 读单调时钟 */
work();
t1 = now();
elapsed = t1 - t0;   /* 墙钟时延，含全部系统噪声 */
\end{lstlisting}

墙钟的优点是定义简单、无处可藏；缺点是它把这段时间里系统发生的全部事情都计了进去：操作系统把 CPU 调度给别的进程、中断处理、定时器滴答、主频按需升降、邻居程序把 cache 挤掉，全部算在读数里。所以墙钟的精度上限不在时钟本身——现代系统的时钟分辨率早已到纳秒级——而在系统噪声：同一程序连测十次，读数差出 10\% 是常态。

第二层是性能计数器（performance counter）：CPU 内部一组可编程的硬件寄存器，先选定事件类型（退休一条指令、一次 cache miss、一次分支误测），硬件此后每检测到一次该事件就加一。它回答“硬件发生了什么”，精确到单个事件；被测进程不在 CPU 上时，它的事件本来就不发生，所以读数几乎不受调度和后台负载污染。

计数器的成本在两头：一是必须先知道该看哪个事件——几十个事件名摆在面前，没有假设就无从下手；二是计数器的物理数量有限，典型同时可编程 4 到 8 个，想看的事件多于这个数就得分几轮测，工具自动多路复用时读数还会带上“本轮未全程计数”的标记。

第三层是采样分析（sampling profiler）：以固定频率打断程序，记录当时的程序计数器和调用栈，事后按函数归并样本。它回答“时间花在哪段代码”：占样本 70\% 的函数就是热点，优化从那里开始。

采样是统计测量，分辨率由频率和总时长决定。典型频率每秒 99 次：跑 10 秒的程序约采 990 个样本，一个占时 12\% 的函数期望分到约 119 个样本，计数本身的统计起伏就有 $\pm\sqrt{119}\approx\pm11$ 个，即占比估计自带约 $\pm1\%$ 的不确定。总共只跑 2 ms 的函数，平均摊不到四分之一个样本——采样找不到它，不等于它不慢，而是这个方法看不见它。

\begin{longtable}{L{0.14\linewidth}L{0.30\linewidth}L{0.20\linewidth}L{0.26\linewidth}}
\toprule
层次 & 回答的问题 & 精度 & 成本与局限 \\
\midrule
墙钟计时 & 用户等了多久 & 时钟分辨率高，读数含全部系统噪声 & 最简单；说不清时间花在哪里 \\\tablerowrule
性能计数器 & 硬件发生了多少事件 & 精确到单个事件 & 要先选对事件；同时可测事件数有限 \\\tablerowrule
采样分析 & 时间花在哪段代码 & 由采样频率与总时长决定 & 短函数、罕见路径会漏；不解释原因 \\
\bottomrule
\end{longtable}

三层的分工是从外到内：先用墙钟确认“确实慢了、慢了多少”，再用采样定位“慢在哪个函数”，最后用计数器回答“为什么慢”。顺序反过来就会走弯路：只报墙钟数字，改进无从下手；一上来就读计数器，可能在只占 2\% 时间的代码上数了半天事件。

三层其实相连：计数器计到指定次数可以溢出触发中断，采样器借此把“哪条指令在 miss”也记录下来——事件采样是计数器的精度加采样的位置信息。教学计算机项目中的教学机是天然的对照组：它没有性能计数器，测一段程序只有秒表和手工数指令两个手段，墙钟加“知道执行了多少条指令”就是全部测量学。计数器不是更高级的玩具，而是把处理器核心相关章节的 CPI 公式做成了可以直接读数的硬件。

墙钟计时还有一个实现细节：读哪只钟。日历时钟（墙上的时间）可能被校时跳变，测时间间隔要用单调时钟（单调递增、不受校时影响）；在教学计算机项目中的教学机上没有操作系统，直接数时钟周期反而更本质——墙钟的“钟”从来就是周期计数，只是被系统调用封装成了时间单位。

三层各有趁手的工具：墙钟一层是 \code{time} 命令和语言自带的计时函数；计数器一层是 \code{perf stat}、VTune 这类直接读硬件寄存器的工具；采样一层是 \code{perf record}、gprof。选工具先看要回答的问题属于哪一层，再看精度够不够用——拿 \code{time} 去定位热点，和拿采样器去报单次时延，是用错了层次。

\topic{消除测量误差}
测量的目标量是“这段代码在稳定状态下花多少时间”，而任何一次真实测量都混着三类与代码无关的成分。误差控制不是把数字测得更“准”，而是把这三类成分分别从读数里剥掉或固定住。

第一类是冷启动。第一次运行时，I-cache、D-cache、TLB、分支预测器全是空的，输入文件可能还躺在磁盘页缓存之外，CPU 主频可能还没升上来。某函数稳态 9.4 ms、首次冷跑 41 ms，4 倍多的差距一点不稀奇。冷启动本身是有意义的量——用户第一次使用就是冷的——但它和稳态是两个量，必须分开报告：要么先跑若干次预热再计时，要么明确给出“冷”“热”两组数字，混着报等于什么都没报。

第二类是随机噪声。调度、中断、后台进程让每次读数略有不同，读数序列常常双峰：大部分时候安静，偶尔受到一次显著干扰。看一组十次读数（ms）：98、99、101、97、142、98、99、100、98、99。均值 103.1，中位数 99，最小值 97——均值被 142 那一个离群点拖高了 4\%，而代码本身什么都没变。

把这十次读数画成直方图，均值如何被拖离群体就一目了然：

\begin{pdfdiagram}
  % 直方图：97:1, 98:3, 99:3, 100:1, 101:1,（断轴）142:1
  \draw[BookMuted, line width=0.6pt, ->] (0,0) -- (8.8,0);
  \draw[BookMuted, line width=0.6pt, ->] (0,0) -- (0,3.4);
  \node[hw label, anchor=north] at (4.2,-0.7) {单次测量读数（ms）};
  \node[hw label, rotate=90, anchor=south] at (-0.7,1.7) {次数};
  \foreach \y in {1, 2, 3} {
    \draw[BookMuted, line width=0.5pt] (0.06,{0.9*\y}) -- (-0.06,{0.9*\y});
    \node[hw label, anchor=east] at (-0.12,{0.9*\y}) {\y};
  }
  \draw[BookTeal, line width=0.8pt, fill=BookCodeBg] (0.4,0) rectangle (1.0,0.9);
  \draw[BookTeal, line width=0.8pt, fill=BookCodeBg] (1.2,0) rectangle (1.8,2.7);
  \draw[BookTeal, line width=0.8pt, fill=BookCodeBg] (2.0,0) rectangle (2.6,2.7);
  \draw[BookTeal, line width=0.8pt, fill=BookCodeBg] (2.8,0) rectangle (3.4,0.9);
  \draw[BookTeal, line width=0.8pt, fill=BookCodeBg] (3.6,0) rectangle (4.2,0.9);
  \node[hw label, anchor=north] at (0.7,-0.08) {97};
  \node[hw label, anchor=north] at (1.5,-0.08) {98};
  \node[hw label, anchor=north] at (2.3,-0.08) {99};
  \node[hw label, anchor=north] at (3.1,-0.08) {100};
  \node[hw label, anchor=north] at (3.9,-0.08) {101};
  % 断轴后的离群点
  \node[hw label] at (4.9,0.15) {//};
  \draw[BookRed, line width=0.8pt, fill=white] (5.6,0) rectangle (6.2,0.9);
  \node[hw label, anchor=north] at (5.9,-0.08) {142};
  % 中位数不动，均值被拖向离群点
  \draw[BookTeal, dashed, line width=0.6pt] (2.3,0) -- (2.3,3.1);
  \node[hw label, anchor=west] at (2.4,3.12) {中位数 99：不动};
  \draw[BookRed, dashed, line width=0.6pt] (4.55,0) -- (4.55,3.1);
  \node[hw label, anchor=west] at (4.65,2.72) {均值 103.1：被拖向离群点};
\end{pdfdiagram}
\diagramnote{十次读数的分布（横轴在 101 之后断轴）：九次聚在 97--101 ms，一次离群到 142 ms。中位数 99 不受离群点影响；均值 103.1 被它拖高 4\%，偏离了九次读数所在的群体。所以这组读数要报告中位数加离散程度；在“机器本身能跑多快”的问题上，最小值反而是受噪声污染最少的估计。}

对这种序列，取均值是最差的汇总方式：既被离群值拖着走，又掩盖双峰。正确做法是测几十次、报告分布：中位数加离散程度；在“这台机器本身能跑多快”的问题上，最小值反而是受噪声污染最少的估计——没有任何机制能让程序“虚假地变快”，却有无数机制让它变慢。次数的底线是中位数要稳定：从 30 次起步，加到读数不再明显移动为止。

第三类是混淆变量：两次比较之间改了不止一个东西。最典型的错误案例是“优化后快了 30\%”，复查发现输入从 $10^{6}$ 个元素换成了 $10^{5}$ 个——时延当然变了，和代码无关。固定做法是一句老生常谈但最容易被违反的话：一次实验只改一个因子。

还要注意系统误差，它不像随机噪声那样能靠多次测量摊掉：地址空间布局随机化让每次运行的 cache 冲突略有不同，链接顺序变化能让同一源码的实测差出几个百分点，主频调节策略（省电还是性能）更是全程一致的偏移。这类误差只能靠固定环境和明确声明来控制，靠平均消不掉。

一个具体例子：同一份源码、同一台机器，仅改变链接顺序，热循环与某个大数据结构的 cache 组映射关系就变了，实测差出 5\%——不是代码变快或变慢，是布局运气变了。对这类误差的诚实写法是报告分布并声明布局条件，而不是挑一个最好看的数字。

把三类误差倒过来看，就得到测量计划的固定写法：动手之前先写假设（“这个改动应把 cache miss 率从 56\% 降到 12\%”），再写要测的事件与次数，再对照清单固定环境，最后才运行程序。计划写在测量之前；“想到什么测什么”，本身就是误差最大的测量方法。

“这组数字能不能比”，是引用任何对比数字之前要先问的问题。下面这张清单建议逐条过：

\begin{longtable}{L{0.20\linewidth}L{0.38\linewidth}L{0.32\linewidth}}
\toprule
检查点 & 关键问题 & 不满足时的后果 \\
\midrule
同一输入 & 两次运行的数据规模和内容是否相同 & 时延随输入变化，比较失去意义 \\\tablerowrule
同一环境 & 机器、主频策略、后台负载是否一致 & 环境噪声可以大于被测差异 \\\tablerowrule
同一构建 & 编译器版本与优化选项是否一致 & 差异来自编译器而非代码改动 \\\tablerowrule
冷热状态 & 两次是否都预热，或都明确测冷启动 & 冷热混测可差出数倍 \\\tablerowrule
样本与分布 & 次数是否足够，报中位数还是均值 & 单次读数可能只是离群点 \\\tablerowrule
单位与定义 & 报的是时延、吞吐还是加速比 & 不同量纲的数字不能互比 \\
\bottomrule
\end{longtable}

工程含义：性能评审里的大量争论——“我的版本快 20\%”“不可能，我这里还慢了”——最后常常发现一边测的是冷启动、一边测的是稳态，或者一边开了 \code{-O2}、一边没开。把测量约定写在数字前面，数字才有意义；清单上任何一条答不上来，结论应先重新测量，不丢人。

\topic{性能计数器能回答什么}
计数器把处理器核心相关章节的公式变成可以直接读数的量。处理器核心相关章节给出 总时间 = 指令数 $\times$ CPI $\times$ 时钟周期，并把有效 CPI 拆成基础值加访存、分支等各项代价；\code{cycles} 和 \code{instructions} 两个事件正好对应公式的前两个因子，其余事件对应各项代价的分子与分母。

最常用的四个比值：IPC（instructions per cycle，\code{instructions}/\code{cycles}，即 CPI 的倒数；四发射机器的理论上限是 4，实测 0.83 意味着近八成发射槽在空转）；cache miss 率（\code{cache-misses}/\code{cache-references}，存储与设备事务相关章节的命中率就是 1 减去它）；分支误预测率（\code{branch-misses}/\code{branch-instructions}，每次误测都要按处理器核心相关章节说的清空错误路径，付出十几拍）；TLB miss 率（\code{dTLB-load-misses}/\code{dTLB-loads}，每次失效意味着额外的页表遍历访存）。

Linux 下一条 \code{perf stat} 就能报出这组数字，输出形态如下（数字为示意）：

\begin{lstlisting}
$ perf stat ./prog
       6,291,456  cycles
       5,242,880  instructions          # 0.83 insn per cycle
       2,097,152  cache-references
          41,943  cache-misses          # 2.0% of all cache refs
       1,048,576  branch-instructions
          31,457  branch-misses         # 3.0% of all branches
\end{lstlisting}

读这组数字就是按处理器核心相关章节的有效 CPI 公式做分解。cache 一项：每条指令平均 0.4 次数据访存，miss 率 2\%，miss 代价 50 拍，贡献 $0.4\times0.02\times50=0.40$ 拍/指令。分支一项：分支占指令两成，误预测率 3\%，每次罚 15 拍，贡献 $0.2\times0.03\times15=0.09$ 拍/指令。基础 CPI 若为 1.0，有效 CPI 就是约 1.49——分解做完，该查哪一侧已经写在算术里。

把 \code{perf stat} 的输出逐行挂上解读，六个数字与有效 CPI 公式的对应关系就固定下来了：

\begin{pdfdiagram}[x=0.85cm]
  % 左：perf stat 输出（示意数字）
  \node[anchor=west, font=\scriptsize\ttfamily, text=BookInk] (p1) at (-7.0,3.0) {6,291,456  cycles};
  \node[anchor=west, font=\scriptsize\ttfamily, text=BookInk] (p2) at (-7.0,2.4) {5,242,880  instructions};
  \node[anchor=west, font=\scriptsize\ttfamily, text=BookInk] (p3) at (-7.0,1.3) {2,097,152  cache-references};
  \node[anchor=west, font=\scriptsize\ttfamily, text=BookInk] (p4) at (-7.0,0.7) {   41,943  cache-misses};
  \node[anchor=west, font=\scriptsize\ttfamily, text=BookInk] (p5) at (-7.0,-0.4) {1,048,576  branch-instructions};
  \node[anchor=west, font=\scriptsize\ttfamily, text=BookInk] (p6) at (-7.0,-1.0) {   31,457  branch-misses};
  % 右：解读框
  \node[hw compute, minimum width=5.2cm, minimum height=1.05cm, align=center] (a1) at (3.6,2.7) {公式前两个因子\\IPC = 0.83，即 CPI $\approx$ 1.20};
  \node[hw compute, minimum width=5.2cm, minimum height=1.05cm, align=center] (a2) at (3.6,1.0) {miss 率 41,943/2,097,152 = 2.0\%\\存储项 0.4$\times$0.02$\times$50 = 0.40 拍/指令};
  \node[hw compute, minimum width=5.2cm, minimum height=1.05cm, align=center] (a3) at (3.6,-0.7) {误测率 31,457/1,048,576 = 3.0\%\\分支项 0.2$\times$0.03$\times$15 = 0.09 拍/指令};
  % 逐组虚线对应
  \draw[hw return, dashed] (p1.east) -- (-0.9,3.0) -- (-0.9,2.85) -- ($(a1.west)+(0,0.15)$);
  \draw[hw return, dashed] (p2.east) -- (-0.5,2.4) -- (-0.5,2.55) -- ($(a1.west)+(0,-0.15)$);
  \draw[hw return, dashed] (p3.east) -- (-0.9,1.3) -- (-0.9,1.15) -- ($(a2.west)+(0,0.15)$);
  \draw[hw return, dashed] (p4.east) -- (-0.5,0.7) -- (-0.5,0.85) -- ($(a2.west)+(0,-0.15)$);
  \draw[hw return, dashed] (p5.east) -- (-0.9,-0.4) -- (-0.9,-0.55) -- ($(a3.west)+(0,0.15)$);
  \draw[hw return, dashed] (p6.east) -- (-0.5,-1.0) -- (-0.5,-0.85) -- ($(a3.west)+(0,-0.15)$);
  \node[hw label] at (0.5,-1.95) {有效 CPI $\approx$ 1.0 + 0.40 + 0.09 = 1.49：先查存储侧，再查分支};
\end{pdfdiagram}
\diagramnote{\code{perf stat} 输出的读法（数字为示意）：\code{cycles} 与 \code{instructions} 相除得 IPC 0.83，即 CPI 约 1.20；\code{cache-misses} 对 \code{cache-references} 得 miss 率 2.0\%，按每条指令 0.4 次访存、miss 代价 50 拍折算，存储项 0.40 拍/指令；\code{branch-misses} 对 \code{branch-instructions} 得误预测率 3.0\%，按分支占两成、每次误测罚 15 拍折算，分支项 0.09 拍/指令。基础 CPI 1.0 加上这两项，有效 CPI 约 1.49——先查存储侧，再查分支。}

IPC 与 CPI 互为倒数，只是同一件事的两种说法：\code{perf} 习惯报 IPC，处理器核心相关章节的公式用 CPI，换算只要取倒数——IPC 0.83 就是 CPI $1/0.83\approx1.20$。报告里用哪个都可以，但一份报告内部要统一，免得读者把 0.83 和 1.20 当成两台机器的数字。

使用计数器还有两个实务提醒。一是读硬件计数器需要权限，容器和虚拟机里事件可能被屏蔽或换成虚拟值，正式测量前先用一个行为已知的小程序标定：比如空转的时钟周期数应约等于时长乘主频，对不上就说明读数通道有问题。二是事件定义因微架构而异，Intel、AMD、ARM 上同名事件的统计边界可能不同，报告里要写清机器型号，否则别人无法复核。

\begin{longtable}{L{0.26\linewidth}L{0.30\linewidth}L{0.34\linewidth}}
\toprule
perf 事件名 & 含义 & 典型比值与判断 \\
\midrule
\code{cycles}、\code{instructions} & 总拍数与退休指令数 & IPC 远低于发射宽度，停顿多，继续查 miss 与误测 \\\tablerowrule
\code{cache-references}、\code{cache-misses} & 末级 cache 访问与失效 & miss 率超过几个百分点，访存是瓶颈候选 \\\tablerowrule
\code{branch-instructions}、\code{branch-misses} & 分支总数与误测次数 & 分支占两成、罚 15 拍时，误预测率每 1\% 抬高 CPI 约 0.03 \\\tablerowrule
\code{dTLB-loads}、\code{dTLB-load-misses} & 数据 TLB 访问与失效 & 失效偏高，查大页与跨页访问模式 \\\tablerowrule
\code{L1-dcache-load-misses} & L1 数据 cache 读失效 & 与末级 miss 对照，定位失效发生在哪一级 \\\tablerowrule
\code{L1-icache-load-misses} & 指令 cache 失效 & 热代码 footprint 超过 L1i 的信号（见案例三） \\
\bottomrule
\end{longtable}

计数器的正确读法是先假设、后验证，而不是把事件全测一遍看哪个大。完整流程是：先根据现象提出一个可证伪的假设（“这段代码慢是因为 cache miss 多”）；再选一组能区分假设成立与否的事件（\code{cache-references} 与 \code{cache-misses}）；最后看读数支持还是推翻。假设被推翻同样是收获：实测 miss 率只有 0.3\%，“cache 瓶颈”这条线就此划掉，转查分支误测或前端取指。没有假设的直接读数，面对几十个事件名只会越看越糊涂——计数器回答“是不是”和“有多少”，假设才决定该问哪个问题。

四个比值还天然排成一个排查顺序：先看 IPC 低不低，确认确实在停顿；再看停顿来自哪一侧——cache miss 率高查数据局部性，分支误预测率高查代码形态，TLB miss 率高查页面布局，I-cache miss 高查代码体积。下面三个案例各按这个顺序走一遍。

\topic{案例一：矩阵转置的 cache 行为}
任务：把 $1024\times1024$ 的双精度矩阵 A 转置写入 B。每个元素 8 字节，整矩阵 8 MiB。机器参数沿用存储与设备事务相关章节的设定：cache line 64 B（恰好装 8 个 double），L1 数据 cache 32 KiB（512 条 line），L2 256 KiB，写策略 write-allocate。两个朴素版本只差循环顺序，第三个版本做 $8\times8$ 分块：

\begin{lstlisting}
/* 版本甲：i 外层 j 内层。A 按行读（连续），B 按列写（跳） */
for (i = 0; i < 1024; i++)
    for (j = 0; j < 1024; j++)
        B[j][i] = A[i][j];

/* 版本乙：j 外层 i 内层。A 按列读（跳），B 按行写（连续） */
for (j = 0; j < 1024; j++)
    for (i = 0; i < 1024; i++)
        B[j][i] = A[i][j];

/* 版本丙：8x8 分块，块内完成小转置 */
for (ii = 0; ii < 1024; ii += 8)
    for (jj = 0; jj < 1024; jj += 8)
        for (i = ii; i < ii + 8; i++)
            for (j = jj; j < jj + 8; j++)
                B[j][i] = A[i][j];
\end{lstlisting}

先按存储与设备事务相关章节的局部性原理算理论值。版本甲：A 按行连续读，一条 line 装 8 个元素，前 7 次访问命中同一条 line，读侧 miss 率 $1/8=12.5\%$，全部是不得不付的强制失效；B 按列写，相邻两次写地址相差一整行 8 KiB，每次写都访问一条不同的 line，一列 1024 个元素没写完，开头的 line 早被挤出 32 KiB 的 L1，写侧 miss 率接近 100\%。版本乙把好坏对调：B 连续写 12.5\%，A 跳跃读接近 100\%。版本丙的每个 $8\times8$ 块内，A 的 8 行各贡献一条被用满的 line，写入 B 的 8 条 line 同样被用满，两侧都回到 12.5\% 的强制失效下限。

把连续与跳跃两个方向的内存足迹画出来，12.5\% 对 100\% 就是“一条 line 服务几次访问”的差别：

\begin{pdfdiagram}[x=0.9cm]
  % 左：连续方向（按行）
  \node[hw label] at (2.3,4.7) {连续方向（按行）};
  \draw[BookTeal, line width=0.8pt, fill=BookCodeBg] (0.6,3.1) rectangle (4.0,3.9);
  \node[hw label] at (2.3,3.5) {一条 cache line（8 个 double）};
  \foreach \k in {1,...,8} {
    \fill[BookAccent] ({0.6+0.425*\k-0.2125},2.5) circle (0.05);
    \draw[hw bus] ({0.6+0.425*\k-0.2125},2.6) -- ({0.6+0.425*\k-0.2125},3.07);
  }
  \node[hw label, anchor=east] at (0.45,2.5) {访问 1..8};
  \node[hw label] at (2.3,1.85) {一条 line 服务 8 次：miss 率 1/8 = 12.5\%};
  % 右：跳跃方向（按列）
  \node[hw label] at (8.6,4.7) {跳跃方向（按列）};
  \draw[BookTeal, line width=0.8pt, fill=BookCodeBg] (7.6,3.6) rectangle (9.6,4.2);
  \node[hw label, anchor=east] at (7.45,3.9) {行 i};
  \fill[BookRed] (8.6,3.9) circle (0.055);
  \draw[BookTeal, line width=0.8pt, fill=BookCodeBg] (7.6,2.5) rectangle (9.6,3.1);
  \node[hw label, anchor=east] at (7.45,2.8) {行 i+1};
  \fill[BookRed] (8.6,2.8) circle (0.055);
  \draw[BookTeal, line width=0.8pt, fill=BookCodeBg] (7.6,1.4) rectangle (9.6,2.0);
  \node[hw label, anchor=east] at (7.45,1.7) {行 i+2};
  \fill[BookRed] (8.6,1.7) circle (0.055);
  \draw[hw bus] (8.6,3.58) -- (8.6,3.12);
  \draw[hw bus] (8.6,2.48) -- (8.6,2.02);
  \node[hw label, anchor=west] at (9.75,2.8) {步长 8 KiB};
  \node[hw label] at (8.6,0.95) {每条 line 只服务 1 次：miss $\approx$ 100\%};
  % 底部注记
  \node[hw label, anchor=north] at (5.0,0.35) {步长 8 KiB 恰为 2 的幂：组映射下集中撞少数几个组，冲突失效把 miss 率顶在高位};
\end{pdfdiagram}
\diagramnote{矩阵转置两个方向的内存足迹。左：按行连续访问，一条 64 B 的 line 装 8 个 double，第 1 次访问 miss、后 7 次命中，miss 率就是强制失效下限 12.5\%。右：按列跳跃访问，相邻两次地址相差一整行 8 KiB，每次都访问一条不同的 line，取回来只用一次就被挤出，miss 率接近 100\%；步长 8 KiB 恰为 2 的幂，组映射下还集中撞少数几个组，冲突失效雪上加霜。}

理论 miss 数可以直接算：读侧、写侧各 $1{,}048{,}576$ 次访问；连续方向 miss $1{,}048{,}576/8=131{,}072$ 次，跳跃方向约 $1{,}048{,}576$ 次。预热 3 次、各测 30 次取中位数，实测与理论吻合（教学测量，示意数据）：

\begin{longtable}{L{0.16\linewidth}L{0.18\linewidth}L{0.18\linewidth}L{0.18\linewidth}L{0.16\linewidth}}
\toprule
版本 & A 侧 miss 率 & B 侧 miss 率 & 总 miss 数 & 墙钟时间 \\
\midrule
甲（A 顺 B 跳） & 12.5\% & $\approx$100\% & $1{,}179{,}648$ & 28.3 ms \\\tablerowrule
乙（A 跳 B 顺） & $\approx$100\% & 12.5\% & $1{,}179{,}648$ & 31.1 ms \\\tablerowrule
丙（$8\times8$ 分块） & 12.5\% & 12.5\% & $262{,}144$ & 9.4 ms \\
\bottomrule
\end{longtable}

第一点值得展开的是那 8 倍：miss 率 12.5\% 对 100\% 不是“快一点”的程度差别，而是一条 line 里装了几个元素的结构性差别。空间局部性按 line 计、不按字节计：顺着 line 走，一条 line 服务 8 次访问；跨着 line 跳，一条 line 只服务 1 次。存储与设备事务相关章节讲的“取回一整行”，在这里变成了可以精确预测的 miss 率。

第二点是甲、乙总 miss 数完全相同（都是 $131{,}072+1{,}048{,}576$），时间却是 28.3 ms 对 31.1 ms。这不是噪声：读 miss 直接挡住等数据的指令，写 miss 可以被写缓冲吸收一部分，同一个 miss 数，读侧的代价更大。调换循环顺序不是免费的优化，要算清哪一侧更怕 miss。

第三点是步长 8 KiB 恰好是 2 的幂：按列访问在组映射下集中命中少数几个组，即便容量还剩，冲突失效也把 miss 率顶在接近 100\%。这也是高性能矩阵代码常给行加 padding、避开 2 幂行距的原因。

TLB 也在同一场景里承压：矩阵占 $8\ \mathrm{MiB}/4\ \mathrm{KiB}=2048$ 个页，按列走时连续访问以 8 KiB 为步长跨页，一列没走完 dTLB 已经换了好几轮，\code{dTLB-load-misses} 随 cache miss 一起上涨。计数器之间会互相印证——这正是“先假设后验证”里假设越来越具体的形态。

跳跃写的代价里还藏着 write-allocate 的一份：写 miss 要先把整行取进 cache 再改，B 侧的实际内存流量是“取一遍加写回一遍”的两倍。版本乙把写侧变连续后，写合并和写缓冲才发挥得了作用。转置这个例子里，存储与设备事务相关章节的每一条 cache 规则都至少出场一次。

硬件预取器在同一场景里的表现也值得一提：对连续流（版本甲的 A 侧、版本乙的 B 侧），预取器能把一部分强制失效的时延提前隐藏，实测时间略好于纯 miss 数推算；对 8 KiB 步长的跳跃流，多数预取器直接放弃。这解释了为什么跳跃方向的代价往往比“miss 数乘 miss 时延”的估算更硬——它连预取这根拐杖都用不上。

回看这个案例该测哪些事件，清单几乎是自己写出来的：\code{L1-dcache-load-misses} 验证读侧，\code{cache-misses}（含写失效）看总量，\code{dTLB-load-misses} 看跨页，再加 \code{cycles}、\code{instructions} 换算 CPI——五条事件足够判定“局部性假设”成立与否。事件清单之所以短，是因为假设先写清楚了。

朴素转置的读写两侧方向天然相反，单层循环顺序怎么换都只能照顾一边；分块改的是数据粒度而不是访问次数。块多大合适？把块尺寸扫一遍就清楚了（其余条件同前）：

\begin{longtable}{L{0.16\linewidth}L{0.24\linewidth}L{0.20\linewidth}L{0.16\linewidth}}
\toprule
块尺寸 & B 侧复用窗口 & 总 miss 数 & 墙钟时间 \\
\midrule
$8\times8$ & 8 条 B line 加 1 条 A line & $262{,}144$ & 9.4 ms \\\tablerowrule
$64\times64$ & 64 条加 8 条 & $262{,}144$ & 9.3 ms \\\tablerowrule
$256\times256$ & 256 条加 32 条 & $262{,}144$ & 9.4 ms \\\tablerowrule
$512\times512$ & 512 条加 64 条，超过 L1 & $\approx4.4\times10^{5}$ & 13.0 ms \\\tablerowrule
$1024\times1024$（不分块） & 1024 条加 128 条 & $1{,}179{,}648$ & 28.3 ms \\
\bottomrule
\end{longtable}

窗口的意思是：对 B 的一条 line 写完 8 个元素，要横跨一整轮块内循环，期间触碰的 line 数就是“块高条 B line 加块宽除以 8 条 A line”。窗口小于 512 条 L1 容量时，每个块都贴着强制失效下限 $2\times131{,}072=262{,}144$；窗口超过容量，line 在复用之前被挤出，miss 率回升；块大到等于不分块，就退回版本甲。分块参数不是调出来的魔法数，是用 cache 容量算出来的。

这正是第四节 roofline 中“cache 友好版”矩阵算法的同一方法：不减少访问次数，改变访问顺序让每条 line 被用满。也请注意时间与 miss 数并非严格正比（miss 数 4.5 倍，时间 3.0 倍）——乱序执行和写缓冲会掩盖一部分停顿，精确的时间模型留给机器实测，计数器负责把方向指对。

\topic{案例二：循环展开换来的 CPI}
任务：对 $1{,}048{,}576$ 个 double 求和。机器 2.0 GHz、四发射，加法延迟 1 拍，L1 命中的 load 可流水供应。基础版循环体每元素 5 条指令：load、add、下标加一、比较、条件跳转——前两条是“正事”，后三条是循环开销。

先算指令数因子。基础版：$5\times1{,}048{,}576=5{,}242{,}880$ 条。展开 4 次后每 4 个元素 11 条指令（4 load、4 add，下标、比较、跳转各 1），共 $262{,}144\times11=2{,}883{,}584$ 条，省了 45.0\%——省的正是被摊薄的循环开销。分支数同步从 $1{,}048{,}576$ 降到 $262{,}144$。

再算 CPI 因子，这一步的关键是用独立累加器 \code{sum0} 到 \code{sum3} 分别累加、出循环再相加。原因在依赖链：基础版的 \code{sum} 每次 add 都等上一次 add，链长等于元素个数，四发射机器被这条链钉在 CPI 约 1.2（每元素 5 条指令，约 6 拍）；四条独立累加链把这条长链拆成四段短链，加法不再等前一次的结果，CPI 降到 0.70——仍高于四发射的理想值，因为 load 与循环开销指令同样要占发射槽，依赖链只是约束之一。

\begin{lstlisting}
sum0 = sum1 = sum2 = sum3 = 0;
for (i = 0; i < N; i += 4) {
    sum0 += a[i];      /* 四条累加链互相独立 */
    sum1 += a[i+1];
    sum2 += a[i+2];
    sum3 += a[i+3];
}
sum = sum0 + sum1 + sum2 + sum3;
\end{lstlisting}

把展开倍数扫一遍，两个因子的此消彼长就全在表里（同机实测，预热后取 30 次中位数）：

\begin{longtable}{L{0.16\linewidth}L{0.20\linewidth}L{0.16\linewidth}L{0.12\linewidth}L{0.14\linewidth}}
\toprule
展开倍数 & 指令数 & 分支数 & CPI & 墙钟时间 \\
\midrule
1 & $5{,}242{,}880$ & $1{,}048{,}576$ & 1.20 & 3.15 ms \\\tablerowrule
2 & $3{,}670{,}016$ & $524{,}288$ & 0.95 & 1.74 ms \\\tablerowrule
4 & $2{,}883{,}584$ & $262{,}144$ & 0.70 & 1.01 ms \\\tablerowrule
8 & $2{,}490{,}368$ & $131{,}072$ & 0.75 & 0.93 ms \\\tablerowrule
16 & $\approx2{,}523{,}000$ & $65{,}536$ & 0.90 & 1.14 ms \\
\bottomrule
\end{longtable}

演算核对：基础版 $5{,}242{,}880\times1.20=6{,}291{,}456$ 拍，除以 2.0 GHz 得 3.15 ms；展开 4 次 $2{,}883{,}584\times0.70\approx2{,}018{,}509$ 拍得 1.01 ms。对基础版的总加速约 3.1 倍，其中指令数贡献 1.82 倍（$5{,}242{,}880/2{,}883{,}584$）、CPI 贡献 1.71 倍（$1.20/0.70$），两个因子相乘、方向一致——这不多见，案例三会给出反例。

分支一项值得说清：循环分支几乎 100\% 可预测，展开前后误预测率都低到可以忽略，收益不在误预测率，而在每元素少付 2.25 条开销指令和依赖链变短。看优化收益时，先分清它作用在公式的哪个因子上。

展开 16 次开始倒贴：理论上每 16 个元素 35 条指令、共 $65{,}536\times35=2{,}293{,}760$ 条，但 16 个累加器加 16 个 load 目标超出 16 个通用寄存器，编译器把中间值溢写到栈上，实测指令数回升约一成，且溢写引入新依赖，CPI 回到 0.90，时间反弹到 1.14 ms。8 次的 0.93 ms 是这台机器上的甜点：8 个累加器加 8 个 load 目标，恰好顶到寄存器堆容量的边缘。

把指令数与墙钟时间都归一到基础版再画成曲线，8 次这个甜点就成了两条线各自的形状：

\begin{pdfdiagram}
  % 横轴展开倍数（log2，每档 1.5）；纵轴相对基础版（=1），y = 相对值 * 3.2
  \draw[BookMuted, line width=0.6pt, ->] (0,0) -- (8.0,0);
  \draw[BookMuted, line width=0.6pt, ->] (0,0) -- (0,4.4);
  \node[hw label, anchor=north] at (3.6,-0.62) {展开倍数};
  \node[hw label, rotate=90, anchor=south] at (-0.95,2.1) {相对基础版（= 1）};
  \foreach \x/\t in {0/1, 1.5/2, 3/4, 4.5/8, 6/16} {
    \draw[BookMuted, line width=0.5pt] (\x,0.06) -- (\x,-0.06);
    \node[hw label, anchor=north] at (\x,-0.1) {\t};
  }
  \foreach \y/\t in {0.8/0.25, 1.6/0.5, 2.4/0.75, 3.2/1.0} {
    \draw[BookMuted, line width=0.5pt] (0.06,\y) -- (-0.06,\y);
    \node[hw label, anchor=east] at (-0.12,\y) {\t};
  }
  % 指令数：1, 0.70, 0.55, 0.475, 0.481
  \draw[BookAccent, line width=1pt] plot[smooth] coordinates {
    (0,3.2) (1.5,2.24) (3,1.76) (4.5,1.52) (6,1.54)
  };
  % 墙钟时间：1, 0.552, 0.321, 0.295, 0.362
  \draw[BookRed, line width=1pt] plot[smooth] coordinates {
    (0,3.2) (1.5,1.77) (3,1.03) (4.5,0.94) (6,1.16)
  };
  \fill[BookInk] (0,3.2) circle (1.6pt);
  \node[hw label, anchor=west] at (0.15,3.45) {基础版：5.24M 条 · 3.15 ms};
  % 甜点：8 次
  \fill[BookRed] (4.5,0.94) circle (1.6pt);
  \draw[BookTeal, dashed, line width=0.5pt] (4.5,0) -- (4.5,0.94);
  \node[hw label, anchor=west] at (4.65,0.7) {甜点：8 次，0.93 ms};
  % 16 次回升
  \fill[BookRed] (6,1.16) circle (1.6pt);
  \node[hw label, anchor=west] at (6.2,1.5) {16 次回升：寄存器溢写增加访存开销};
  % 图例
  \draw[BookAccent, line width=1pt] (0.4,4.15) -- (1.1,4.15);
  \node[hw label, anchor=west] at (1.25,4.15) {指令数};
  \draw[BookRed, line width=1pt] (2.9,4.15) -- (3.6,4.15);
  \node[hw label, anchor=west] at (3.75,4.15) {墙钟时间};
\end{pdfdiagram}
\diagramnote{循环展开 1、2、4、8、16 倍的两条归一化曲线（基础版 = 1，对应 5.24M 条指令、3.15 ms）。指令数随展开单调下降——循环开销被摊薄；墙钟时间先降后升，8 次的 0.93 ms 是甜点：8 个累加器加 8 个 load 目标恰好顶到寄存器堆容量。展开 16 次后寄存器不够用，中间值溢写到栈上，指令数与 CPI 双双回升。}

收益边界由两个物理量固定：寄存器数量和代码体积。累加器与 load 目标各占一个寄存器，展开倍数越过容量就开始因 load/store 溢写而增加访存开销；代码体积按倍数乘循环体大小增长，涨过 I-cache 容量就是案例三的剧情。还有一点：结论不可随输入类型外推，整数加法与双精度加法的延迟不同、链的限制不同，换数据类型要重测——这正是控制变量那一节的要求。

两个相关问题顺手说清。其一，现代编译器在 \code{-O2} 下本来就会做适度展开和调度，手写展开前要先看编译器已经做到了哪一步——测“编译器自动展开”与“手写四累加器”的差，才是手写这一笔的真实收益，否则容易把编译器的功劳算在自己名下。其二，展开常与软件流水混淆：展开是复制循环体、摊薄开销，软件流水是把不同迭代的阶段错开重叠；两者都打破依赖，但流水不改变指令总数，展开改变。

展开还是向量化的前奏：四条独立累加链的下一步，往往是一条指令同时算四路的 SIMD。从这个意义上说，先把依赖链打断、再谈更宽的执行单元，顺序不能反——链不断，多宽的部件都在等上一次的 \code{sum}。

\topic{案例三：一次错误的“优化”}
某协议栈里有一个 24 字节的小循环——逐字节查表更新 CRC——调用非常频繁，采样分析也确认它是热点。优化者决定把它内联到全部约 40 个调用点，顺手把循环完全展开。推导无懈可击：省掉每次调用的 call/ret、参数搬运和栈帧调整，动态指令数确实降了 8\%，从 $1.00\times10^{8}$ 条降到 $9.2\times10^{7}$ 条。

\begin{longtable}{L{0.22\linewidth}L{0.16\linewidth}L{0.16\linewidth}L{0.12\linewidth}L{0.14\linewidth}}
\toprule
版本 & 指令数 & L1i miss & CPI & 墙钟时间 \\
\midrule
内联前 & $1.00\times10^{8}$ & $0.2\times10^{6}$ & 1.10 & 55.0 ms \\\tablerowrule
内联加全展开 & $0.92\times10^{8}$ & $3.1\times10^{6}$ & 1.50 & 69.0 ms \\
\bottomrule
\end{longtable}

慢了 25.5\%。演算核对：内联前 $1.00\times10^{8}\times1.10=1.10\times10^{8}$ 拍，55.0 ms；内联后 $0.92\times10^{8}\times1.50=1.38\times10^{8}$ 拍，69.0 ms。指令数降了 8\%，CPI 却从 1.10 涨到 1.50，总账由盈转亏。

CPI 为什么涨：内联加完全展开后，主循环的指令 footprint——含全部约 40 个调用点展开后的代码总量，不是那个 24 字节的小循环本身——从 24 KiB 涨到 40 KiB，超过 32 KiB 的 L1 I-cache，每一轮主循环都要从 L2 重新取一部分指令。\code{L1-icache-load-misses} 从 $0.2\times10^{6}$ 涨到 $3.1\times10^{6}$，每次 miss 从 L2 取回约付 10 拍，新增 $2.9\times10^{6}$ 次就是约 $2.9\times10^{7}$ 拍——比省下的 $8\times10^{6}$ 条指令折合的全部收益还多。存储与设备事务相关章节讲 D-cache 的那套局部性逻辑对 I-cache 一字不差地成立：取指也是访存，代码也有 footprint。

\begin{lstlisting}
$ perf stat -e cycles,instructions,L1-icache-load-misses ./before
    110,034,512  cycles    100,012,440  instructions      198,744 L1i
$ perf stat -e cycles,instructions,L1-icache-load-misses ./after
    138,120,977  cycles     92,011,305  instructions    3,104,882 L1i
\end{lstlisting}

教训有两条。其一，“优化”是个测量命题，不是逻辑命题：推导再漂亮也只是假设，必须按本节的流程回测——先测基线，只改一个因子，再测，对照核对清单逐条过。其二，假设要选对因子：如果预先的假设是“省指令数就是省时间”，\code{instructions} 下降 8\% 会漂亮地“确认”假设；只有同时读 \code{cycles} 和前端取指事件，才会发现假设漏掉了 I-cache 这个变量。指令数和 CPI 是两个独立的因子，只算其中一个的优化，连方向都不能保证。

同类事故还有别的形态：头文件里的大宏在每个调用点展开、模板实例化让同一份逻辑复制几十份、过度循环展开——机制相同，都是热路径的代码 footprint 悄悄涨过 I-cache。正确做法是给代码体积设预算：改动前后各测一次 \code{L1-icache-load-misses}，热路径 footprint 不许越过 I-cache 容量这条线；真要内联，只内联最热的一两个调用点，而不是全部四十个。

把这次失败按本节的流程复盘一遍，每一步都对应得上：采样确认热点（第三层）→ 提出假设“内联省指令就是省时间” → 只改代码结构这一个因子（控制变量做对了）→ 回测墙钟发现变慢（第一层）→ 用计数器定位到 L1i（第二层）→ 假设被推翻，撤回改动。流程本身没有失败，失败的是没走流程就合入——这正是“先假设后验证”最值钱的场景。

工程上的默认动作因此只有一条：任何以“快”为目标的改动，合入前必须附回测数据，且按报告模板的五段写全。没有回测数据的“优化”一律按未优化处理——这条规矩看起来慢，实际省掉的是案例三这种来回。

\topic{性能报告模板}
测量做到最后，要变成别人能复核的文字。一份最低限度可信的性能报告固定五段：环境、方法、数据、分析、结论。

环境段写机器与主频、cache 层次参数、编译器与选项、输入规模——换一台机器时，读者据此判断哪些数字会跟着变。方法段写测了什么事件、冷热如何处理、测了多少次、报的是中位数还是均值——这一段决定数字能不能和别人的比。数据段放原始测量表，不加评论。分析段做比值演算，明说每个假设被支持还是被推翻。结论段只写可执行的判断，不写感想。

把五段画成结构图，每一段回答什么问题、不回答什么问题，就都标在它自己的格子里：

\begin{pdfdiagram}[x=0.9cm]
  \node[hw state, minimum width=2.6cm, minimum height=0.9cm] (s1) at (1.5,4.4) {环境};
  \node[hw state, minimum width=2.6cm, minimum height=0.9cm] (s2) at (1.5,3.3) {方法};
  \node[hw state, minimum width=2.6cm, minimum height=0.9cm] (s3) at (1.5,2.2) {数据};
  \node[hw state, minimum width=2.6cm, minimum height=0.9cm] (s4) at (1.5,1.1) {分析};
  \node[hw state, minimum width=2.6cm, minimum height=0.9cm] (s5) at (1.5,0.0) {结论};
  \draw[hw bus] (s1.south) -- (s2.north);
  \draw[hw bus] (s2.south) -- (s3.north);
  \draw[hw bus] (s3.south) -- (s4.north);
  \draw[hw bus] (s4.south) -- (s5.north);
  \node[hw label, anchor=west] at (3.3,4.4) {机器、cache 参数、编译选项、输入规模：换机器时判断哪些数字会变};
  \node[hw label, anchor=west] at (3.3,3.3) {测什么事件、冷热处理、次数、中位数还是均值：决定数字能不能互比};
  \node[hw label, anchor=west] at (3.3,2.2) {原始测量表，不加评论：报告中唯一不随立场改变的部分};
  \node[hw label, anchor=west] at (3.3,1.1) {比值演算：明说每个假设被支持还是被推翻};
  \node[hw label, anchor=west] at (3.3,0.0) {可执行判断：条件 + 版本 + 预期收益，不写感想};
  \node[hw label, anchor=north] at (5.2,-0.7) {五段一段不省，后来者才能不重测就复核结论在什么条件下成立};
\end{pdfdiagram}
\diagramnote{性能报告的五段结构：环境段让读者判断数字在什么条件下成立，方法段决定数字能不能与别人的互比，数据段只放原始测量、不加评论，分析段做比值演算并明说每个假设被支持还是被推翻，结论段只写可执行的判断——条件、版本、预期收益。五段一段不省，后来者才能不重测就复核；缺了任何一段，结论都只能算作者的私见。}

各段都有典型错误：环境段漏写主频调节策略；方法段只写“测了几次取平均”；数据段只给加速比不给原始值；分析段把假设成立当成理所当然，不给演算；结论段写“明显更快”这种无法执行的句子。方法段的合格写法可以直接照抄这个句式：“预热 3 次后连续测量 30 次，表中数值为中位数；两次运行之间仅修改循环顺序，编译选项固定为 \code{-O2}，机器无其他负载。”

数据段不加评论是有道理的：原始数字是报告中唯一不随作者立场改变的部分，评论写进分析段，读者才能把“你测到了什么”和“你怎么解释”分开复核。两者混在一起的报告，复核者要先把它们剥开才能检查。

结论段的合格句式是“在什么条件下、采用什么版本、预期多少收益”，例如：“在 32 KiB L1d、line 64 B 的机器上，矩阵转置采用 $8\times8$ 分块版，预期比朴素版快约 3 倍。”条件写在前面，是因为条件一变结论就可能失效——这正是环境段存在的意义。

按案例一填写的一份完整样例：

\begin{longtable}{L{0.10\linewidth}L{0.28\linewidth}L{0.50\linewidth}}
\toprule
段落 & 写什么 & 样例（按案例一填写） \\
\midrule
环境 & 机器、cache、编译选项、输入 & 2.0 GHz 四发射；L1d 32 KiB、line 64 B、write-allocate；\code{-O2}；$1024\times1024$ double \\\tablerowrule
方法 & 测什么、冷热、次数、统计量 & \code{perf stat} 测 \code{cache-misses} 等；预热 3 次后测 30 次，报中位数 \\\tablerowrule
数据 & 原始测量表 & 三版总 miss 数 $1{,}179{,}648$ / $1{,}179{,}648$ / $262{,}144$；时间 28.3 / 31.1 / 9.4 ms \\\tablerowrule
分析 & 比值演算与假设检验 & 假设“跳跃方向 miss 率接近 100\%”，实测与 $1/8$、$8/8$ 的理论值吻合，假设成立；分块回到强制失效下限 \\\tablerowrule
结论 & 可执行判断 & 采用 $8\times8$ 分块版；单层调换循环顺序只改好坏归属，不改总 miss 数 \\
\bottomrule
\end{longtable}

报告还有一个跨版本的作用：同一台机器、同一份方法，性能数字才能按版本追踪，回归才有人认。模板五段一段不省，后来者才能在不重测的情况下，判断你的结论在什么条件下成立。




























\section{七、结构、数据与控制冒险}

冒险表示下一拍若照常推进，就会使用不存在、过期或错误路径上的信息。

\begin{longtable}{L{0.18\linewidth}L{0.36\linewidth}L{0.36\linewidth}}
\toprule
类型 & 产生原因 & 典型处理 \\
\midrule
结构冒险 & 两个阶段同时争同一硬件端口 & 增加端口、分离资源或停顿 \\\tablerowrule
数据冒险 & 消费者早于生产者提交结果 & 转发、互锁停顿、编译器调度 \\\tablerowrule
控制冒险 & 下一 PC 尚未确定 & 预测执行；错误时冲刷并重定向 \\
\bottomrule
\end{longtable}

考虑相邻两条指令：

\begin{lstlisting}
ADD R3, R1, R2
SUB R5, R3, R4
\end{lstlisting}

ADD 的结果在 EX 末端已经算出，但要到 WB 才写入寄存器堆。SUB 在下一拍进入 EX 时若只读 ID 阶段取得的旧 R3 就会出错。EX/MEM 到 ALU 输入的旁路可以把结果直接送给 SUB；旁路选择条件必须同时比较目标寄存器号、源寄存器号和“前条指令确实写寄存器”三项。

load-use 更难：

\begin{lstlisting}
LOAD R3, [R1]
ADD  R5, R3, R4
\end{lstlisting}

加载数据通常到 MEM 末端才出现，紧随其后的 ADD 在同一拍开头已经需要输入，即使有旁路也来不及，必须插入至少一个气泡。若数据 Cache miss，固定的一拍停顿还要升级为“保持请求状态直到响应真正返回”。

\section{八、精确异常要求按程序顺序提交}

流水线里可能同时有多条指令处于不同阶段。发生异常时，架构状态必须看起来像：异常指令之前的都已完成，异常指令及之后的都没有产生可见副作用。实现上应先阻止年轻指令写寄存器、写存储器或发出不可撤销 MMIO，再保存异常 PC 和原因，最后跳到异常入口。

乱序核心进一步把“执行完成”和“按序提交”分开。重排序缓冲记录程序顺序，执行单元可以先算就绪的独立操作，但只有队首指令确认无异常时才能把结果变成架构可见状态。store buffer 也必须等到提交条件成立，才允许写操作越过不可恢复的边界。

\begin{keyidea}
性能优化可以改变内部执行顺序，不能改变程序看到的提交顺序、异常位置和设备副作用顺序。
\end{keyidea}

\section*{章末练习}

\begin{chapterexercise}{时延/吞吐区分}
某服务单核处理一个请求的时延为 2 ms，请求一到就处理、无排队。求单核吞吐（请求/秒）；增加一个相同的核并行后，吞吐与单请求时延各变为多少；再举出一个“吞吐翻倍而时延不变”的硬件机制
\exercisecheck{学习目标}{分清两个量纲，知道各自何时独立变化}
\end{chapterexercise}

\begin{chapterexercise}{加权 CPI}
某程序指令构成为 ALU 50\%、load 20\%、store 10\%、branch 20\%，四类 CPI 分别为 1、2、1、3。求平均 CPI；分支预测改进使 branch 类 CPI 降为 2 后，求新的平均 CPI 与总时间变化百分比（指令数与主频不变）
\exercisecheck{学习目标}{平均 CPI 按占比加权，总时间同比缩放}
\end{chapterexercise}

\begin{chapterexercise}{Amdahl 排序}
某程序中浮点运算占 40\% 时间、cache 访问占 30\%、分支占 20\%；三个候选优化分别把对应部分加速 4 倍、2 倍、1.5 倍，且作用互不重叠。用 Amdahl 定律分别计算整体加速比并给出实施排序
\exercisecheck{学习目标}{占比决定优先级，倍数不等于收益}
\end{chapterexercise}

\begin{chapterexercise}{roofline 落点}
某机器峰值 64 GFLOP/s、内存带宽 16 GB/s。计算拐点算术强度；kernel 甲算术强度 0.5 FLOP/B、kernel 乙 8 FLOP/B，分别判断位于哪一侧并给出可达性能
\exercisecheck{学习目标}{可达性能取峰值与强度乘带宽的较小者}
\end{chapterexercise}

\begin{chapterexercise}{Little 定律在途数}
DDR 读时延 80 ns，每个请求返回一条 64 B 的 cache line。要维持 12.8 GB/s 读带宽，求请求率（请求/秒）与所需平均在途请求数；若时延降到 40 ns，在途数如何变化
\exercisecheck{学习目标}{$L=\lambda W$，带宽靠并发支撑}
\end{chapterexercise}

\begin{chapterexercise}{测量误差识别}
同事报告“优化把函数从 10.2 ms 降到 9.8 ms，快 4\%”：两次各测一次，第一次含冷启动，两台机器后台负载未知。指出至少三个使结论不可靠的因素，并给出能得到可信结论的测法
\exercisecheck{学习目标}{先核对测量约定，再比较数字}
\end{chapterexercise}

\section*{参考解答与要点}

\begin{chapteranswer}{时延/吞吐区分}
单核串行：吞吐 $=1/0.002\ \mathrm{s}=500$ 请求/秒，单请求时延 2 ms。增加一核后两核并行，理想吞吐约 1000 请求/秒，单请求时延仍为 2 ms——时延由单个请求的处理路径决定，与并行度无关。“吞吐翻倍而时延不变”的硬件例子：存储与设备事务相关章节的多 bank 交错与总线流水化，单次访问时延不变，单位时间完成数成倍增加
\answercheck{必须掌握的要点}{时延与吞吐是两个独立的量；提吞吐不等于降时延}
\end{chapteranswer}

\begin{chapteranswer}{加权 CPI}
平均 CPI $=0.5\times1+0.2\times2+0.1\times1+0.2\times3=0.5+0.4+0.1+0.6=1.6$。branch 类 CPI 降为 2 后：$0.5+0.4+0.1+0.2\times2=1.4$。总时间与平均 CPI 成正比，变为 $1.4/1.6=0.875$，即缩短 12.5\%
\answercheck{必须掌握的要点}{权重是指令占比；改进只作用在对应那一项上}
\end{chapteranswer}

\begin{chapteranswer}{Amdahl 排序}
$S=1/(1-f+f/s)$。浮点：$1/(0.6+0.4/4)=1/0.7\approx1.43$；cache：$1/(0.7+0.3/2)=1/0.85\approx1.18$；分支：$1/(0.8+0.2/1.5)\approx1/0.933\approx1.07$。排序：浮点 $>$ cache $>$ 分支。分支优化 1.5 倍听着不小，但占比 20\% 使整体收益只剩约 7\%
\answercheck{必须掌握的要点}{比较候选优化先看占比 $f$，再看倍数 $s$}
\end{chapteranswer}

\begin{chapteranswer}{roofline 落点}
拐点 $=$ 峰值算力 $/$ 带宽 $=64/16=4$ FLOP/B。甲强度 $0.5<4$，位于带宽侧：可达 $=0.5\times16=8$ GFLOP/s，只有峰值的八分之一；乙强度 $8>4$，位于算力侧：可达 $=\min(64,\ 8\times16)=64$ GFLOP/s，用满峰值
\answercheck{必须掌握的要点}{先算拐点再比强度；强度不够时加算力无用}
\end{chapteranswer}

\begin{chapteranswer}{Little 定律在途数}
请求率 $\lambda=12.8\ \mathrm{GB/s}\div64\ \mathrm{B}=2\times10^{8}$ 请求/秒；$L=\lambda W=2\times10^{8}\times80\times10^{-9}=16$ 个在途请求。时延降到 40 ns 则 $L=8$：带宽目标不变时，时延减半，所需并发减半；反之时延降不下来，只能靠增加在途数（更多 MSHR、预取、多 bank）补
\answercheck{必须掌握的要点}{$W$ 用秒、$\lambda$ 用每秒，量纲先配平再代入}
\end{chapteranswer}

\begin{chapteranswer}{测量误差识别}
不可靠因素（任答其三）：两次各测一次，无法区分真实差异与随机噪声；第一次含冷启动而第二次未说明，冷热混测；两台机器后台负载未知，环境不一致；且 4\% 的差异本身就在典型系统噪声幅度之内。可信测法：同一台机器或完整声明环境；预热后各测不少于 30 次；报中位数与分布而非单次值；两次之间只改“是否优化”这一个因子
\answercheck{必须掌握的要点}{核对清单：输入、环境、构建、冷热、样本、量纲}
\end{chapteranswer}
